% Amélioration de la production de l'énergie par l'organisme — SVT 1ère C — Cours signé OnBuch+
% Programme officiel MINESEC. Compiler avec : tectonic NCT01-amelioration-production-energie-organisme.tex
\newcommand{\DOCMATIERE}{SVT}
\newcommand{\DOCNIVEAU}{1ère C}
\newcommand{\DOCMODULE}{Module 1 : Le Monde Vivant}
\newcommand{\DOCLECON}{NCT01}
\newcommand{\DOCTITRE}{Amélioration de la production de l'énergie par l'organisme}
% OnBuch+ — préambule « cours signés » ENRICHI (compilé avec tectonic / XeLaTeX)
% Système visuel « Pop » d'OnBuch+ : fond crème (jamais blanc), bordures encre
% épaisses, ombres « dures » décalées, titres Archivo Black en capitales.
% Version MATHÉMATIQUES : graphes (pgfplots/tikz), boîtes pédagogiques
% (définition, propriété, méthode, attention, le savais-tu, exercice résolu,
% exercice, TP, bilan), page de garde et signature OnBuch+ ancrée en bas.
%
% Macros attendues dans main.tex AVANT \input{preamble} :
%   \DOCMATIERE  \DOCNIVEAU  \DOCMODULE  \DOCLECON (numéro)  \DOCTITRE
\documentclass[11pt]{article}

\usepackage{fontspec}
\usepackage[french]{babel}
\usepackage[a4paper,margin=2cm,top=3.4cm,bottom=2.8cm,headheight=1.4cm,footskip=1.3cm]{geometry}
\usepackage{amsmath,amssymb,amsfonts}
\usepackage{mathtools} % \xmapsto, \coloneqq... souvent utilisés par les modèles
\usepackage{cancel} % simplifications barrées (bilans de forces)
% \abs{z} (module, valeur absolue) et \norm{u} (norme), utilisés par les modèles
\providecommand{\abs}{}\let\abs\relax\DeclarePairedDelimiter{\abs}{\lvert}{\rvert}
\providecommand{\norm}{}\let\norm\relax\DeclarePairedDelimiter{\norm}{\lVert}{\rVert}
\usepackage[table]{xcolor} % [table] : fournit aussi \rowcolors (alternance de lignes)
\usepackage{pagecolor}
\usepackage{tikz}
\usetikzlibrary{arrows.meta,positioning,calc,shapes.geometric,shapes.misc,shapes.arrows,decorations.pathmorphing,decorations.pathreplacing,decorations.markings,patterns,shadows,mindmap,trees,fit,backgrounds,matrix,intersections,angles,quotes}
\usepackage{pgfplots}
\usepackage[version=4]{mhchem} % équations nucléaires \ce{^{235}_{92}U}
\usepackage[european,siunitx]{circuitikz} % schémas électriques
\pgfplotsset{compat=1.18}
\usepgfplotslibrary{fillbetween,groupplots} % aires entre courbes (calcul intégral)
\usepackage{enumitem}
\usepackage{fancyhdr}
% [most] charge toutes les bibliothèques tcolorbox (listings, minted, raster,
% documentation...) et gonfle inutilement la mémoire TeX sur un document long
% (~300 boîtes) : on ne charge que ce qui est réellement utilisé.
\usepackage[skins,breakable]{tcolorbox}
\usepackage{graphicx}
\usepackage{colortbl}
\usepackage{booktabs}
\usepackage{tabularx}
\usepackage{diagbox} % case d'en-tête diagonale des tableaux à double entrée
% Colonnes extensibles centrée / gauche / droite (C, L, R), souvent utilisées par les modèles
\newcolumntype{C}{>{\centering\arraybackslash}X}
\newcolumntype{L}{>{\raggedright\arraybackslash}X}
\newcolumntype{R}{>{\raggedleft\arraybackslash}X}
\usepackage{array}
\usepackage{multicol}
\usepackage{siunitx}
% parse-numbers=false : accepte aussi \SI{2,5 \times 10^{-3}}{...}
\sisetup{locale=FR,output-decimal-marker={,},group-minimum-digits=4,parse-numbers=false}
\DeclareSIUnit\ohm{\ensuremath{\symup{\Omega}}} % U+2126 absent de la police
\DeclareSIUnit\division{div} % réglages d'oscilloscope (ms/div, V/div)
\DeclareSIUnit\tr{tr} % tours (vitesse de rotation, ex. \SI{1500}{\tr\per\minute})
\DeclareSIUnit\ch{ch} % cheval-vapeur (puissance moteur, unité usuelle hors SI)
\DeclareSIUnit\franccfa{FCFA} % franc CFA (coûts, contexte camerounais)
\DeclareSIUnit\dioptre{\ensuremath{\delta}} % vergence d'une lentille (dioptrie)
\usepackage{qrcode}
% \degree : commande fréquemment utilisée par les modèles pour un angle ou une
% température en dehors de \SI{}{} (non fournie par les paquets chargés).
\providecommand{\degree}{^{\circ}}
% \qed (fin de démonstration), utilisé par les modèles sans amsthm
\providecommand{\qed}{\hfill$\square$}
% \barre : double barre des valeurs interdites dans un tableau de variations (tkz-tab)
\providecommand{\barre}{\ensuremath{\|}}
% environnement proof (amsthm non chargé) utilisé par les modèles
\ifdefined\proof\else\newenvironment{proof}[1][Démonstration]{\par\noindent\textit{#1.}\ }{\qed\par}\fi
\usepackage{lastpage}
\usepackage{hyperref}
\hypersetup{hidelinks}

% ============================================================
% Palette « Pop » OnBuch+ — mêmes hex que la classe P (plus_ui.dart)
% ============================================================
\definecolor{poporange}{HTML}{F59321}
\definecolor{popdark}{HTML}{E07A0C}
\definecolor{popcream}{HTML}{FFF6E8}
\definecolor{popink}{HTML}{1C1714}
\definecolor{poppurple}{HTML}{7A5AE0}
\definecolor{popgreen}{HTML}{2E9E6A}
\definecolor{poppink}{HTML}{E0457B}
\definecolor{popblue}{HTML}{2F7DE1}
\definecolor{popgold}{HTML}{C98A12}
\definecolor{popmuted}{HTML}{8A7F76}
\definecolor{popline}{HTML}{EADFD2}
\definecolor{popgray}{HTML}{8A7F76}
\colorlet{popgrey}{popgray}
% Alias tolérants (le modèle nomme parfois les couleurs différemment)
\colorlet{popwhite}{white}
\colorlet{popblack}{popink}
\colorlet{popred}{poppink}
\colorlet{popyellow}{popgold}
% Teintes claires pour fonds de tableaux / zones
\colorlet{popblueL}{popblue!10!white}
\colorlet{poporangeL}{poporange!14!white}
\colorlet{popgreenL}{popgreen!12!white}
\colorlet{poppurpleL}{poppurple!12!white}
\colorlet{poppinkL}{poppink!12!white}
\colorlet{popgoldL}{popgold!14!white}

\pagecolor{popcream}

% ============================================================
% Typographie
% ============================================================
\defaultfontfeatures{Ligatures=TeX}
\setmainfont{PlusJakartaSans-Medium.ttf}[Path=../../../fonts/,BoldFont=PlusJakartaSans-Bold.ttf]
% Maths en sans-serif (Fira Math) avec chiffres et lettres droites de Plus
% Jakarta, pour que les formules se fondent dans le texte.
\usepackage[math-style=ISO,bold-style=ISO]{unicode-math}
\setmathfont{FiraMath-Regular.otf}[Path=../../../fonts/]
\setmathfont{PlusJakartaSans-Medium.ttf}[Path=../../../fonts/,range={up/num,up/{Latin,latin},\mathup}]
\usepackage{icomma} % virgule décimale sans espace en mode math
% Caractères mathématiques Unicode absents des polices de texte (Plus Jakarta,
% Archivo Black) : rendus par leur équivalent mathématique, en texte comme en
% formule — sinon ils s'affichent en carré vide (ex. « Arithmétique dans ℤ »).
\usepackage{newunicodechar}
\newunicodechar{ℝ}{\ensuremath{\mathbb{R}}}
\newunicodechar{ℤ}{\ensuremath{\mathbb{Z}}}
\newunicodechar{ℂ}{\ensuremath{\mathbb{C}}}
\newunicodechar{ℕ}{\ensuremath{\mathbb{N}}}
\newunicodechar{ℚ}{\ensuremath{\mathbb{Q}}}
\newunicodechar{𝔻}{\ensuremath{\mathbb{D}}}
\newunicodechar{α}{\ensuremath{\alpha}}
\newunicodechar{β}{\ensuremath{\beta}}
\newunicodechar{γ}{\ensuremath{\gamma}}
\newunicodechar{δ}{\ensuremath{\delta}}
\newunicodechar{ε}{\ensuremath{\varepsilon}}
\newunicodechar{θ}{\ensuremath{\theta}}
\newunicodechar{λ}{\ensuremath{\lambda}}
\newunicodechar{μ}{\ensuremath{\mu}}
\newunicodechar{σ}{\ensuremath{\sigma}}
\newunicodechar{τ}{\ensuremath{\tau}}
\newunicodechar{φ}{\ensuremath{\varphi}}
\newunicodechar{ψ}{\ensuremath{\psi}}
\newunicodechar{ω}{\ensuremath{\omega}}
\newunicodechar{Σ}{\ensuremath{\Sigma}}
% majuscules produites par \MakeUppercase dans les titres (α-aminés -> Α)
\newunicodechar{Α}{\ensuremath{\alpha}}
\newunicodechar{Β}{\ensuremath{\beta}}
\newunicodechar{Γ}{\ensuremath{\Gamma}}
\newunicodechar{Δ}{\ensuremath{\Delta}}
\newunicodechar{Ε}{\ensuremath{\varepsilon}}
\newunicodechar{Θ}{\ensuremath{\Theta}}
\newunicodechar{Λ}{\ensuremath{\Lambda}}
\newunicodechar{Μ}{\ensuremath{\mu}}
\newunicodechar{Τ}{\ensuremath{\tau}}
\newunicodechar{Φ}{\ensuremath{\Phi}}
\newunicodechar{Ψ}{\ensuremath{\Psi}}
\newunicodechar{Ω}{\ensuremath{\Omega}}
\newunicodechar{↦}{\ensuremath{\mapsto}}
\newunicodechar{∈}{\ensuremath{\in}}
\newunicodechar{∉}{\ensuremath{\notin}}
\newunicodechar{∧}{\ensuremath{\wedge}}
\newunicodechar{⇔}{\ensuremath{\Leftrightarrow}}
\newunicodechar{⟺}{\ensuremath{\Longleftrightarrow}}
\newunicodechar{⇒}{\ensuremath{\Rightarrow}}
\newunicodechar{⟹}{\ensuremath{\Longrightarrow}}
\newunicodechar{∘}{\ensuremath{\circ}}
\newunicodechar{∪}{\ensuremath{\cup}}
\newunicodechar{∩}{\ensuremath{\cap}}
\newunicodechar{≡}{\ensuremath{\equiv}}
\newunicodechar{⊂}{\ensuremath{\subset}}
\newunicodechar{⊥}{\ensuremath{\perp}}
\newunicodechar{∃}{\ensuremath{\exists}}
\newunicodechar{∀}{\ensuremath{\forall}}
\newunicodechar{∛}{\ensuremath{\sqrt[3]{\,}}}
\newunicodechar{∜}{\ensuremath{\sqrt[4]{\,}}}
\newunicodechar{⃗}{\ensuremath{{}^{\to}}}
\newunicodechar{ⁿ}{\ifmmode^{n}\else\textsuperscript{\ensuremath{n}}\fi}
\newunicodechar{ˣ}{\ifmmode^{x}\else\textsuperscript{\ensuremath{x}}\fi}
\newunicodechar{ᵇ}{\ifmmode^{b}\else\textsuperscript{\ensuremath{b}}\fi}
\newunicodechar{ʳ}{\ifmmode^{r}\else\textsuperscript{\ensuremath{r}}\fi}
\newunicodechar{ᵘ}{\ifmmode^{u}\else\textsuperscript{\ensuremath{u}}\fi}
\newunicodechar{ʸ}{\ifmmode^{y}\else\textsuperscript{\ensuremath{y}}\fi}
\newunicodechar{ᵃ}{\ifmmode^{a}\else\textsuperscript{\ensuremath{a}}\fi}
\newunicodechar{ᵉ}{\ifmmode^{e}\else\textsuperscript{\ensuremath{e}}\fi}
\newunicodechar{ᵏ}{\ifmmode^{k}\else\textsuperscript{\ensuremath{k}}\fi}
\newunicodechar{ᵖ}{\ifmmode^{p}\else\textsuperscript{\ensuremath{p}}\fi}
\newunicodechar{ᴺ}{\ifmmode^{N}\else\textsuperscript{\ensuremath{N}}\fi}
\newunicodechar{ᶜ}{\ifmmode^{c}\else\textsuperscript{\ensuremath{c}}\fi}
\newunicodechar{ʰ}{\ifmmode^{h}\else\textsuperscript{\ensuremath{h}}\fi}
\newunicodechar{ᵠ}{\ifmmode^{q}\else\textsuperscript{\ensuremath{q}}\fi}
\newunicodechar{ₙ}{\ifmmode_{n}\else\textsubscript{\ensuremath{n}}\fi}
\newunicodechar{ₐ}{\ifmmode_{a}\else\textsubscript{\ensuremath{a}}\fi}
\newunicodechar{ᵢ}{\ifmmode_{i}\else\textsubscript{\ensuremath{i}}\fi}
\newunicodechar{ᵣ}{\ifmmode_{r}\else\textsubscript{\ensuremath{r}}\fi}
\newunicodechar{ₖ}{\ifmmode_{k}\else\textsubscript{\ensuremath{k}}\fi}
\newunicodechar{ᵦ}{\ifmmode_{\beta}\else\textsubscript{\ensuremath{\beta}}\fi}
\newunicodechar{ₓ}{\ifmmode_{x}\else\textsubscript{\ensuremath{x}}\fi}
\newfontfamily\popdisplay{ArchivoBlack-Regular.ttf}[Path=../../../fonts/]
\newfontfamily\poplogo{SpaceGrotesk-Bold.ttf}[Path=../../../fonts/]
\newfontfamily\popsemi{PlusJakartaSans-SemiBold.ttf}[Path=../../../fonts/]
\newfontfamily\popxbold{PlusJakartaSans-ExtraBold.ttf}[Path=../../../fonts/]
\newfontfamily\popxboldls{PlusJakartaSans-ExtraBold.ttf}[Path=../../../fonts/,LetterSpace=3.0]

\setlist[enumerate]{leftmargin=1.7em,itemsep=5pt,topsep=4pt}
\setlist[itemize]{leftmargin=1.7em,itemsep=4pt,topsep=4pt,label={\color{poporange}\textbullet}}
\setlength{\parindent}{0pt}
\setlength{\parskip}{5pt}
\renewcommand{\arraystretch}{1.25}

% pgfplots : style Pop par défaut
\pgfplotsset{
  every axis/.append style={axis line style={popink,line width=1pt},tick style={popink},
    grid style={popline,line width=0.5pt},label style={font=\small},tick label style={font=\footnotesize}},
  popcurve/.style={poporange,line width=1.8pt,no markers,smooth},
}
% Même style disponible dans un \draw TikZ simple (hors pgfplots)
\tikzset{popcurve/.style={poporange,line width=1.8pt}}
\tikzset{popgrid/.style={popline,very thin}}
\colorlet{popcurve}{poporange} % le nom du style est parfois utilisé comme couleur

% ============================================================
% Pill / squiggle
% ============================================================
\newcommand{\poppill}[3]{%
  \tikz[baseline=-0.55ex]\node[draw=popink,line width=1.2pt,rounded corners=2.6pt,%
    fill=#2,inner xsep=6.5pt,inner ysep=3pt]{%
    \popxboldls\fontsize{7}{7}\selectfont\color{#3}\MakeUppercase{#1}};%
}
\newcommand{\popsquiggle}[1]{%
  \tikz[baseline=-0.4ex]{
    \draw[#1,line width=2.6pt,line cap=round]
      plot [smooth] coordinates {(0,0.09) (0.34,0.01) (0.68,0.21) (1.02,0.01) (1.36,0.10)};
  }%
}
% Mot-clé surligné (marqueur orange sous le texte)
\newcommand{\cle}[1]{{\popxbold\color{popdark}#1}}

% ============================================================
% En-tête / pied
% ============================================================
\pagestyle{fancy}
\fancyhf{}
\renewcommand{\headrulewidth}{0pt}
\renewcommand{\footrulewidth}{0pt}
\lhead{\raisebox{-0.15ex}{\poplogo\fontsize{13}{13}\selectfont\color{popink}OnBuch\color{poporange}+}}
\rhead{\poppill{\DOCMATIERE\ \textbullet\ \DOCNIVEAU\ \textbullet\ Leçon \DOCLECON}{white}{popink}}
\lfoot{\popxbold\fontsize{7.6}{7.6}\selectfont\color{popmuted}\MakeUppercase{Cours signé OnBuch+}\ \textbullet\ {\popsemi\DOCTITRE}}
\rfoot{\tikz[baseline=-0.6ex]\node[rounded corners=8.5pt,fill=popink,minimum height=17pt,minimum width=17pt,inner xsep=5pt,inner sep=0]{\popxbold\fontsize{8}{8}\selectfont\color{popcream}\thepage\,/\,\pageref*{LastPage}};}

% ============================================================
% Titres
% ============================================================
\newcommand{\chaptitre}[2]{%
  \begin{tcolorbox}[enhanced,colback=poporange,colframe=popink,boxrule=2.6pt,arc=11pt,
      drop shadow=popink,left=15pt,right=15pt,top=12pt,bottom=11pt,before skip=3pt,after skip=18pt]
    \poppill{#2}{popink}{popcream}\par\vspace{9pt}
    {\raggedright\hyphenpenalty=10000\exhyphenpenalty=10000\popdisplay\fontsize{22}{24}\selectfont\color{popcream}\MakeUppercase{#1}\par}
  \end{tcolorbox}
}
% Section numérotée : pastille orange avec le numéro + titre Archivo
\newcounter{popsec}
\newcommand{\coursec}[1]{%
  \stepcounter{popsec}\setcounter{popsub}{0}%
  \par\vspace{16pt}\noindent
  \begin{minipage}{\linewidth}
  \tikz[baseline=-0.7ex]\node[draw=popink,line width=1.6pt,fill=poporange,rounded corners=4pt,minimum size=22pt,inner sep=0]{\popdisplay\fontsize{12}{12}\selectfont\color{popcream}\thepopsec};%
  \hspace{8pt}{\popdisplay\fontsize{15}{17}\selectfont\color{popink}\MakeUppercase{#1}}\par
  \vspace{4pt}\noindent{\color{poporange}\rule{36pt}{3.6pt}}
  \end{minipage}\par\nopagebreak\vspace{8pt}
}
\newcounter{popsub}
\newcommand{\courssub}[1]{%
  \stepcounter{popsub}%
  \par\vspace{9pt}\noindent{\popxbold\fontsize{11.5}{13}\selectfont\color{popink}{\color{poporange}\thepopsec.\thepopsub}\hspace{6pt}#1}\par\nopagebreak\vspace{4pt}
}

% ============================================================
% Boîtes pédagogiques — même recette : fond blanc, bordure épaisse colorée,
% ombre dure encre, badge plein. Titre optionnel [#1].
% ============================================================
\tcbset{popbase/.style={breakable,enhanced,colback=white,boxrule=2.3pt,arc=9pt,
  shadow={3pt}{-3pt}{0pt}{popink},left=14pt,right=14pt,top=11pt,bottom=11pt,before skip=11pt,after skip=15pt,
  fontupper=\popsemi\fontsize{10.6}{14.5}\selectfont\color{popink},
  fontlower=\popsemi\fontsize{10.6}{14.5}\selectfont\color{popink}}}
% Coche dessinée (le glyphe ✓ n'existe pas dans les polices du document)
\newcommand{\popcheck}[1]{\tikz[baseline=-0.5ex]\draw[#1,line width=1.6pt,line cap=round,line join=round](-0.13,0.0)--(-0.04,-0.09)--(0.13,0.1);}
\providecommand{\coche}{\popcheck{popgreen}}
\providecommand{\resultat}[1]{\par\smallskip\noindent\fcolorbox{popgreen}{popgreenL}{\parbox{\dimexpr\linewidth-2\fboxsep-2\fboxrule\relax}{\textbf{Résultat.} #1}}\par\smallskip}
\newcommand{\popboxhead}[3]{\poppill{#1}{#2}{white}\ifx\relax#3\relax\else\hspace{6pt}{\popxbold\fontsize{10.6}{12}\selectfont\color{popink}#3}\fi\par\vspace{7pt}}

\newtcolorbox{aretenir}[1][]{popbase,colframe=popgreen,before upper={\popboxhead{À retenir}{popgreen}{#1}}}
\newtcolorbox{exemplebox}[1][]{popbase,colframe=poporange,before upper={\popboxhead{Exemple}{poporange}{#1}}}
\newtcolorbox{definition}[1][]{popbase,colframe=popblue,before upper={\popboxhead{Définition}{popblue}{#1}}}
\newtcolorbox{propriete}[1][]{popbase,colframe=poppurple,before upper={\popboxhead{Propriété}{poppurple}{#1}}}
\newtcolorbox{methode}[1][]{popbase,colframe=popgold,before upper={\popboxhead{Méthode}{popgold}{#1}}}
\newtcolorbox{attention}[1][]{popbase,colframe=poppink,before upper={\popboxhead{Attention}{poppink}{#1}}}
\newtcolorbox{experience}[1][]{popbase,colframe=popblue,colback=popblueL,before upper={\popboxhead{Expérience}{popblue}{#1}}}
% « Le savais-tu ? » : carte violette pleine (contexte, culture, Cameroun)
\newtcolorbox{savaistu}[1][]{popbase,colback=poppurple,colframe=popink,
  fontupper=\popsemi\fontsize{10.4}{14.2}\selectfont\color{white},
  before upper={\poppill{Le savais-tu\,?}{white}{poppurple}\ifx\relax#1\relax\else\hspace{6pt}{\popxbold\color{white}#1}\fi\par\vspace{7pt}}}
% Exercice résolu : énoncé en haut, \tcblower puis la solution sur fond crème
\newcounter{popexo}\newcounter{popexr}
\newtcolorbox{exoresolu}[1][]{popbase,colframe=poporange,colbacklower=popcream,
  segmentation style={popink,line width=1.2pt,dashed},
  before upper={\stepcounter{popexr}\popboxhead{Exercice résolu \thepopexr}{poporange}{#1}},
  before lower={\poppill{Solution}{popink}{popcream}\par\vspace{6pt}}}
% Exercice (non résolu en place) — la correction est dans la partie Corrigés
\newtcolorbox{exercice}[1][]{popbase,colframe=popink,
  before upper={\stepcounter{popexo}\popboxhead{Exercice \thepopexo}{popink}{#1}}}
% Corrigé
\newtcolorbox{corrige}[1][]{popbase,colframe=popgreen,colback=popgreenL,
  before upper={\popboxhead{Corrigé}{popgreen}{#1}}}
% Objectifs / prérequis / bilan
\newtcolorbox{objectifs}{popbase,colframe=popink,colback=poporangeL,
  before upper={\popboxhead{Objectifs de la leçon}{popink}{}}}
\newtcolorbox{prerequis}{popbase,colframe=popink,colback=popblueL,
  before upper={\popboxhead{Prérequis}{popblue}{}}}
\newtcolorbox{bilan}{popbase,colframe=popink,colback=white,boxrule=2.8pt,
  before upper={\popboxhead{Fiche bilan}{poporange}{L'essentiel à connaître pour l'examen}}}
% Figure encadrée avec légende
\newtcolorbox{popfigure}[1][]{enhanced,breakable=false,colback=white,colframe=popline,boxrule=1.4pt,arc=8pt,
  left=10pt,right=10pt,top=10pt,bottom=8pt,before skip=10pt,after skip=12pt,halign=center,
  lower separated=false,halign lower=center,
  fontlower=\popsemi\fontsize{9}{11.5}\selectfont\color{popmuted},
  #1}
\newcommand{\legende}[1]{\par\vspace{6pt}{\popsemi\fontsize{9}{11.5}\selectfont\color{popmuted}#1}}

% ============================================================
% Signature OnBuch+
% ============================================================
\newcommand{\poplogoinline}[1]{{\poplogo\fontsize{#1}{#1}\selectfont\color{popink}OnBuch\color{poporange}+}}
\newcommand{\signatureonbuch}{%
  \begin{tcolorbox}[enhanced,colback=popink,colframe=popink,boxrule=2.2pt,arc=10pt,
      drop shadow=poporange,left=14pt,right=14pt,top=10pt,bottom=10pt,before skip=0pt,after skip=0pt]
    \tikz[baseline=-0.7ex]\node[circle,fill=popgreen,draw=popcream,line width=1.3pt,minimum size=22pt,inner sep=0]{\popcheck{white}};%
    \hspace{9pt}\begin{minipage}[c]{0.62\linewidth}
      {\popxbold\fontsize{12}{13}\selectfont\color{popcream}Signé\ \textbullet\ {\poplogo OnBuch\color{poporange}+}}\par\vspace{2pt}
      {\raggedright\popsemi\fontsize{8}{10.5}\selectfont\color{popline}\MakeUppercase{Équipe pédagogique OnBuch+}\\\MakeUppercase{Conforme au programme officiel MINESEC}\par}
    \end{minipage}\hfill
    {\poplogo\fontsize{22}{22}\selectfont\color{popcream}OnBuch\color{poporange}+}
  \end{tcolorbox}%
}

% ============================================================
% Page de garde
% #1 titre · #2 matière · #3 niveau · #4 module · #5 n° leçon · #6 durée ·
% #7 description / objectifs (texte ou itemize)
% ============================================================
\newcommand{\pagegarde}[7]{%
  \thispagestyle{empty}
  \newgeometry{a4paper,margin=1.7cm,top=1.6cm,bottom=1.4cm}%
  \begin{tikzpicture}[remember picture,overlay]
    \fill[poporange!18] ([xshift=-3.2cm,yshift=-3.6cm]current page.north east) circle (4.6cm);
    \fill[poppurple!14] ([xshift=2.4cm,yshift=7.6cm]current page.south west) circle (3.8cm);
  \end{tikzpicture}%
  {\poplogo\fontsize{30}{30}\selectfont\color{popink}OnBuch\color{poporange}+}\hfill
  \raisebox{0.6ex}{\poppill{Cours signé \textbullet\ Premium}{popink}{popcream}}\par
  \vspace{1pt}\popsquiggle{poppurple}\par
  \vspace{8pt}
  \begin{tcolorbox}[enhanced,colback=poporange,colframe=popink,boxrule=2.8pt,arc=13pt,
      drop shadow=popink,left=18pt,right=18pt,top=12pt,bottom=13pt,after skip=10pt]
    \poppill{#2 \textbullet\ #3}{popink}{popcream}\hspace{5pt}\poppill{#4}{white}{popink}\par\vspace{8pt}
    {\popdisplay\fontsize{14}{14}\selectfont\color{popink}LEÇON #5}\par\vspace{4pt}
    {\raggedright\hyphenpenalty=10000\exhyphenpenalty=10000\popdisplay\fontsize{25}{27}\selectfont\color{popcream}\MakeUppercase{#1}\par}
    \vspace{8pt}
    {\popxbold\fontsize{9.5}{9.5}\selectfont\color{popink}\MakeUppercase{Durée conseillée : #6}}
  \end{tcolorbox}
  \begin{tcolorbox}[enhanced,colback=white,colframe=popink,boxrule=2.2pt,arc=10pt,
      drop shadow=popink,left=16pt,right=16pt,top=10pt,bottom=10pt,after skip=8pt]
    \poppill{Dans cette leçon}{poppurple}{white}\par\vspace{5pt}
    \popsemi\fontsize{9.3}{12.6}\selectfont\color{popink}#7
  \end{tcolorbox}
  \noindent\begin{minipage}[t]{0.487\linewidth}
    \begin{tcolorbox}[enhanced,colback=popgreen,colframe=popink,boxrule=2.2pt,arc=10pt,
        drop shadow=popink,left=11pt,right=11pt,top=10pt,bottom=10pt,height=3.1cm,valign=center]
      \tcbox[colback=white,colframe=popink,boxrule=1.3pt,arc=5pt,boxsep=0pt,nobeforeafter,
        left=3pt,right=3pt,top=3pt,bottom=3pt,valign=center]{\qrcode[height=1.9cm]{https://play.google.com/store/apps/details?id=com.luvvix.onbuch&hl=fr}}%
      \hspace{8pt}\begin{minipage}[c]{0.5\linewidth}
        {\popdisplay\fontsize{11}{12.5}\selectfont\color{white}\MakeUppercase{Télécharge OnBuch}}\par\vspace{4pt}
        {\raggedright\popsemi\fontsize{8.4}{11}\selectfont\color{white}Gratuit \textbullet\ Google Play\\Tuteur IA, annales, concours, résultats\par}
      \end{minipage}
    \end{tcolorbox}
  \end{minipage}\hfill
  \begin{minipage}[t]{0.487\linewidth}
    \begin{tcolorbox}[enhanced,colback=poppurple,colframe=popink,boxrule=2.2pt,arc=10pt,
        drop shadow=popink,left=11pt,right=11pt,top=10pt,bottom=10pt,height=3.1cm,valign=center]
      \tikz[baseline=-0.7ex]\node[circle,fill=white,draw=popink,line width=1.3pt,minimum size=1.6cm,inner sep=0]{\popdisplay\fontsize{24}{24}\selectfont\color{poppurple}+};%
      \hspace{8pt}\begin{minipage}[c]{0.6\linewidth}
        {\popdisplay\fontsize{11}{12.5}\selectfont\color{white}\MakeUppercase{Passe à OnBuch+}}\par\vspace{4pt}
        {\raggedright\popsemi\fontsize{8.4}{11}\selectfont\color{white}Tous les cours signés, Tuteur IA illimité, fiches PDF — onglet OnBuch+ de l'app\par}
      \end{minipage}
    \end{tcolorbox}
  \end{minipage}
  \vspace{8pt}
  \popcontenu
  \vfill
  \signatureonbuch
  \restoregeometry
  \newpage
}

% Rangée « Ce cours contient » (page de garde)
\newcommand{\poptile}[3]{% #1 couleur #2 gros texte #3 légende
  \begin{tcolorbox}[enhanced,colback=#1,colframe=popink,boxrule=2pt,arc=9pt,shadow={2.5pt}{-2.5pt}{0pt}{popink},
      left=6pt,right=6pt,top=9pt,bottom=9pt,halign=center,valign=center,height=1.75cm,width=\linewidth,nobeforeafter]
    {\popdisplay\fontsize{12.5}{13}\selectfont\color{white}\MakeUppercase{#2}}\par\vspace{3pt}
    {\centering\popsemi\fontsize{7.8}{9.5}\selectfont\color{white}#3\par}
  \end{tcolorbox}}
\newcommand{\popcontenu}{%
  \par\noindent{\popxboldls\fontsize{7.5}{7.5}\selectfont\color{popmuted}CE COURS SIGNÉ CONTIENT}\par\vspace{7pt}
  \noindent\begin{minipage}[t]{0.235\linewidth}\poptile{popblue}{Cours}{complet, illustré, pas à pas}\end{minipage}\hfill
  \begin{minipage}[t]{0.235\linewidth}\poptile{poporange}{Méthodes}{et exercices résolus}\end{minipage}\hfill
  \begin{minipage}[t]{0.235\linewidth}\poptile{poppink}{Exercices}{type examen + corrigés}\end{minipage}\hfill
  \begin{minipage}[t]{0.235\linewidth}\poptile{popgreen}{Fiche bilan}{l'essentiel à retenir}\end{minipage}\par}

% Sommaire
\newtcolorbox{sommaire}{enhanced,colback=white,colframe=popink,boxrule=2.2pt,arc=10pt,
  drop shadow=popink,left=18pt,right=18pt,top=13pt,bottom=13pt,before skip=6pt,after skip=20pt,
  before upper={\poppill{Sommaire}{poppurple}{white}\par\vspace{9pt}\popsemi\fontsize{10.6}{14.5}\selectfont\color{popink}}}

% Signature de fin de cours — poussée tout en bas de la dernière page
\newcommand{\findecours}{%
  \par\vfill\nopagebreak
  \begin{center}\popsquiggle{poporange}\end{center}\vspace{4pt}
  \signatureonbuch
}
\AtBeginDocument{\def\llbracket{\mathopen{[\mkern-3.5mu[}}\def\rrbracket{\mathclose{]\mkern-3.5mu]}}} % intervalles d'entiers (glyphe ⟦ absent de la police)
% Glyphes absents de Fira Math : redéfinis à partir de glyphes disponibles
\AtBeginDocument{%
  \def\setminus{\mathbin{\backslash}}\let\smallsetminus\setminus
  \def\vdots{\mathord{\vbox{\baselineskip3.2pt\lineskiplimit0pt\kern4pt\hbox{.}\hbox{.}\hbox{.}}}}%
  \def\ddots{\mathinner{\mkern1mu\raise6.5pt\hbox{.}\mkern2mu\raise3.6pt\hbox{.}\mkern2mu\raise0.7pt\hbox{.}\mkern1mu}}%
  \def\triangle{\mathord{\scalebox{0.9}{$\symup{\Delta}$}}}%
  \def\nabla{\mathord{\raisebox{\depth}{\scalebox{1}[-1]{$\symup{\Delta}$}}}}%
  \def\checkmark{\popcheck{popdark}}%
  \def\star{\mathbin{\ast}}\def\bigstar{\mathord{\ast}}%
  \def\diamond{\mathbin{\scalebox{0.6}{$\Diamond$}}}%
  \def\bigcup{\mathop{\mathchoice{\raisebox{-0.3em}{\scalebox{1.7}{$\cup$}}}{\scalebox{1.25}{$\cup$}}{\cup}{\cup}}\displaylimits}%
  \def\bigcap{\mathop{\mathchoice{\raisebox{-0.3em}{\scalebox{1.7}{$\cap$}}}{\scalebox{1.25}{$\cap$}}{\cap}{\cap}}\displaylimits}%
  \def\lesseqqgtr{\mathrel{\lessgtr}}\def\bigoplus{\mathop{\oplus}\displaylimits}%
}
\providecommand{\ssi}{si et seulement si} % abréviation fréquente
\AtBeginDocument{\providecommand{\wideparen}{\overparen}} % arc de cercle

%% FIGFIX-BEGIN v4 — correctifs globaux des figures (docs/onbuch-generation/tools/figfix)
\usepackage{adjustbox}\usepackage{etoolbox}
% toute figure TikZ/pgfplots plus large que la ligne est réduite à la largeur disponible
\BeforeBeginEnvironment{tikzpicture}{\begin{adjustbox}{max width=\linewidth}}
\AfterEndEnvironment{tikzpicture}{\end{adjustbox}}
% légendes pgfplots lisibles par-dessus les courbes
\pgfplotsset{every axis/.append style={legend style={fill=white,fill opacity=0.92,text opacity=1,draw=black!15,inner sep=3pt}}}
% étiquettes d'arêtes (syntaxe edge["..."]) avec fond blanc
\tikzset{every edge quotes/.append style={fill=white,inner sep=1.2pt}}
%% FIGFIX-END
\begin{document}

\pagegarde{Amélioration de la production de l'énergie par l'organisme}{SVT}{1ère C}{Module 1 : Le Monde Vivant}{NCT01}{8 h}{Cette leçon explique comment la cellule transforme le glucose en énergie utilisable (ATP) grâce à la respiration cellulaire dans les mitochondries, et comment elle s'adapte en l'absence d'oxygène par la fermentation. Elle permet aussi de comparer les rendements énergétiques de ces voies et d'exploiter des données expérimentales sur les échanges gazeux.
\vspace{4pt}
\begin{multicols}{2}\small
\begin{itemize}
\item À la fin de cette leçon, je sais définir la respiration cellulaire et situer ses trois étapes (glycolyse, cycle de Krebs, chaîne respiratoire)
\item À la fin de cette leçon, je sais écrire le bilan simplifié de la respiration cellulaire
\item À la fin de cette leçon, je sais décrire la structure d'une mitochondrie et relier chaque compartiment à une étape de la respiration
\item À la fin de cette leçon, je sais définir la fermentation et citer ses deux types (lactique et alcoolique)
\item À la fin de cette leçon, je sais comparer respiration et fermentation en termes de rendement énergétique (ATP produits)
\item À la fin de cette leçon, je sais exploiter des données expérimentales sur la consommation d'O2 et la production de CO2 (courbes, tableaux)
\end{itemize}\end{multicols}}

\begin{sommaire}
\begin{multicols}{2}
\begin{enumerate}
\item Mise en évidence des échanges gazeux cellulaires
\item La mitochondrie, centrale énergétique de la cellule
\item Les trois étapes de la respiration cellulaire
\item La fermentation, voie alternative sans oxygène
\item Facteurs influençant le rendement énergétique
\item Synthèse et applications
\item Activité d'intégration
\item Méthodes et exercices résolus
\item Exercices
\item Corrigés des exercices
\item Fiche bilan
\end{enumerate}
\end{multicols}
\end{sommaire}

\begin{objectifs}
\begin{itemize}
\item À la fin de cette leçon, je sais définir la respiration cellulaire et situer ses trois étapes (glycolyse, cycle de Krebs, chaîne respiratoire)
\item À la fin de cette leçon, je sais écrire le bilan simplifié de la respiration cellulaire
\item À la fin de cette leçon, je sais décrire la structure d'une mitochondrie et relier chaque compartiment à une étape de la respiration
\item À la fin de cette leçon, je sais définir la fermentation et citer ses deux types (lactique et alcoolique)
\item À la fin de cette leçon, je sais comparer respiration et fermentation en termes de rendement énergétique (ATP produits)
\item À la fin de cette leçon, je sais exploiter des données expérimentales sur la consommation d'O2 et la production de CO2 (courbes, tableaux)
\item À la fin de cette leçon, je sais expliquer l'influence de l'entraînement, de l'alimentation et de l'oxygénation sur le rendement énergétique de la cellule
\item À la fin de cette leçon, je sais interpréter des phénomènes du quotidien (crampe, essoufflement, levée du vin de palme) à l'aide des notions de la leçon
\end{itemize}
\end{objectifs}

\begin{prerequis}
\begin{itemize}
\item Notion de cellule et d'organites (noyau, cytoplasme, membrane) vue en classe de Seconde
\item Notion de réaction chimique et d'équation-bilan (enseignement de chimie)
\item Nutriments énergétiques : glucides, lipides, protides (alimentation, Seconde)
\item Notion d'enzyme et de catalyse biologique
\item Lecture et exploitation de tableaux et de courbes simples
\end{itemize}
\end{prerequis}

\newpage

\coursec{Mise en évidence des échanges gazeux cellulaires}

\textbf{Situation d'accroche.} Lors de la grande course annuelle de Yaoundé, deux athlètes s'élancent ensemble. Le premier sprinte dès le départ : au bout de quelques centaines de mètres, il s'effondre, essoufflé, les jambes prises de crampes violentes. Le second gère son effort, respire régulièrement, et franchit la ligne d'arrivée sans douleur. Pourquoi les muscles du premier produisent-ils une substance douloureuse alors que ceux du second fonctionnent efficacement ? Pourquoi notre fréquence respiratoire augmente-t-elle pendant l'effort ? Et pourquoi les footballeurs camerounais consomment-ils des boissons sucrées à la mi-temps ?

Toutes ces questions ont une racine commune : la manière dont nos \cle{cellules} produisent l'\cle{énergie} indispensable à leur fonctionnement. Comprendre ce mécanisme commence par une observation simple : tout être vivant échange des gaz avec son milieu. C'est cette observation que nous allons transformer en démarche scientifique rigoureuse.

\textbf{Problématique.} Comment mettre en évidence expérimentalement les échanges gazeux des organismes vivants, et que nous apprennent-ils sur la production d'énergie par la cellule ?

\courssub{Un constat expérimental : les êtres vivants consomment du dioxygène et rejettent du dioxyde de carbone}

\textbf{Observation de la vie courante.} Une bougie placée sous un bocal retourné finit par s'éteindre. Un sac en plastique fermé contenant des fruits mûrit et se remplit d'une atmosphère « étouffante ». Un animal enfermé dans une boîte étanche meurt. Ces observations suggèrent que les êtres vivants, comme la flamme, « consomment » quelque chose dans l'air et « rejettent » autre chose.

\textbf{Hypothèse.} Les organismes vivants absorbent du \cle{dioxygène} (\ce{O2}) et rejettent du \cle{dioxyde de carbone} (\ce{CO2}).

Pour tester cette hypothèse, il faut des expériences contrôlées, avec des témoins. C'est l'objet des deux expériences de référence ci-dessous.

\courssub{Expérience de référence 1 : les graines germées dans un respiromètre (manomètre)}

\begin{experience}[Mise en évidence de la consommation de \ce{O2} et de la production de \ce{CO2} par des graines germées]
\textbf{Matériel.} Deux respiromètres identiques (flacon fermé muni d'un tube capillaire gradué et d'une coupelle) ; des graines de haricot germées (vivantes) ; des graines bouillies (mortes, témoin) ; de la potasse (KOH) solide ou en solution concentrée ; de l'eau colorée (pour le manomètre) ; de l'eau de chaux.

\textbf{Protocole.}
\begin{enumerate}
\item Dans le respiromètre A (expérimental), on place des graines germées vivantes au fond du flacon et de la potasse dans la coupelle. Dans le respiromètre B (témoin), on place des graines bouillies et de la potasse.
\item On ferme hermétiquement les deux appareils et on met le bout du capillaire dans un bain d'eau colorée. On note la position initiale du ménisque (temps $t_0$).
\item On attend 24 à 48 heures à température constante. La potasse absorbe le \ce{CO2} produit : \ce{CO2 + 2KOH -> K2CO3 + H2O}.
\item On observe le déplacement du ménisque dans le capillaire (mesure de la consommation de \ce{O2}).
\item En parallèle, on peut prélever l'air d'un troisième flacon contenant des graines vivantes \textbf{sans potasse} et le faire barboter dans de l'eau de chaux pour mettre en évidence le \ce{CO2}.
\end{enumerate}

\textbf{Observations.}
\begin{itemize}
\item Respiromètre A (graines vivantes + KOH) : le ménisque se déplace \textbf{vers l'intérieur} du flacon (pression interne diminue). L'eau de chaux du dispositif parallèle se \textbf{trouble}.
\item Respiromètre B (graines mortes + KOH) : le ménisque reste \textbf{immobile}. L'eau de chaux reste \textbf{limpide}.
\end{itemize}

\textbf{Interprétation.} La potasse absorbe le \ce{CO2} rejeté par les graines vivantes. La consommation nette de \ce{O2} (non compensée par une production de gaz) crée une dépression qui aspire le ménisque. Seules les graines vivantes consomment du \ce{O2} et rejettent du \ce{CO2}. Le trouble de l'eau de chaux (dispositif sans potasse) confirme la production de \ce{CO2} : $\ce{CO2 + Ca(OH)2 -> CaCO3 v + H2O}$.

\textbf{Conclusion.} Les organismes vivants (ici des végétaux en germination) absorbent du dioxygène et rejettent du dioxyde de carbone : ce sont les \cle{échanges gazeux respiratoires}.
\end{experience}

\begin{popfigure}
\begin{center}
\begin{tikzpicture}[scale=0.9, every node/.style={font=\small}]
% Flacon principal
\draw[thick, popdark] (0,0) -- (0,3.5) -- (0.4,4.0) -- (0.4,4.5);
\draw[thick, popdark] (3.6,0) -- (3.6,3.5) -- (3.2,4.0) -- (3.2,4.5);
\draw[thick, popdark] (0,0) -- (3.6,0);
% Bouchon percé
\draw[thick, popdark, fill=popgoldL] (0.4,4.4) rectangle (3.2,4.7);
% Coupelle KOH
\draw[thick, poppurple, fill=poppurpleL] (0.6,2.8) rectangle (1.6,3.2);
\node[font=\scriptsize, poppurple, align=center] at (1.1,3.0) {Potasse (KOH)\\(absorbe \ce{CO2})};
% Graines germées
\foreach \x/\y in {0.7/0.3, 1.3/0.4, 1.9/0.3, 2.5/0.4, 1.0/0.7, 1.6/0.65, 2.2/0.7, 2.8/0.6}{
  \draw[fill=popgreenL, draw=popgreen] (\x,\y) ellipse (0.18 and 0.12);
  \draw[popgreen, thick] (\x,\y+0.1) .. controls (\x+0.08,\y+0.3) .. (\x+0.15,\y+0.4);
}
% Capillaire (manomètre)
\draw[thick, popblue] (1.9,4.5) -- (1.9,6.5);
\draw[thick, popblue] (2.1,4.5) -- (2.1,6.5);
\draw[fill=popblue!30] (1.9,5.0) rectangle (2.1,5.3);
\draw[thin, popmuted] (1.9,5.0) -- (2.1,5.0);
\node[font=\scriptsize, popblue, right] at (2.15,5.15) {Ménisque};
\draw[<->, thin, popdark] (2.3,4.5) -- (2.3,5.15) node[midway, right, font=\scriptsize] {$\Delta h$};
% Bain d'eau colorée
\draw[thick, popdark, fill=popblue!10] (1.0,6.5) rectangle (3.0,6.8);
\draw[thin, popmuted] (1.9,6.5) -- (1.9,6.8);
% Légendes
\draw[thin, popmuted] (0.5,0.5) -- (-1.0,0.5) node[left, align=right, text=popdark] {Graines germées\\(vivantes)};
\draw[thin, popmuted] (1.9,6.5) -- (3.8,6.5) node[right, text=popdark] {Bain d'eau colorée};
\draw[thin, popmuted] (1.9,4.5) -- (3.8,4.5) node[right, text=popdark] {Tube capillaire gradué};
\end{tikzpicture}
\end{center}
\legende{Respiromètre (manomètre à KOH) pour mesurer la consommation de dioxygène. La potasse absorbe le \ce{CO2} produit ; la baisse de pression due à la consommation de \ce{O2} fait monter le ménisque dans le capillaire. $\Delta h$ est proportionnelle au volume d'\ce{O2} consommé.}
\end{popfigure}

Les mesures quantitatives réalisées sur un tel dispositif (ou par analyse chromatographique de l'air) donnent l'évolution suivante des compositions gazeuses :

\begin{popfigure}
\begin{center}
\begin{tikzpicture}
\begin{axis}[
  width=0.9\linewidth, height=6.5cm,
  xlabel={Temps (heures)}, ylabel={Taux de gaz (\%)},
  xmin=0, xmax=48, ymin=0, ymax=22,
  xtick={0,12,24,36,48},
  ytick={0,5,10,15,20,21},
  legend style={at={(0.98,0.98)}, anchor=north east, font=\small, fill=popcream, draw=popdark},
  grid=major, grid style={popline!40},
  axis line style={popdark},
  tick style={popdark},
]
\addplot[popcurve, popblue, very thick, domain=0:48, samples=100]{21*exp(-x/30)};
\addlegendentry{Taux de \ce{O2} (décroît)}
\addplot[popcurve, poporange, very thick, domain=0:48, samples=100]{0.03+18*(1-exp(-x/30))};
\addlegendentry{Taux de \ce{CO2} (croît)}
\end{axis}
\end{tikzpicture}
\end{center}
\legende{Évolution des taux de dioxygène (décroissant) et de dioxyde de carbone (croissant) dans un milieu confiné contenant des graines germées vivantes. La consommation de \ce{O2} s'accompagne d'une production de \ce{CO2} en quantités comparables.}
\end{popfigure}

\courssub{Expérience de référence 2 : un animal en milieu confiné}

\begin{experience}[Survie d'un criquet (ou d'une souris) dans un milieu fermé]
\textbf{Protocole.} On place un criquet vivant dans un bocal A hermétiquement fermé, et un criquet identique dans un bocal B identique mais dont l'air est renouvelé en permanence (témoin). On peut aussi utiliser une souris dans une cloche de verre sur bac à eau de chaux.

\textbf{Observations.} Dans le bocal A, le criquet devient de moins en moins actif, puis meurt au bout d'un temps limité (quelques heures). Dans le bocal B, il survit normalement. L'analyse de l'air du bocal A (ou le test à l'eau de chaux sur le bac) montre une chute du taux de \ce{O2} et une élévation du taux de \ce{CO2}.

\textbf{Interprétation et conclusion.} L'animal consomme le dioxygène du milieu et rejette du dioxyde de carbone ; lorsque le \ce{O2} s'épuise, les cellules ne peuvent plus produire suffisamment d'énergie et l'organisme meurt. Le dioxygène est donc \textbf{indispensable à la vie} de la plupart des organismes (aérobies stricts).
\end{experience}

\begin{methode}[Tests d'identification des gaz : \ce{O2} et \ce{CO2} — Incontournables au Probatoire]
Pour exploiter des résultats expérimentaux, retiens les deux tests de caractérisation :
\begin{enumerate}
\item \textbf{Mise en évidence du dioxygène \ce{O2}} : on approche une \cle{bûchette incandescente} (ou une braise) de l'échantillon de gaz. Si la bûchette \textbf{se rallume} (flamme vive) ou que la braise \textbf{se ravive vivement}, le gaz contient du dioxygène.
\item \textbf{Mise en évidence du dioxyde de carbone \ce{CO2}} : on fait barboter le gaz dans de l'\cle{eau de chaux} (solution saturée d'hydroxyde de calcium \ce{Ca(OH)2}). Si l'eau de chaux \textbf{se trouble} (formation d'un précipité blanc de carbonate de calcium \ce{CaCO3}), le gaz contient du dioxyde de carbone.
\end{enumerate}
\textbf{Règle d'or :} Toujours comparer à un \textbf{témoin} (air ambiant, ou organisme tué/bouilli) pour conclure valablement que les variations sont dues à l'activité vitale.
\end{methode}

\courssub{Les échanges gazeux s'accompagnent d'un dégagement de chaleur}

\textbf{Observation.} Un tas de fumier ou de compost « chauffe » ; un silo de grains mal ventilé peut s'échauffer dangereusement et s'enflammer. D'où vient cette chaleur ?

\begin{experience}[Dégagement de chaleur par des graines germées (calorimétrie)]
\textbf{Protocole.} On place une masse connue de graines germées dans un vase calorimétrique (vase Dewar, à double paroi argentée sous vide, qui limite au maximum les échanges thermiques avec l'extérieur) muni d'un thermomètre précis. Un vase témoin identique contient la même masse de graines bouillies (mortes).

\textbf{Observations.} La température du vase contenant les graines vivantes s'élève de plusieurs degrés Celsius en quelques heures (ex. une élévation de \SI{4}{\celsius} en 24 h) ; celle du témoin reste constante (variation inférieure à \SI{0,1}{\celsius}).

\textbf{Conclusion.} L'activité vitale des graines \textbf{dégage de la chaleur}. Les échanges gazeux sont donc associés à une \cle{libération d'énergie} : les organismes dégradent des substances organiques de réserve (amidon, sucres, lipides) pour en extraire de l'énergie, dont une partie est dissipée sous forme de chaleur (rendement non nul).
\end{experience}

\courssub{Du constat global au mécanisme cellulaire : la notion de métabolisme énergétique}

Où se produisent ces échanges gazeux ? Des expériences complémentaires (mesure de la consommation de \ce{O2} par des tissus isolés — foie, muscle —, puis par des cellules isolées en culture) montrent que ce sont les \textbf{cellules} elles-mêmes qui consomment le \ce{O2} et produisent le \ce{CO2}. Chez un animal pluricellulaire, la respiration pulmonaire (ventilation) ne fait qu'acheminer le \ce{O2} de l'air alvéolaire vers le sang, puis vers les cellules, et évacuer le \ce{CO2} qu'elles rejettent.

La cellule fonctionne comme un moteur thermique microscopique : elle dégrade des \cle{molécules organiques} riches en énergie — principalement le \cle{glucose} (\ce{C6H12O6}), apporté par l'alimentation (d'où les boissons sucrées des footballeurs à la mi-temps !) — pour libérer l'énergie chimique stockée dans ses liaisons. L'ensemble de ces réactions de dégradation productrices d'énergie constitue le \cle{métabolisme énergétique}.

\begin{definition}[Respiration cellulaire]
La \cle{respiration cellulaire} est l'ensemble des réactions chimiques d'\textbf{oxydation complète du glucose} se déroulant dans les cellules \textbf{en présence de dioxygène} (voie aérobie). Elle produit du dioxyde de carbone, de l'eau et de l'\textbf{énergie}, stockée sous forme d'ATP. Son bilan simplifié (somme des trois étapes) s'écrit :
\[ \ce{C6H12O6 + 6O2 -> 6CO2 + 6H2O} + \text{énergie (environ 38 ATP)} \]
\end{definition}

\begin{attention}[Ne pas confondre respiration cellulaire et respiration pulmonaire — \textbf{Erreur éliminatoire au Probatoire !}]
\begin{itemize}
\item La \textbf{respiration pulmonaire} (ou ventilation pulmonaire) est un phénomène \textbf{macroscopique, mécanique} : ce sont les mouvements inspiratoires et expiratoires qui renouvellent l'air des alvéoles pulmonaires. Elle se voit, elle s'accélère à l'effort (fréquence et amplitude).
\item La \textbf{respiration cellulaire} est un phénomène \textbf{microscopique et biochimique} : ce sont des réactions chimiques d'oxydation du glucose \emph{à l'intérieur de chaque cellule}, dans le cytoplasme et les mitochondries. Elle est invisible à l'œil nu.
\end{itemize}
Les deux sont liées : la respiration pulmonaire fournit le \ce{O2} substrat et évacue le \ce{CO2} produit par la respiration cellulaire. Mais dans une copie, écrire « la cellule respire par les poumons » ou « les poumons produisent de l'ATP » est une \textbf{faute grave} qui coûte cher.
\end{attention}

\courssub{L'ATP : la « monnaie d'énergie » de la cellule}

\textbf{Intuition.} L'énergie libérée par l'oxydation du glucose ($\approx 2870\ \text{kJ/mol}$) est trop brutale, trop importante pour être utilisée directement par les micromachines cellulaires (pompes à ions, moteurs moléculaires, enzymes de synthèse). Comme on ne paie pas un pain au marché avec un lingot d'or, il faut une « monnaie » d'échange, utilisable en petites quantités, partout dans la cellule, de façon contrôlée. Cette monnaie, c'est l'\cle{ATP}.

\begin{definition}[ATP, ADP et hydrolyse de l'ATP]
L'\cle{ATP} (adénosine triphosphate) est un nucléotide formé d'une base azotée (adénine), d'un sucre (ribose) et de \textbf{trois groupes phosphate} liés entre eux par des liaisons anhydride à haut potentiel de transfert d'énergie ($\Delta G^\circ \approx -30,5\ \text{kJ/mol}$ en conditions standard). Elle joue le rôle de \textbf{vecteur d'énergie} universel de la cellule : l'énergie libérée par la dégradation du glucose est captée transitoirement dans l'ATP (phosphorylation), puis restituée là où la cellule en a besoin.

L'énergie est libérée par \textbf{hydrolyse} de l'ATP (réaction exergonique) :
\[ \ce{ATP + H2O -> ADP + Pi} + \text{énergie (utilisable)} \]
où \cle{ADP} est l'adénosine diphosphate et \cle{Pi} un phosphate inorganique ($\ce{HPO4^{2-}}$). La réaction inverse (resynthèse de l'ATP à partir d'ADP et \ce{Pi}, réaction endergonique) est couplée aux réactions exergoniques de la respiration cellulaire : le couple \ce{ATP/ADP} fonctionne comme une \textbf{batterie rechargeable} en permanence.
\end{definition}

\begin{savaistu}[Des millions d'ATP par seconde !]
Pendant un effort intense, une cellule musculaire humaine peut hydrolyser et resynthétiser jusqu'à \textbf{10 millions de molécules d'ATP par seconde} ! Or le stock d'ATP libre d'un muscle n'assure que \textbf{2 à 3 secondes} d'effort maximal : il faut donc le régénérer en permanence. Cela explique l'accélération immédiate de la respiration pulmonaire (apport massif de \ce{O2}) et les crampes du sprinteur de notre situation d'accroche : ses cellules, privées d'oxygène suffisant par rapport à la demande brutale, basculent vers un autre mode de production d'énergie, moins rentable mais plus rapide et anaérobie : la \textbf{fermentation lactique}, que nous étudierons dans la suite de la leçon.
\end{savaistu}

\courssub{Le Quotient Respiratoire (QR) : un outil d'exploitation quantitative}

\begin{propriete}[Quotient Respiratoire (QR)]
Le \cle{Quotient Respiratoire (QR)} est le rapport entre le volume de dioxyde de carbone produit et le volume de dioxygène consommé sur une même durée, dans les mêmes conditions de température et de pression :
\[ \text{QR} = \frac{V_{\ce{CO2}}\ \text{produit}}{V_{\ce{O2}}\ \text{consommé}} = \frac{n_{\ce{CO2}}}{n_{\ce{O2}}} \]
Il permet d'identifier le substrat respiratoire principal oxydé par la cellule :
\begin{center}
\begin{tabularx}{\linewidth}{|l|X|c|}\hline
\textbf{Substrat principal} & \textbf{Bilan d'oxydation simplifié} & \textbf{QR théorique} \\ \hline
Glucides (glucose) & $\ce{C6H12O6 + 6O2 -> 6CO2 + 6H2O}$ & $6/6 = \mathbf{1,0}$ \\ \hline
Lipides (acides gras) & $\ce{C57H104O6 + 80O2 -> 57CO2 + 52H2O}$ (ex. tripalmitine) & $\approx 57/80 = \mathbf{0,71}$ \\ \hline
Protéines (moyenne) & Composition variable & $\approx \mathbf{0,80\text{--}0,85}$ \\ \hline
\end{tabularx}
\end{center}
\textbf{Application :} Un QR mesuré proche de 1 indique une oxydation quasi exclusive de glucides (jeune sportif à l'effort, graines germées riches en amidon). Un QR proche de 0,7 indique une oxydation lipidique (jeûne prolongé, migration d'oiseaux, effort d'endurance modéré).
\end{propriete}

\begin{exoresolu}[Exploitation de données expérimentales — Type Probatoire]
\textbf{Énoncé.} On place 50 g de graines germées dans un flacon fermé de 1 L contenant initialement 21\% de \ce{O2} et 0,03\% de \ce{CO2}. Après 24 h, on mesure 15\% de \ce{O2} et 6\% de \ce{CO2}. Un flacon témoin contenant 50 g de graines bouillies ne montre aucune variation.
\begin{enumerate}
\item Calculer les variations des taux de \ce{O2} et de \ce{CO2} dans le flacon expérimental.
\item Calculer le Quotient Respiratoire (QR) approximatif. En déduire le substrat principal respiré.
\item Que peut-on conclure sur l'activité des graines ?
\item Quel est le rôle du flacon témoin ?
\end{enumerate}
\tcblower
\textbf{Solution détaillée.}
\begin{enumerate}
\item Variation du taux de \ce{O2} : $\Delta[\ce{O2}] = 15 - 21 = -6$ points de pourcentage (diminution de $6\%$).\\
Variation du taux de \ce{CO2} : $\Delta[\ce{CO2}] = 6 - 0,03 = +5,97 \approx +6$ points de pourcentage (augmentation de $\approx 6\%$).
\item Comme le volume et la température sont constants, les variations de pourcentages sont proportionnelles aux variations de quantités de matière (loi des gaz parfaits).
\[ \text{QR} = \frac{\Delta n_{\ce{CO2}}}{\Delta n_{\ce{O2}}} \approx \frac{+6}{-6} = \mathbf{1,0} \]
(On prend les valeurs absolues pour le rapport). Un QR = 1,0 correspond à l'oxydation quasi exclusive de \textbf{glucides} (amidon/sucres de réserve des graines).
\item Les graines germées vivantes ont \textbf{consommé du dioxygène} et \textbf{rejeté du dioxyde de carbone} en quantités molaires égales (QR=1) : elles réalisent une respiration cellulaire aérobie utilisant leurs réserves glucidiques pour produire de l'énergie (ATP).
\item Le flacon témoin (graines bouillies, donc mortes, mêmes conditions physiques) permet de vérifier que les variations observées sont bien dues à l'\textbf{activité vitale (métabolisme)} des graines et non à un artefact du dispositif (fuite, dissolution des gaz dans l'humidité, réaction chimique abiotique). Sans témoin, aucune conclusion scientifique valable n'est possible.
\end{enumerate}
\end{exoresolu}

\begin{aretenir}[L'essentiel de la section — À connaître par cœur]
\begin{itemize}
\item Tout organisme vivant aérobie \textbf{consomme du \ce{O2}} et \textbf{rejette du \ce{CO2}} : ce sont les échanges gazeux cellulaires, mis en évidence par le test à la bûchette incandescente (\ce{O2}) et à l'eau de chaux (\ce{CO2}).
\item Ces échanges s'accompagnent d'un \textbf{dégagement de chaleur} mesurable (calorimétrie) : la cellule dégrade le glucose pour libérer de l'énergie.
\item L'énergie libérée est stockée dans l'\textbf{ATP} (adénosine triphosphate), monnaie énergétique universelle : hydrolyse $\ce{ATP + H2O -> ADP + Pi} + \text{énergie}$.
\item La \textbf{respiration cellulaire} (biochimique, mitochondriale) a pour bilan global : $\ce{C6H12O6 + 6O2 -> 6CO2 + 6H2O}$ ($\approx 38$ ATP).
\item \textbf{Ne jamais confondre} respiration \textbf{cellulaire} (microscopique, productrice d'ATP) et respiration \textbf{pulmonaire} (macroscopique, ventilation, transport des gaz).
\item Le \textbf{Quotient Respiratoire (QR = $V_{\ce{CO2}}/V_{\ce{O2}}$)} permet d'identifier le substrat : QR=1 (glucides), QR$\approx$0,7 (lipides), QR$\approx$0,8 (protéines).
\end{itemize}
\end{aretenir}

\coursec{La mitochondrie, centrale énergétique de la cellule}

Dans la section précédente, nous avons mis en évidence les échanges gazeux cellulaires : toute cellule vivante consomme du \ce{O2} et rejette du \ce{CO2}, comme l'a montré l'expérience des graines germées ou celle du sportif dont la respiration s'accélère à l'effort. Mais une question fondamentale demeure : \emph{où}, précisément, à l'intérieur de la cellule, le dioxygène est-il utilisé ? Et \emph{comment} l'énergie contenue dans le \ce{glucose} est-elle transformée en énergie utilisable, l'\ce{ATP} ? La réponse tient en un mot : la \cle{mitochondrie}. C'est elle que nous allons explorer maintenant, de sa localisation à sa structure ultrafine, pour comprendre pourquoi on la surnomme la « centrale énergétique » de la cellule.

\courssub{Localisation et abondance des mitochondries : une observation qui pose question}

\textbf{Une situation concrète pour démarrer.}\par
Observe un match de football au stade Ahmadou Ahidjo de Yaoundé : après 90 minutes de course intense, les joueurs sont épuisés, alors qu'un spectateur resté assis n'a ressenti aucune fatigue musculaire. De même, le \ce{c\oe ur}, qui bat environ \num{100000} fois par jour sans jamais s'arrêter, ne se fatigue pas. Qu'est-ce qui distingue les cellules de ces organes infatigables ?

\textbf{Démarche d'investigation.}\par
Des biologistes ont comparé, au microscope électronique, des coupes de différents types cellulaires et compté les mitochondries présentes :

\begin{center}
\begin{tabularx}{\linewidth}{|l|X|}\hline
\rowcolor{popblueL} \textbf{Type de cellule} & \textbf{Nombre approximatif de mitochondries} \\ \hline
Cellule musculaire squelettique & 300 à 400 \\ \hline
Cellule musculaire cardiaque & 1 000 à 2 000 (jusqu'à 35 \% du volume cellulaire) \\ \hline
Cellule hépatique (foie) & 1 000 à 2 500 \\ \hline
Neurone & plusieurs centaines, concentrées près des synapses \\ \hline
Hématie (globule rouge humain) & 0 \\ \hline
\end{tabularx}
\end{center}

\textbf{Interprétation.} Les cellules les plus consommatrices d'énergie (muscles, foie, neurones) sont aussi les plus riches en mitochondries. À l'inverse, l'hématie, qui n'utilise pas le dioxygène qu'elle transporte, en est totalement dépourvue : elle ne peut réaliser que la glycolyse, dans son cytoplasme.

\textbf{Conclusion.} Il existe une \cle{corrélation} entre l'abondance des mitochondries et les besoins énergétiques de la cellule : la mitochondrie est bien l'organite de la production d'énergie.

\begin{definition}[Mitochondrie]
La \cle{mitochondrie} est un organite du \cle{cytoplasme} des cellules eucaryotes, de forme généralement ovoïde ou bâtonnet (longueur de 1 à \SI{10}{\micro m}, diamètre de 0,5 à \SI{1}{\micro m}), entourée d'une \textbf{double membrane}. Elle est le siège des étapes oxydatives de la \cle{respiration cellulaire}, c'est-à-dire de la production de la majeure partie de l'\cle{ATP} (adénosine triphosphate), la « molécule d'énergie » de la cellule.
\end{definition}

\begin{exemplebox}[Où trouve-t-on beaucoup de mitochondries ?]
\begin{itemize}
\item Les \textbf{cellules musculaires} : elles doivent fournir un travail mécanique intense et rapide (contraction).
\item Le \textbf{foie} : il réalise des milliers de réactions chimiques (détoxication, synthèses) très coûteuses en énergie.
\item Les \textbf{neurones} : ils maintiennent en permanence des pompes ioniques (pompe \ce{Na+}/\ce{K+}) indispensables à la transmission de l'influx nerveux.
\item Les \textbf{cellules des trompes de Fallope} et les \textbf{spermatozoïdes} : leurs mouvements (cils, flagelle) exigent beaucoup d'\ce{ATP} ; chez le spermatozoïde, les mitochondries s'enroulent autour de la base du flagelle.
\end{itemize}
\end{exemplebox}

\courssub{Observation au microscope électronique : la structure de la mitochondrie}

La mitochondrie est trop petite (environ \SI{1}{\micro m}) pour être observée en détail au microscope optique : sa structure fine n'a été révélée que par le \cle{microscope électronique à transmission} (MET), qui atteint des grossissements de plusieurs centaines de milliers de fois.

\textbf{Description d'un cliché de mitochondrie au MET.}\par
Sur une coupe longitudinale, on observe :
\begin{itemize}
\item un organite ovoïde limité par \textbf{deux membranes} nettement distinctes ;
\item une \textbf{membrane externe lisse}, régulière, qui isole l'organite du cytoplasme ;
\item une \textbf{membrane interne} qui se replie profondément vers l'intérieur, formant des lames ou des tubules appelés \cle{crêtes mitochondriales} ;
\item un espace intérieur sombre et homogène, la \cle{matrice}, dans lequel on distingue parfois de petits grains (ribosomes) et un filament circulaire (l'ADN mitochondrial) ;
\item un espace étroit entre les deux membranes, l'\cle{espace intermembranaire}.
\end{itemize}

\begin{definition}[Les compartiments de la mitochondrie]
La mitochondrie comprend \textbf{deux membranes} et \textbf{deux compartiments} :
\begin{itemize}
\item la \cle{membrane externe}, lisse, perméable aux petites molécules ;
\item la \cle{membrane interne}, repliée en \cle{crêtes mitochondriales}, qui porte les enzymes de la chaîne respiratoire ;
\item la \cle{matrice mitochondriale}, le compartiment interne délimité par la membrane interne : elle contient les enzymes du cycle de Krebs, l'ADN mitochondrial et des ribosomes ;
\item l'\cle{espace intermembranaire}, l'espace étroit situé entre la membrane externe et la membrane interne.
\end{itemize}
\end{definition}

\begin{popfigure}
\begin{center}
\begin{tikzpicture}[scale=1.0]
% Membrane externe (forme ovoïde ondulée)
\draw[popblue, very thick, rounded corners=14pt] (0,0) .. controls (0.3,1.6) and (1.5,2.1) .. (3.2,2.0)
 .. controls (5.2,1.9) and (7.2,2.1) .. (8.6,1.2)
 .. controls (9.6,0.5) and (9.6,-0.9) .. (8.4,-1.5)
 .. controls (6.8,-2.2) and (4.6,-2.0) .. (3.0,-2.0)
 .. controls (1.4,-2.0) and (0,-1.4) .. (0,0) -- cycle;
% Membrane interne avec crêtes
\draw[poporange, very thick, rounded corners=10pt] (0.5,0.1) .. controls (0.8,1.2) and (1.7,1.6) .. (3.2,1.55)
 .. controls (5.0,1.5) and (6.9,1.65) .. (8.0,0.95)
 .. controls (8.8,0.45) and (8.8,-0.7) .. (7.9,-1.15)
 .. controls (6.5,-1.75) and (4.6,-1.55) .. (3.0,-1.55)
 .. controls (1.6,-1.55) and (0.4,-1.1) .. (0.5,0.1) -- cycle;
% Crêtes (replis de la membrane interne)
\draw[poporange, very thick] (1.6,1.5) -- (1.6,0.55) .. controls (1.6,0.2) and (2.3,0.2) .. (2.3,0.55) -- (2.3,1.5);
\draw[poporange, very thick] (3.6,-1.5) -- (3.6,-0.5) .. controls (3.6,-0.15) and (4.3,-0.15) .. (4.3,-0.5) -- (4.3,-1.5);
\draw[poporange, very thick] (5.4,1.55) -- (5.4,0.6) .. controls (5.4,0.25) and (6.1,0.25) .. (6.1,0.6) -- (6.1,1.55);
\draw[poporange, very thick] (6.9,-1.6) -- (6.9,-0.6) .. controls (6.9,-0.25) and (7.6,-0.25) .. (7.6,-0.6) -- (7.6,-1.45);
% Matrice : ADN et ribosomes
\draw[popgreen, thick] (3.0,0.6) circle (0.28);
\node[popgreen] at (3.0,0.6) {\scriptsize ADN};
\fill[popgreen] (4.9,-0.6) circle (0.07);
\fill[popgreen] (5.2,-0.35) circle (0.07);
\fill[popgreen] (2.6,-0.7) circle (0.07);
% Légendes (traits horizontaux)
\draw[thin, popdark] (8.9,1.6) -- (11.2,1.6) node[right, align=left] {\scriptsize Membrane externe (lisse)};
\draw[thin, popdark] (8.15,0.95) -- (11.2,0.95) node[right, align=left] {\scriptsize Membrane interne};
\draw[thin, popdark] (6.1,1.9) -- (11.2,2.3) node[right, align=left] {\scriptsize Crête mitochondriale};
\draw[thin, popdark] (8.5,-0.4) -- (11.2,-0.4) node[right, align=left] {\scriptsize Espace intermembranaire};
\draw[thin, popdark] (0.6,-1.2) -- (-1.4,-1.6) node[left, align=right] {\scriptsize Matrice mitochondriale\\ \scriptsize (cycle de Krebs)};
\draw[thin, popdark] (3.28,0.6) -- (-1.4,0.6) node[left, align=right] {\scriptsize ADN mitochondrial\\ \scriptsize + ribosomes};
\node[align=center, poppurple] at (4.7,-2.6) {\scriptsize \textbf{Chaîne respiratoire : sur la membrane interne (crêtes)}};
\end{tikzpicture}
\end{center}
\legende{Schéma d'une mitochondrie en coupe longitudinale (observation au microscope électronique). La membrane externe est lisse ; la membrane interne forme des crêtes qui délimitent la matrice. L'espace intermembranaire sépare les deux membranes.}
\end{popfigure}

\begin{methode}[Légender correctement un schéma de mitochondrie à l'examen]
Au Probatoire, la légende d'un schéma est souvent notée sur plusieurs points. Applique rigoureusement ces étapes :
\begin{enumerate}
\item \textbf{Donne un titre} au schéma, en dessous ou au-dessus : « Mitochondrie en coupe longitudinale (M.E.T.) », en précisant la technique d'observation.
\item \textbf{Trace des traits de légende horizontaux}, à la règle, qui partent exactement de la structure désignée (jamais de flèches obliques croisées, jamais de trait qui traverse une autre structure).
\item \textbf{Écris les noms à l'horizontale}, alignés de part et d'autre du schéma, sans abréviation fantaisiste.
\item \textbf{Sois précis} : écris « membrane interne repliée en crêtes » et non simplement « membrane » ; « matrice mitochondriale » et non « intérieur ».
\item \textbf{Ne légende que ce qui est visible} sur le schéma ou le cliché : n'ajoute pas « ATP » ou « glucose », qui sont des molécules invisibles au microscope.
\end{enumerate}
\end{methode}

\courssub{Relation structure-fonction : pourquoi des crêtes ?}

\textbf{Une analogie pour comprendre.}\par
Pour chauffer une maison à Garoua pendant la saison sèche, on n'utilise pas un simple fil électrique tendu : on utilise un radiateur à \emph{ailettes}. Pourquoi ? Parce que les ailettes augmentent considérablement la \textbf{surface d'échange} entre le radiateur et l'air, sans augmenter l'encombrement de l'appareil. De même, les intestins ne sont pas un tube lisse mais un tube tapissé de millions de \emph{villosités} : la surface d'absorption des nutriments est multipliée par 30 environ.

\textbf{Application à la mitochondrie.} Les \cle{crêtes mitochondriales} jouent exactement ce rôle : en se repliant, la membrane interne \textbf{multiplie sa surface} à volume constant. Or, c'est sur cette membrane interne que sont fixées les \textbf{enzymes de la chaîne respiratoire} et les complexes de synthèse de l'\ce{ATP} (\ce{ATP} synthétase). Plus la surface est grande, plus la cellule peut installer de « machines » à fabriquer de l'\ce{ATP}, et plus sa production énergétique est élevée.

\begin{propriete}[La membrane interne, siège majeur de la synthèse d'ATP (admise)]
La \cle{membrane interne} des mitochondries est le siège de la synthèse de la \textbf{majorité de l'ATP} produit par la cellule (environ 90 \%). Elle porte les enzymes de la \cle{chaîne respiratoire} (transporteurs d'électrons) et l'\cle{ATP synthétase}, l'enzyme qui fabrique l'\ce{ATP}. Cette propriété est \textbf{admise} au niveau Première : sa démonstration n'est pas exigible au Probatoire, mais tu dois savoir l'énoncer et la relier à la présence des crêtes.
\end{propriete}

La conséquence logique est remarquable : les cellules très consommatrices d'énergie possèdent des mitochondries dont la membrane interne est \emph{particulièrement riche en crêtes}. Une mitochondrie de cellule cardiaque présente des crêtes denses et nombreuses ; celle d'une cellule peu active en possède peu. La structure épouse la fonction.

\courssub{Répartition des étapes de la respiration cellulaire : qui fait quoi, et où ?}

Nous pouvons maintenant cartographier précisément les trois étapes de la respiration cellulaire (que nous détaillerons dans la section suivante) :

\begin{center}
\begin{tabularx}{\linewidth}{|l|X|X|}\hline
\rowcolor{popblueL} \textbf{Étape} & \textbf{Localisation exacte} & \textbf{Rôle essentiel} \\ \hline
Glycolyse & \textbf{Cytoplasme} (hyaloplasme), hors mitochondrie & Dégradation du \ce{glucose} en \ce{pyruvate} ; faible gain en \ce{ATP} \\ \hline
Cycle de Krebs & \textbf{Matrice mitochondriale} & Oxydation complète ; production de \ce{CO2} et de transporteurs d'électrons chargés \\ \hline
Chaîne respiratoire & \textbf{Membrane interne} (crêtes) & Consommation de \ce{O2} ; synthèse de la majorité de l'\ce{ATP} \\ \hline
\end{tabularx}
\end{center}

Cette répartition explique les échanges gazeux observés dans la section 1 : le \ce{CO2} rejeté par la cellule naît dans la \textbf{matrice} (cycle de Krebs), et le \ce{O2} consommé est fixé sur la \textbf{membrane interne} (chaîne respiratoire), où il est réduit en eau.

\begin{attention}[Erreur classique au Probatoire : la glycolyse n'a PAS lieu dans la mitochondrie]
De nombreux candidats écrivent : « la glycolyse se déroule dans la mitochondrie ». C'est \textbf{faux}. La glycolyse se déroule dans le \cle{cytoplasme} (hyaloplasme), à l'extérieur de la mitochondrie. Seules le \textbf{cycle de Krebs} (matrice) et la \textbf{chaîne respiratoire} (membrane interne) sont mitochondriales. Retiens le schéma logique : le \ce{glucose} entre dans la cellule, est découpé dans le cytoplasme, et seuls ses produits (\ce{pyruvate}) pénètrent dans la mitochondrie pour y être oxydés complètement. Preuve indirecte : l'hématie, \textbf{dépourvue de mitochondries}, réalise pourtant la glycolyse — mais elle ne peut pas aller plus loin.
\end{attention}

\courssub{Abondance des mitochondries et entraînement : une comparaison parlante}

L'observation de coupes de muscle au microscope électronique, chez un sujet sédentaire et chez un athlète entraîné, montre une différence spectaculaire : les fibres musculaires du sportif contiennent nettement plus de mitochondries, et celles-ci sont plus riches en crêtes. C'est l'une des bases biologiques de l'amélioration des performances par l'entraînement, que nous étudierons en détail dans la section 5.

\begin{popfigure}
\begin{center}
\begin{tikzpicture}[scale=0.95]
% Cellule non entraînée
\draw[popblue, thick, rounded corners=8pt] (0,0) rectangle (5,2.6);
\foreach \x/\y in {0.8/0.7, 1.8/1.6, 2.9/0.9, 4.0/1.8, 3.9/0.5}{
  \draw[poporange, thick] (\x,\y) ellipse (0.35 and 0.18);
}
\node[align=center] at (2.5,-0.6) {\scriptsize \textbf{Cellule musculaire non entraînée}\\ \scriptsize (environ 5 mitochondries par champ)};
% Cellule entraînée
\draw[popblue, thick, rounded corners=8pt] (6.5,0) rectangle (11.5,2.6);
\foreach \x/\y in {7.1/0.6, 7.3/1.5, 8.0/2.1, 8.2/0.9, 8.9/1.7, 9.0/0.4, 9.7/1.1, 9.9/2.2, 10.5/0.6, 10.6/1.6, 11.0/2.3, 7.6/2.3}{
  \draw[poporange, thick] (\x,\y) ellipse (0.32 and 0.16);
}
\node[align=center] at (9,-0.6) {\scriptsize \textbf{Cellule musculaire entraînée}\\ \scriptsize (mitochondries plus nombreuses et plus riches en crêtes)};
% Légende
\draw[poporange, thick] (3.2,3.3) ellipse (0.32 and 0.16);
\draw[thin, popdark] (3.55,3.3) -- (4.2,3.3) node[right] {\scriptsize Mitochondrie};
\end{tikzpicture}
\end{center}
\legende{Comparaison schématique de deux fibres musculaires observées au microscope électronique : l'entraînement augmente le nombre de mitochondries par cellule, donc la capacité de production d'ATP.}
\end{popfigure}

\begin{exoresolu}[Exploiter des données de microscopie électronique]
\textbf{Énoncé.} On compare au microscope électronique des coupes de fibres musculaires de deux sujets : un sujet A sédentaire et un sujet B, coureur de fond entraîné depuis 5 ans. On compte en moyenne 6 mitochondries par coupe de fibre chez A, contre 15 chez B ; de plus, les mitochondries de B présentent environ deux fois plus de crêtes. (1) Quelle hypothèse peut-on formuler pour expliquer ces différences ? (2) Quelle conséquence ces différences ont-elles sur les capacités physiques des deux sujets ? Justifier.
\tcblower
\textbf{Solution détaillée.}
(1) \textbf{Hypothèse} : l'entraînement régulier stimule la multiplication des mitochondries dans les fibres musculaires (\cle{biogenèse mitochondriale}) et l'accroissement de leur membrane interne (davantage de crêtes). Les cellules musculaires du sujet B, sollicitées quotidiennement, ont adapté leur « parc énergétique » à la demande.
(2) \textbf{Conséquence} : les mitochondries produisent l'\ce{ATP} par la respiration cellulaire, et la majorité de cet \ce{ATP} est synthétisée sur la membrane interne (crêtes). Le sujet B possède donc, à masse musculaire égale, une \textbf{capacité de production d'ATP nettement supérieure} : ses muscles peuvent maintenir un effort intense plus longtemps sans épuiser leurs réserves ni basculer prématurément vers la fermentation lactique (source de fatigue et de crampes). Il présente donc une meilleure \textbf{endurance} que le sujet A. C'est la preuve que le rendement énergétique de la cellule n'est pas figé : il dépend de facteurs comme l'entraînement.
\end{exoresolu}

\courssub{La mitochondrie, un organite semi-autonome : ADN et ribosomes propres}

Un détail de l'observation au MET intrigue les biologistes depuis longtemps : la matrice mitochondriale contient de \textbf{petits grains semblables à des ribosomes} et un \textbf{filament d'ADN circulaire}. La mitochondrie possède donc son \cle{propre matériel génétique} et fabrique elle-même une partie de ses protéines. Elle peut aussi se diviser en deux, indépendamment de la division de la cellule. On dit qu'elle est \cle{semi-autonome}.

\begin{savaistu}[Les mitochondries descendent-elles de bactéries ?]
La \cle{théorie endosymbiotique}, proposée notamment par la biologiste américaine Lynn Margulis et aujourd'hui \textbf{admise} par la communauté scientifique, suggère que les mitochondries descendent de \textbf{bactéries libres} capables de respirer, qu'une cellule ancestrale aurait \emph{englouties sans les digérer} il y a environ \textbf{1,5 milliard d'années}. Plusieurs indices soutiennent cette hypothèse : l'ADN mitochondrial est \textbf{circulaire} comme celui des bactéries (alors que l'ADN du noyau est linéaire) ; les ribosomes mitochondriaux ressemblent aux ribosomes bactériens ; la mitochondrie possède une \textbf{double membrane} (la membrane externe correspondrait à la membrane de la vacuole qui a englouti la bactérie) ; enfin, elle se divise par \textbf{scissiparité}, comme les bactéries. Cette association ancienne a été si profitable — la bactérie fournissait de l'énergie, la cellule hôte offrait protection et nutriments — qu'elle est devenue définitive. Autre fait remarquable : tu as hérité de \emph{toutes} tes mitochondries de ta \textbf{mère}, car seul l'ovule transmet ses mitochondries à l'embryon ; c'est pourquoi l'ADN mitochondrial sert en généalogie à retracer les lignées maternelles.
\end{savaistu}

\begin{aretenir}[L'essentiel de la section]
\begin{itemize}
\item La \cle{mitochondrie} est un organite du cytoplasme, très abondant dans les cellules consommatrices d'énergie (muscles, foie, neurones) et absente des hématies.
\item Elle est limitée par une \textbf{membrane externe lisse} et une \textbf{membrane interne repliée en crêtes}, qui délimitent deux compartiments : la \textbf{matrice} et l'\textbf{espace intermembranaire}.
\item Les crêtes \textbf{augmentent la surface} de la membrane interne, qui porte les enzymes de la chaîne respiratoire : c'est là qu'est synthétisée la majorité de l'\ce{ATP}.
\item Répartition des étapes : \textbf{glycolyse} dans le cytoplasme, \textbf{cycle de Krebs} dans la matrice, \textbf{chaîne respiratoire} sur la membrane interne.
\item La mitochondrie possède son propre \textbf{ADN circulaire} et ses propres \textbf{ribosomes} : elle est semi-autonome, héritage probable d'une bactérie endosymbiotique.
\end{itemize}
\end{aretenir}

Maintenant que nous connaissons l'usine, ses ateliers et ses machines, il reste à suivre le cheminement complet de la matière première — le \ce{glucose} — jusqu'au produit fini, l'\ce{ATP}. C'est l'objet de la section suivante : les trois étapes de la respiration cellulaire.

\coursec{Les trois étapes de la respiration cellulaire}

Nous venons de voir que la mitochondrie est la centrale énergétique de la cellule : c'est en elle que s'achève la dégradation du glucose et que se produit la majeure partie de l'ATP. Mais une question demeure : comment, concrètement, une molécule de glucose est-elle transformée en énergie utilisable ? La réponse tient en trois étapes enchaînées, comme trois ateliers d'une même usine : la \cle{glycolyse}, le \cle{cycle de Krebs} et la \cle{chaîne respiratoire}. C'est l'objet de cette section.

\courssub{Vue d'ensemble : le devenir d'une molécule de glucose}

\begin{exemplebox}[Une situation concrète pour comprendre]
Un athlète camerounais court le 100 mètres au stade Ahmadou Ahidjo. Pendant l'effort, ses muscles consomment énormément d'ATP. D'où vient-il ? Du glucose stocké dans ses cellules musculaires (sous forme de glycogène). Mais ce glucose ne « brûle » pas d'un coup : il est dégradé progressivement, étape par étape, pour récupérer un maximum d'énergie sans « gaspillage » en chaleur. C'est exactement la stratégie de la respiration cellulaire.
\end{exemplebox}

La dégradation complète d'une molécule de glucose suit un chemin précis, en trois étapes principales, précédées d'une étape de liaison décisive :

\begin{enumerate}
\item La \cle{glycolyse}, qui se déroule dans le \textbf{cytoplasme} (hyaloplasme), en \textbf{absence d'oxygène} : le glucose (6 atomes de carbone) est scindé en deux molécules de \textbf{pyruvate} (3 atomes de carbone chacune).
\item La \cle{décarboxylation oxydative du pyruvate} (ou \emph{réaction de liaison}), qui se déroule dans la \textbf{matrice mitochondriale} : chaque pyruvate perd un carbone (\ce{CO2}) et devient un \textbf{acétyl-CoA} (2 carbones), en réduisant un \ce{NAD+} en \ce{NADH}.
\item Le \cle{cycle de Krebs}, qui se déroule dans la \textbf{matrice mitochondriale} : l'acétyl-CoA y est totalement oxydé, avec dégagement de \ce{CO2} et production de molécules réduites (\ce{NADH}, \ce{FADH2}) et d'ATP.
\item La \cle{chaîne respiratoire}, qui se déroule sur la \textbf{membrane interne} de la mitochondrie (crêtes) : c'est là que l'oxygène intervient enfin, comme accepteur final d'électrons, et que se produit la grande majorité de l'ATP par phosphorylation oxydative.
\end{enumerate}

\begin{popfigure}
\begin{center}
\begin{tikzpicture}[font=\small, >=Stealth, node distance=0.8cm and 1.2cm]
% Styles
\tikzset{
  boxmain/.style={draw, rounded corners=4pt, minimum height=1cm, align=center, font=\small},
  boxstep/.style={boxmain, fill=popgreenL, minimum width=3cm, font=\bfseries\small},
  boxmol/.style={boxmain, fill=popgoldL, minimum width=2.5cm},
  boxloc/.style={font=\footnotesize, text=popmuted, anchor=north},
  arrowmain/.style={->, very thick, popdark},
  arrowred/.style={->, thick, popink},
}
% --- CYTOPLASME ---
\node[boxmol, align=center] (glu){Glucose \\(6 C)};
\node[boxstep, right=of glu] (gly) {Glycolyse};
\node[boxloc, below=2pt of gly, align=center]{Cytoplasme \\ 2 ATP nets + 2 NADH};
\node[boxmol, right=of gly, align=center] (pyr){2 Pyruvate \\(3 C)};

% --- MITOCHONDRIE (Matrice) ---
\node[boxstep, below right=1.5cm and 0.5cm of pyr, align=center] (link){Décarboxylation\\ oxydative du pyruvate};
\node[boxloc, below=2pt of link, align=center]{Matrice \\ 2 NADH + 2 CO$_2$};
\node[boxmol, right=of link, align=center] (accoa){2 Acétyl-CoA \\(2 C)};
\node[boxstep, below=of accoa] (krebs) {Cycle de Krebs};
\node[boxloc, below=2pt of krebs, align=center]{Matrice \\ 2 ATP + 6 NADH + 2 FADH$_2$ + 4 CO$_2$};

% --- MITOCHONDRIE (Membrane interne / Chaîne respiratoire) ---
\node[boxstep, below left=1.5cm and 1cm of krebs] (chain) {Chaîne respiratoire};
\node[boxloc, below=2pt of chain, align=center]{Membrane interne \\ O$_2$ + 4e$^-$ + 4H$^+$ $\rightarrow$ 2 H$_2$O \\ $\approx$ 32--34 ATP};

% Flèches principales
\draw[arrowmain] (glu) -- (gly);
\draw[arrowmain] (gly) -- (pyr);
\draw[arrowmain] (pyr) -- (link);
\draw[arrowmain] (link) -- (accoa);
\draw[arrowmain] (accoa) -- (krebs);
\draw[arrowmain] (krebs) -- (chain);

% Flèches NADH/FADH2 vers chaîne (représentation schématique)
\draw[arrowred, dashed] (gly.east) .. controls +(1,0) and +(0,1) .. (chain.north) node[midway, above, sloped, font=\footnotesize] {2 NADH (navettes)};
\draw[arrowred, dashed] (link.east) .. controls +(0.5,0) and +(0,1) .. (chain.north) node[midway, above, sloped, font=\footnotesize] {2 NADH};
\draw[arrowred, dashed] (krebs.west) .. controls +(-0.5,0) and +(0,-0.5) .. (chain.north) node[midway, left, font=\footnotesize] {6 NADH + 2 FADH$_2$};

% Cadre Mitochondrie
\begin{scope}[on background layer]
\draw[fill=popblueL!15, rounded corners=15pt, thick, popblue] 
  ($(link.north west)+(-0.5,0.5)$) rectangle 
  ($(chain.south east)+(0.5,-0.5)$);
\node[popblue, anchor=north west, font=\bfseries\small] at ($(link.north west)+(-0.5,0.6)$) {Mitochondrie};
\node[popmuted, font=\footnotesize] at ($(accoa)!0.5!(krebs)+(0,0.8)$) {Matrice};
\node[popmuted, font=\footnotesize] at ($(chain)+(0,-1.2)$) {Membrane interne (crêtes)};
\end{scope}
\end{tikzpicture}
\end{center}
\legende{Schéma synoptique de la respiration cellulaire : le glucose est dégradé en quatre phases — glycolyse (cytoplasme), décarboxylation oxydative du pyruvate (matrice), cycle de Krebs (matrice) et chaîne respiratoire (membrane interne) — avec production d'ATP et de molécules réduites (NADH, FADH$_2$) à chaque étape.}
\end{popfigure}

\begin{aretenir}[L'idée directrice]
La respiration cellulaire est une \textbf{oxydation ménagée} : au lieu de brûler le glucose d'un seul coup (ce qui dissiperait l'énergie en chaleur), la cellule le dégrade par petites étapes et récupère l'énergie sous deux formes : directement en \textbf{ATP} (glycolyse, cycle de Krebs) et sous forme de \textbf{molécules réduites} (\textbf{NADH} et \textbf{FADH$_2$}), véritables « navettes » d'électrons à haut potentiel énergétique qui alimenteront la chaîne respiratoire.
\end{aretenir}

\courssub{Première étape : la glycolyse}

\begin{definition}[Glycolyse]
La \cle{glycolyse} est une suite de dix réactions enzymatiques qui se déroule dans le \textbf{cytoplasme} de toutes les cellules, \textbf{sans utiliser d'oxygène} (voie anaérobie). Au cours de la glycolyse, une molécule de glucose (6 atomes de carbone) est scindée en \textbf{deux molécules de pyruvate} (3 atomes de carbone chacune).
\end{definition}

Le mot vient du grec \emph{glykus} (sucre) et \emph{lusis} (rupture) : c'est littéralement la « cassure du sucre ». Cette voie est très ancienne dans l'histoire du vivant : elle est commune à presque tous les organismes, ce qui témoigne de son importance fondamentale.

\textbf{Bilan de la glycolyse pour une molécule de glucose :}
\begin{itemize}
\item \textbf{2 molécules de pyruvate} (qui entreront ensuite dans les mitochondries si l'oxygène est présent) ;
\item \textbf{2 ATP nets} (la glycolyse consomme 2 ATP en phase d'amorçage puis en produit 4 en phase de rendement, d'où un gain net de 2) ;
\item \textbf{2 molécules de NADH} (réduction de 2 \ce{NAD+}) qui transportent des électrons de haute énergie vers la chaîne respiratoire.
\end{itemize}

\begin{attention}[L'oxygène n'intervient pas dans la glycolyse]
Erreur fréquente au Probatoire : écrire que la glycolyse « consomme de l'oxygène » ou qu'elle a lieu dans la mitochondrie. C'est faux. La glycolyse est une voie \textbf{anaérobie} : elle fonctionne aussi bien en présence qu'en absence d'\ce{O2}. C'est d'ailleurs pour cela qu'elle est commune à la respiration et à la fermentation. L'oxygène n'intervient directement qu'à la \textbf{toute dernière étape}, la chaîne respiratoire.
\end{attention}

\courssub{Étape de liaison : la décarboxylation oxydative du pyruvate}

Avant d'entrer dans le cycle de Krebs, le pyruvate produit dans le cytoplasme doit franchir la double membrane mitochondriale (via un transporteur spécifique) pour atteindre la matrice. Il y subit une réaction de décarboxylation oxydative catalysée par le complexe \textbf{pyruvate déshydrogénase}.

\begin{propriete}[Bilan de la décarboxylation oxydative (pour 1 glucose = 2 pyruvate)]
\begin{itemize}
\item \textbf{2 CO$_2$} libérés (premier dégagement de CO$_2$ de la respiration) ;
\item \textbf{2 NADH} formés (réduction de 2 \ce{NAD+}) ;
\item \textbf{2 Acétyl-CoA} formés (groupement acétyle 2C porté par la coenzyme A).
\end{itemize}
\end{propriete}

\begin{attention}[Ne pas confondre pyruvate et acide pyruvique]
Au pH physiologique (environ 7,4), l'acide pyruvique est quasi entièrement sous sa forme base conjuguée : le \textbf{pyruvate} (ion \ce{CH3COCOO-}). Dans le rédaction scientifique rigoureuse (et au Probatoire), on utilise le terme \textbf{pyruvate}.
\end{attention}

\courssub{Deuxième étape : le cycle de Krebs}

L'acétyl-CoA produit dans la matrice entre dans une série de réactions \textbf{cycliques} : le cycle de Krebs (ou cycle de l'acide citrique, ou cycle des acides tricarboxyliques).

\begin{definition}[Cycle de Krebs]
Le \cle{cycle de Krebs} est une suite de huit réactions enzymatiques cycliques qui se déroule dans la \textbf{matrice mitochondriale}. Il achève l'oxydation de l'acétyl-CoA : les deux atomes de carbone de l'acétyl-CoA sont entièrement libérés sous forme de \ce{CO2}, tandis que l'énergie est récupérée sous forme de molécules réduites (\textbf{NADH} et \textbf{FADH$_2$}) et d'un \textbf{ATP} (ou GTP) par tour.
\end{definition}

Comme une molécule de glucose donne 2 acétyl-CoA, le cycle tourne \textbf{deux fois} par glucose.

\textbf{Bilan du cycle de Krebs pour une molécule de glucose (2 tours) :}
\begin{itemize}
\item \textbf{4 CO$_2$} libérés (s'ajoutant aux 2 CO$_2$ de la décarboxylation oxydative = 6 CO$_2$ total, soit tout le carbone du glucose) ;
\item \textbf{2 ATP} (ou GTP) produits directement (phosphorylation au niveau du substrat) ;
\item \textbf{6 NADH} et \textbf{2 FADH$_2$} formés, qui constituent le véritable « trésor énergétique » de cette étape pour la chaîne respiratoire.
\end{itemize}

\begin{savaistu}[Hans Krebs, un prix Nobel pour un cycle]
C'est le biochimiste britannique d'origine allemande Hans Krebs qui a élucidé ce cycle en 1937, travail couronné par le prix Nobel de physiologie ou médecine en 1953. Le cycle porte son nom. Il est parfois appelé « cycle de l'acide citrique » car le citrate (acide citrique, celui des agrumes !) est le premier composé formé dans le cycle par condensation de l'acétyl-CoA sur l'oxaloacétate.
\end{savaistu}

\courssub{Troisième étape : la chaîne respiratoire et la phosphorylation oxydative}

Voici l'étape décisive, celle qui produit l'essentiel de l'énergie et celle où l'oxygène entre enfin en scène. Elle se déroule sur la \textbf{membrane interne} de la mitochondrie (au niveau des crêtes).

\begin{definition}[Chaîne respiratoire et phosphorylation oxydative]
La \cle{chaîne respiratoire} est un ensemble de \textbf{complexes protéiques transporteurs d'électrons} (Complexes I à IV) fixés sur la membrane interne de la mitochondrie. Les molécules réduites \ce{NADH} et \ce{FADH2} y cèdent leurs électrons, qui circulent de complexe en complexe en libérant de l'énergie. Cette énergie est utilisée pour pomper des protons (\ce{H+}) de la matrice vers l'espace intermembranaire, créant un \textbf{gradient électrochimique de protons} (force protonmotrice). Le retour des protons dans la matrice à travers l'\textbf{ATP-synthase} (Complexe V) alimente la synthèse massive d'ATP : c'est la \textbf{phosphorylation oxydative}. L'\textbf{oxygène} est l'accepteur final d'électrons : il capte les électrons en bout de chaîne (au Complexe IV, cytochrome oxydase) et, avec des protons, forme de \textbf{l'eau}.
\end{definition}

\begin{experience}[Démarche d'investigation : couplage consommation O$_2$ / synthèse ATP]
\textbf{Observation.} Sur des mitochondries isolées, on mesure la consommation d'\ce{O2} (électrode à O$_2$) et la production d'ATP en présence de substrats (pyruvate + malate). \textbf{Expérience.} On ajoute un \textbf{inhibiteur de l'ATP-synthase} (oligomycine) : la production d'ATP s'arrête \emph{et} la consommation d'\ce{O2} chute brutalement. On ajoute ensuite un \textbf{découpleur} (DNP, qui rend la membrane perméable aux H$^+$) : la consommation d'\ce{O2} repart à la hausse (maximale) mais l'ATP n'est pas synthétisé. \textbf{Interprétation.} Le flux d'électrons (consommation O$_2$) et la synthèse d'ATP sont couplés par le gradient de protons. Bloquer le retour des H$^+$ (oligomycine) arrête la pompe à protons et donc le flux d'électrons. Détruire le gradient (DNP) libère le flux d'électrons mais empêche la synthèse ATP. \textbf{Conclusion.} L'énergie libérée par la chaîne respiratoire est stockée temporairement sous forme de gradient de protons ; ce gradient est la source d'énergie directe de l'ATP-synthase (chimiosmose, théorie de Peter Mitchell, Nobel 1978).
\end{experience}

\begin{popfigure}
\begin{center}
\begin{tikzpicture}[font=\small, >=Stealth]
% Membrane interne (double trait pour suggérer bicouche)
\draw[fill=popblue!15, thick] (0,0) rectangle (14,1.2);
\draw[thick] (0,0.6) -- (14,0.6); % Milieu membrane
\node[popblue, font=\footnotesize] at (7,1.5) {Espace intermembranaire (fort [H$^+$], pH bas)};
\node[popblue, font=\footnotesize] at (7,-0.8) {Matrice mitochondriale (faible [H$^+$], pH haut)};

% Complexes
\node[draw, fill=poporangeL, rounded corners=3pt, minimum width=1.3cm, minimum height=0.9cm] (CI) at (1.5,0.6) {Complexe I};
\node[draw, fill=poporangeL, rounded corners=3pt, minimum width=1.3cm, minimum height=0.9cm] (CII) at (3.5,0.6) {Complexe II};
\node[draw, fill=poporangeL, rounded corners=3pt, minimum width=1.3cm, minimum height=0.9cm] (CIII) at (6.0,0.6) {Complexe III};
\node[draw, fill=poporangeL, rounded corners=3pt, minimum width=1.3cm, minimum height=0.9cm] (CIV) at (8.5,0.6) {Complexe IV};
\node[draw, fill=popgoldL, rounded corners=3pt, minimum width=1.6cm, minimum height=0.9cm] (CV) at (11.5,0.6) {ATP-synthase};

% NADH -> CI
\node[draw, fill=popgreenL, rounded corners=3pt, minimum width=1.2cm] (NADH) at (1.5,2.2) {NADH};
\draw[->, thick, popgreen] (NADH.south) -- (CI.north) node[midway, right, font=\footnotesize] {2 e$^-$};
\draw[->, thick, popgreen, dashed] (1.5,0.6) -- (1.5,-0.5) node[midway, right, font=\footnotesize] {4 H$^+$ pompés};

% FADH2 -> CII
\node[draw, fill=popgreenL, rounded corners=3pt, minimum width=1.2cm] (FADH2) at (3.5,2.2) {FADH$_2$};
\draw[->, thick, popgreen] (FADH2.south) -- (CII.north) node[midway, right, font=\footnotesize] {2 e$^-$};

% Flux électrons
\draw[->, very thick, popink] (CI.east) -- (CIII.west);
\draw[->, very thick, popink] (CII.east) -- (3.8,0.6) -- (3.8,0.2) -- (5.7,0.2) -- (5.7,0.6) -- (CIII.west); % Via Q (ubiquinone)
\node[popmuted, font=\footnotesize] at (4.7,0.1) {Ubiquinone (Q)};
\draw[->, very thick, popink] (CIII.east) -- (CIV.west);
\node[popmuted, font=\footnotesize] at (7.25,1.0) {Cytochrome c};

% Pompe à protons CIII, CIV
\draw[->, thick, poporange, dashed] (CIII) -- ++(0,-0.5) node[midway, right, font=\footnotesize] {4 H$^+$};
\draw[->, thick, poporange, dashed] (CIV) -- ++(0,-0.5) node[midway, right, font=\footnotesize] {2 H$^+$};

% O2 -> H2O
\node[draw, fill=poppinkL, rounded corners=3pt, minimum width=1.2cm] (O2) at (8.5,2.2) {\ce{O2} (1/2 par e$^-$)};
\draw[->, thick, poppink] (O2.south) -- (CIV.north) node[midway, right, font=\footnotesize] {4 e$^-$};
\node[draw, fill=popblueL, rounded corners=3pt, minimum width=1.2cm] (H2O) at (8.5,-1.6) {2 H$_2$O};
\draw[->, thick, poppink] (CIV.south) -- (H2O.north) node[midway, right, font=\footnotesize] {4 H$^+$};

% ATP Synthase
\draw[->, very thick, popgold] (11.5,0) -- (11.5,-0.8) node[midway, right, font=\footnotesize] {H$^+$ flux (chimiosmose)};
\node[draw, fill=popgoldL, rounded corners=3pt, minimum width=1.5cm] (ATP) at (11.5,-1.8) {ADP + Pi $\rightarrow$ ATP};
\draw[->, thick, popgold] (11.5,-0.8) -- (ATP.north);

% Légendes flux
\node[popink, font=\footnotesize, align=center] at (5,1.8) {Flux d'électrons \\ (libération énergie)};
\node[poporange, font=\footnotesize, align=center] at (10, -0.5) {Gradient $\Delta \mu_{H^+}$ \\ (force protonmotrice)};
\end{tikzpicture}
\end{center}
\legende{Schéma simplifié de la chaîne respiratoire et de la phosphorylation oxydative (chimiosmose). Les électrons du NADH (Complexe I) et du FADH$_2$ (Complexe II) traversent les complexes I--III--IV, pompant des H$^+$ vers l'espace intermembranaire. L'O$_2$ est l'accepteur final (Complexe IV). Le retour des H$^+$ par l'ATP-synthase (Complexe V) synthétise l'ATP.}
\end{popfigure}

\begin{savaistu}[Le cyanure, un poison de la chaîne respiratoire]
Le cyanure (ion \ce{CN-}) est un poison mortel qui agit en quelques minutes : il se fixe sur le fer du \textbf{Complexe IV (cytochrome oxydase)} et bloque le transfert des électrons à l'oxygène. Les électrons ne peuvent plus atteindre l'\ce{O2} : la chaîne s'arrête, le gradient de protons s'effondre, la production d'ATP s'arrête, et les cellules — surtout celles du cerveau et du cœur, très gourmandes en énergie — meurent asphyxiées \emph{alors même que le sang est plein d'oxygène}. C'est une preuve dramatique du rôle indispensable de la chaîne respiratoire et du couplage chimiosmotique.
\end{savaistu}

\courssub{Bilan énergétique global et équation de la respiration}

\begin{propriete}[Rendement énergétique de la respiration cellulaire (valeur admise au programme)]
L'oxydation complète d'une molécule de glucose par la respiration cellulaire produit au total \textbf{environ 36 à 38 molécules d'ATP} (valeur théorique classique admise au Probatoire). Ce rendement se répartit ainsi :
\begin{itemize}
\item \textbf{2 ATP} produits par la glycolyse (phosphorylation au niveau du substrat) ;
\item \textbf{2 ATP} (ou GTP) produits par le cycle de Krebs (phosphorylation au niveau du substrat) ;
\item \textbf{environ 32 à 34 ATP} produits par la phosphorylation oxydative (chaîne respiratoire + ATP-synthase) grâce à l'oxydation des 10 NADH et 2 FADH$_2$ totaux.
\end{itemize}
\end{propriete}

\begin{exemplebox}[Répartition chiffrée des ATP (rappel)]
Pour une molécule de glucose totalement oxydée :
\[ \underbrace{2\ \text{ATP}}_{\text{glycolyse}} + \underbrace{2\ \text{ATP}}_{\text{cycle de Krebs}} + \underbrace{32 \text{ à } 34\ \text{ATP}}_{\text{phosphorylation oxydative}} = \textbf{environ 36 à 38 ATP} \]
On constate que c'est la \textbf{chaîne respiratoire} qui fournit l'essentiel de l'énergie (environ 90 \% du rendement) : sans oxygène, la cellule perd cette quasi-totalité.
\end{exemplebox}

\begin{methode}[Écrire et équilibrer le bilan simplifié de la respiration]
Pour rédiger correctement l'équation-bilan au Probatoire, procède en quatre étapes :
\begin{enumerate}
\item \textbf{Identifie les réactifs} : le glucose \ce{C6H12O6} et le dioxygène \ce{O2} (à gauche de la flèche).
\item \textbf{Identifie les produits} : le dioxyde de carbone \ce{CO2}, l'eau \ce{H2O} et l'énergie (ATP) (à droite).
\item \textbf{Équilibre les atomes de carbone} : 6 C à gauche, donc 6 \ce{CO2} à droite.
\item \textbf{Équilibre H puis O} : 12 H à gauche, donc 6 \ce{H2O} ; on compte alors les O à droite ($6\times2 + 6\times1 = 18$), dont 6 fournis par le glucose, il en manque 12, soit 6 \ce{O2}.
\end{enumerate}
D'où l'équation-bilan simplifiée :
\[ \ce{C6H12O6 + 6O2 -> 6CO2 + 6H2O} + \text{énergie (ATP)} \]
\end{methode}

\begin{attention}[Pièges de rédaction fréquents au Probatoire]
\begin{itemize}
\item Ne jamais écrire que la respiration « produit de l'oxygène » : c'est la \textbf{photosynthèse} qui le produit ; la respiration le \textbf{consomme}.
\item Le \ce{CO2} rejeté provient du \textbf{carbone du glucose} (décarboxylation oxydative + cycle de Krebs), pas de la transformation de l'\ce{O2} : l'oxygène inspiré se retrouve exclusivement dans \textbf{l'eau} formée au Complexe IV.
\item L'équation-bilan est un \textbf{bilan global} : elle masque les quatre étapes (glycolyse, liaison, Krebs, chaîne). Dire « le glucose réagit directement avec l'\ce{O2} » est faux : l'\ce{O2} n'intervient qu'au Complexe IV.
\item Attention à l'orthographe : \textbf{glycolyse} (pas « glycolise »), \textbf{Krebs} (avec majuscule), \textbf{phosphorylation oxydative}.
\end{itemize}
\end{attention}

\begin{exoresolu}[Exploiter un bilan énergétique et gazeux]
\textbf{Énoncé.} Une cellule musculaire d'un sportif de Yaoundé dégrade 10 molécules de glucose en présence d'oxygène abondant. (a) Combien de molécules d'ATP sont produites au maximum ? (b) Quel volume de réaction en \ce{CO2} cela représente-t-il en molécules de \ce{CO2} ? (c) Que se passe-t-il si l'apport en \ce{O2} est interrompu brutalement ?
\tcblower
\textbf{Solution.} 
(a) Le rendement théorique admis est de 36 à 38 ATP par glucose. Pour 10 molécules : $10 \times 36 = 360$ à $10 \times 38 = 380$ molécules d'ATP.
(b) D'après l'équation-bilan, 1 glucose donne 6 \ce{CO2}, donc 10 glucose donnent $10 \times 6 = 60$ molécules de \ce{CO2}.
(c) Sans \ce{O2}, le Complexe IV n'a plus d'accepteur final d'électrons. La chaîne respiratoire s'arrête immédiatement : plus de pompage de protons, plus de gradient, plus de phosphorylation oxydative. Les NADH et FADH$_2$ s'accumulent sous forme réduite, ce qui inhibe les déshydrogénases du cycle de Krebs (rétroinhibition) : le cycle de Krebs s'arrête. Seule la glycolyse continue (voie anaérobie), avec un rendement de $10 \times 2 = 20$ ATP : la cellule perd environ 95 \% de sa production énergétique. Pour régénérer le \ce{NAD+} indispensable à la glycolyse, la cellule doit recourir à la \textbf{fermentation lactique}, voie étudiée dans la section suivante.
\end{exoresolu}

\begin{aretenir}[L'essentiel de la section]
\begin{itemize}
\item Quatre phases successives : \textbf{glycolyse} (cytoplasme, anaérobie, 2 ATP + 2 NADH + 2 pyruvate), \textbf{décarboxylation oxydative du pyruvate} (matrice, 2 NADH + 2 CO$_2$ + 2 Acétyl-CoA), \textbf{cycle de Krebs} (matrice, 2 ATP + 6 NADH + 2 FADH$_2$ + 4 CO$_2$), \textbf{chaîne respiratoire + phosphorylation oxydative} (membrane interne, \ce{O2} accepteur final $\rightarrow$ H$_2$O, $\approx$ 32--34 ATP).
\item Bilan global : \textbf{environ 36 à 38 ATP} par molécule de glucose.
\item Équation-bilan : \ce{C6H12O6 + 6O2 -> 6CO2 + 6H2O} + énergie (ATP).
\item L'oxygène est indispensable au rendement maximal : sans lui, chaîne respiratoire et cycle de Krebs s'arrêtent, seul subsiste le faible rendement de la glycolyse (fermentation).
\end{itemize}
\end{aretenir}

\coursec{La fermentation, voie alternative sans oxygène}

Dans la section précédente, nous avons suivi le glucose pas à pas à travers les trois étapes de la respiration cellulaire : la glycolyse dans le cytoplasme, le cycle de Krebs dans la matrice mitochondriale, puis la chaîne respiratoire sur les crêtes mitochondriales, où l'oxygène joue le rôle d'accepteur final d'électrons. Ce long cheminement est très rentable : il libère environ 36 à 38 molécules d'ATP par molécule de glucose. Mais tout ce dispositif repose sur une condition absolue : la présence d'\cle{oxygène}. Or, dans certaines situations de la vie courante — un sprint pour attraper le car au carrefour, un match de football disputé sous le soleil de Douala — les muscles consomment l'oxygène plus vite que le sang ne peut le leur apporter. Que devient alors la production d'énergie ? La cellule est-elle condamnée à s'arrêter ? C'est toute la question de cette section.

\courssub{Situation problème : les crampes du sprinteur}

\textbf{Observation.} Un élève de Première court un 100 mètres à fond. À l'arrivée, il est essoufflé, ses cuisses « brûlent », et quelques heures plus tard il ressent des douleurs musculaires (courbatures). Pourtant, pendant l'effort, ses muscles ont continué à fonctionner et à produire de l'énergie.

\textbf{Hypothèse.} Faute d'oxygène suffisant, la cellule musculaire utiliserait un autre mode de dégradation du glucose, moins rentable et produisant une substance responsable des douleurs.

\textbf{Données expérimentales.} On mesure la concentration sanguine en acide lactique chez un athlète avant et après un sprint de 400 m :

\begin{center}
\begin{tabularx}{\linewidth}{|l|X|X|}\hline
\rowcolor{popblueL} \textbf{Moment du prélèvement} & \textbf{Acide lactique sanguin (mmol/L)} & \textbf{Consommation d'O$_2$} \\ \hline
Au repos & environ 1 & normale \\ \hline
Juste après le sprint & 12 à 15 & fortement augmentée (essoufflement) \\ \hline
30 min après l'effort & retour progressif vers 1 & reste élevée puis normale \\ \hline
\end{tabularx}
\end{center}

\textbf{Interprétation.} Pendant l'effort intense, l'apport en O$_2$ est insuffisant : les cellules musculaires produisent de l'acide lactique, qui s'accumule dans le sang. Après l'effort, l'oxygène respiré en excès (essoufflement prolongé) permet d'éliminer cet acide lactique.

\textbf{Conclusion.} Il existe bien une voie de production d'énergie \emph{sans oxygène} : c'est la \cle{fermentation}.

\courssub{Définition et caractéristiques générales de la fermentation}

\begin{definition}[La fermentation]
La \cle{fermentation} est un mode de dégradation \textbf{incomplète} du glucose qui se déroule en \textbf{absence d'oxygène} (on dit en anaérobiose), \textbf{uniquement dans le cytoplasme} (hyaloplasme), sans intervention de la mitochondrie. Elle ne fait intervenir ni le cycle de Krebs, ni la chaîne respiratoire. On distingue deux types principaux :
\begin{itemize}
\item la \cle{fermentation lactique}, qui produit de l'\textbf{acide lactique} (cellules musculaires humaines en effort intense, certaines bactéries comme les ferments du yaourt) ;
\item la \cle{fermentation alcoolique}, qui produit de l'\textbf{éthanol} et du \textbf{dioxyde de carbone} (levures).
\end{itemize}
\end{definition}

\begin{attention}[La fermentation n'est pas une « respiration sans oxygène »]
Erreur classique au Probatoire : écrire que la fermentation est « la respiration en anaérobiose ». C'est faux. La respiration cellulaire exige par définition l'oxygène, la mitochondrie et la chaîne respiratoire. La fermentation se passe \textbf{entièrement dans le cytoplasme}, ne met en jeu \textbf{aucune mitochondrie} et \textbf{aucune chaîne respiratoire}. Ce sont deux processus distincts qui ne partagent qu'une seule étape : la glycolyse.
\end{attention}

\courssub{Le point de départ commun : la glycolyse}

Que l'oxygène soit présent ou non, la dégradation du glucose commence toujours par la \cle{glycolyse}, dans le cytoplasme : une molécule de glucose (6 atomes de carbone) est scindée en deux molécules d'acide pyruvique (3 atomes de carbone), avec un gain net de \textbf{2 ATP} et la production de molécules de NADH (forme réduite du transporteur NAD$^+$).

C'est à la sortie de la glycolyse que tout se joue :

\begin{itemize}
\item \textbf{Si l'oxygène est présent} : l'acide pyruvique entre dans la mitochondrie, est dégradé complètement par le cycle de Krebs et la chaîne respiratoire, jusqu'à CO$_2$ et H$_2$O. Bilan total : environ 36 à 38 ATP.
\item \textbf{Si l'oxygène est absent} : l'acide pyruvique reste dans le cytoplasme et subit une transformation qui ne rapporte \textbf{aucun ATP supplémentaire}. Le bilan de la fermentation se limite donc aux \textbf{2 ATP de la glycolyse}.
\end{itemize}

\begin{popfigure}
\begin{center}
\begin{tikzpicture}[font=\small, >=Stealth]
% Glucose
\node[draw=popdark, fill=popgoldL, rounded corners, minimum width=2.4cm, minimum height=0.9cm] (glu) at (0,0) {\textbf{Glucose}};
\node[draw=popblue, fill=popblueL, rounded corners, minimum width=3.2cm, minimum height=1cm, align=center] (gly) at (4.6,0) {\textbf{Glycolyse}\\ (cytoplasme) : 2 ATP};
\node[draw=popdark, fill=popgoldL, rounded corners, minimum width=2.8cm, minimum height=0.9cm] (pyr) at (9.2,0) {\textbf{Acide pyruvique}};
\draw[->, thick, popdark] (glu) -- (gly);
\draw[->, thick, popdark] (gly) -- (pyr);
% Branche avec O2
\node[draw=popgreen, fill=popgreenL, rounded corners, minimum width=4.6cm, minimum height=1.3cm, align=center] (resp) at (9.2,2.6) {\textbf{RESPIRATION} (mitochondrie)\\ CO$_2$ + H$_2$O\\ \textbf{+ 34 à 36 ATP}};
\draw[->, very thick, popgreen] (pyr.north) -- (resp.south) node[midway, right, popgreen!60!black]{avec O$_2$};
% Branche sans O2
\node[draw=poporange, fill=poporangeL, rounded corners, minimum width=4.2cm, minimum height=1.1cm, align=center] (lact) at (6.4,-2.6) {\textbf{Fermentation lactique}\\ acide lactique (0 ATP de plus)};
\node[draw=poppurple, fill=poppurpleL, rounded corners, minimum width=4.2cm, minimum height=1.1cm, align=center] (alc) at (12.4,-2.6) {\textbf{Fermentation alcoolique}\\ éthanol + CO$_2$ (0 ATP de plus)};
\draw[->, very thick, poporange] (pyr.south) -- (lact.north) node[midway, left, poporange!70!black]{sans O$_2$};
\draw[->, very thick, poppurple] (pyr.south) -- (alc.north);
% Bilans
\node[draw=popgreen, fill=popgreenL!50, rounded corners, align=center] at (14.6,2.6) {\textbf{Bilan : 36--38 ATP}\\ par glucose};
\node[draw=popdark, fill=popcream, rounded corners, align=center] at (9.4,-4.3) {\textbf{Bilan fermentation : 2 ATP} par glucose (glycolyse seule)};
\end{tikzpicture}
\end{center}
\legende{Les deux destins de l'acide pyruvique après la glycolyse. Avec oxygène, il entre dans la mitochondrie (respiration, 36 à 38 ATP au total). Sans oxygène, il est transformé dans le cytoplasme en acide lactique (fermentation lactique) ou en éthanol et CO$_2$ (fermentation alcoolique), sans gain d'ATP supplémentaire : seules les 2 ATP de la glycolyse sont produites.}
\end{popfigure}

\courssub{La fermentation lactique}

\begin{definition}[Fermentation lactique]
La \cle{fermentation lactique} est la transformation de l'acide pyruvique en \textbf{acide lactique}, dans le cytoplasme, en absence d'oxygène. Elle a lieu chez l'Homme dans les \textbf{cellules musculaires} lors d'un effort intense et bref, ainsi que chez certaines bactéries (lactobacilles, utilisés pour fabriquer le yaourt et les laits fermentés).
\end{definition}

Son équation simplifiée s'écrit :
\[ \ce{C6H12O6 -> 2 CH3CH(OH)COOH} \quad \textbf{+ 2 ATP} \]
(glucose $\rightarrow$ 2 molécules d'acide lactique).

\begin{exemplebox}[Pourquoi les muscles « brûlent » pendant un sprint]
Lors d'un 100 m, la vitesse de contraction musculaire est telle que la consommation d'ATP dépasse ce que la respiration peut fournir : le sang n'apporte pas l'oxygène assez vite. Les fibres musculaires basculent alors en fermentation lactique. L'acide lactique s'accumule dans le muscle puis diffuse dans le sang : il est responsable de la sensation de brûlure, des \textbf{crampes} et, en partie, des \textbf{courbatures} ressenties après l'effort. C'est ce mécanisme qui limite la durée d'un effort maximal à quelques dizaines de secondes.
\end{exemplebox}

\courssub{La fermentation alcoolique}

\begin{definition}[Fermentation alcoolique]
La \cle{fermentation alcoolique} est la transformation de l'acide pyruvique en \textbf{éthanol} (alcool éthylique) et en \textbf{dioxyde de carbone}, dans le cytoplasme des \textbf{levures} (champignons unicellulaires, par exemple \emph{Saccharomyces cerevisiae}), en absence d'oxygène.
\end{definition}

Son équation simplifiée s'écrit :
\[ \ce{C6H12O6 -> 2 C2H5OH + 2 CO2} \quad \textbf{+ 2 ATP} \]

\begin{experience}[Mise en évidence de la fermentation alcoolique par la levure de boulanger]
\textbf{Matériel :} deux tubes à essai A et B, eau tiède, glucose, levure de boulanger, eau de chaux, huile végétale.

\textbf{Protocole :} dans le tube A, on mélange eau tiède + glucose + levure, et l'on verse une couche d'huile en surface (pour isoler le milieu de l'air). Dans le tube B (témoin), on met le même mélange sans levure. On laisse les tubes à l'étuve à 35 °C pendant 30 minutes, puis on recueille le gaz dégagé dans de l'eau de chaux.

\textbf{Observations :} dans le tube A, on observe un dégagement gazeux abondant ; le gaz trouble l'eau de chaux (présence de CO$_2$) ; le liquide dégage une odeur d'alcool. Dans le tube B, rien ne se produit.

\textbf{Interprétation :} en absence d'oxygène (sous l'huile), la levure dégrade le glucose en éthanol et en CO$_2$ : c'est la fermentation alcoolique. Sans levure, aucune réaction : ce sont bien des cellules vivantes qui réalisent la transformation.
\end{experience}

\begin{savaistu}[Le vin de palme et le « bil-bil » : la fermentation au quotidien au Cameroun]
La fermentation alcoolique des levures est à la base de nombreuses productions traditionnelles camerounaises. Le \textbf{vin de palme}, récolté comme sève sucrée du palmier ou du raphia, fermente spontanément grâce aux levures présentes dans l'air : en quelques heures, les sucres se transforment en éthanol et CO$_2$ (d'où son pétillement). Le \textbf{« bil-bil »}, bière de mil ou de maïs brassée dans le Grand Nord, résulte de la fermentation des sucres du grain par des levures. Même principe pour la \textbf{levée de la pâte à beignets} (« puff-puff ») : le CO$_2$ dégagé par les levures forme les bulles qui rendent la pâte légère et aérée. La biotechnologie traditionnelle camerounaise est donc une application directe... de ce chapitre !
\end{savaistu}

\courssub{Le rôle profond de la fermentation : régénérer le NAD$^+$}

On peut se demander : si la fermentation ne rapporte aucun ATP supplémentaire, à quoi sert-elle ? Pourquoi la cellule transforme-t-elle l'acide pyruvique au lieu de le laisser s'accumuler ?

La réponse tient en une notion simple que nous admettrons : la glycolyse a besoin, pour fonctionner en continu, d'une molécule « navette » appelée \cle{NAD$^+$}, qui accepte temporairement des électrons (elle devient NADH). Or la cellule ne possède qu'un stock limité de NAD$^+$. En respiration, la chaîne respiratoire mitochondriale régénère le NAD$^+$ grâce à l'oxygène. \textbf{Sans oxygène, cette régénération est impossible} : tout le NAD$^+$ serait bientôt bloqué sous forme de NADH, et la glycolyse — seule source d'ATP restante — s'arrêterait.

\begin{aretenir}[Rôle essentiel de la fermentation]
La fermentation (lactique ou alcoolique) permet de \textbf{régénérer le NAD$^+$} à partir du NADH produit par la glycolyse. Ce faisant, elle garantit la \textbf{poursuite de la glycolyse} et donc la production des 2 ATP par glucose, même en l'absence totale d'oxygène. La fermentation est donc une « roue de secours » : peu rentable, mais vitale à court terme.
\end{aretenir}

\courssub{Comparaison des rendements : respiration contre fermentation}

\begin{propriete}[Rendement énergétique comparé]
Pour une molécule de glucose dégradée :
\begin{itemize}
\item la \textbf{respiration cellulaire} produit environ \textbf{36 à 38 ATP} (dégradation complète du glucose en CO$_2$ et H$_2$O) ;
\item la \textbf{fermentation} ne produit que \textbf{2 ATP} (ceux de la glycolyse ; dégradation incomplète).
\end{itemize}
La respiration est donc environ \textbf{18 fois plus rentable} que la fermentation ($36 / 2 = 18$).
\end{propriete}

\begin{popfigure}
\begin{center}
\begin{tikzpicture}
\begin{axis}[
    ybar, bar width=1.6cm,
    width=0.85\linewidth, height=6.5cm,
    ymin=0, ymax=42,
    ylabel={ATP produits par molécule de glucose},
    symbolic x coords={Respiration, Fermentation},
    xtick=data,
    nodes near coords,
    every node near coord/.append style={font=\bfseries},
    ymajorgrids=true, grid style={popline!40},
    axis line style={popdark},
    ylabel style={font=\small},
]
\addplot[fill=popgreen!70, draw=popgreen!60!black] coordinates {(Respiration,37)};
\addplot[fill=poporange!80, draw=poporange!70!black] coordinates {(Fermentation,2)};
\end{axis}
\end{tikzpicture}
\end{center}
\legende{Comparaison du nombre d'ATP produits par molécule de glucose en respiration (environ 37 ATP) et en fermentation (2 ATP). La respiration est environ 18 fois plus rentable.}
\end{popfigure}

\begin{exemplebox}[Un exemple chiffré pour bien comprendre l'écart]
Pour obtenir l'équivalent de 36 ATP, une cellule en fermentation doit dégrader :
\[ \frac{36 \text{ ATP}}{2 \text{ ATP par glucose}} = 18 \text{ molécules de glucose} \]
là où une cellule en respiration n'en dégrade qu'\textbf{une seule}. La fermentation « gaspille » donc les réserves de glucose : c'est pourquoi un effort anaérobie épuise très vite les stocks de glycogène musculaire, et pourquoi l'organisme ne peut pas fonctionner longtemps en fermentation.
\end{exemplebox}

\begin{methode}[Comparer respiration et fermentation en tableau]
Pour répondre à une question du type « comparez respiration et fermentation », construis toujours un tableau à 4 ou 5 critères :
\begin{enumerate}
\item \textbf{Lieu} dans la cellule ;
\item \textbf{Besoin en O$_2$} (présence ou absence) ;
\item \textbf{Produits} de la dégradation ;
\item \textbf{Rendement} en ATP ;
\item éventuellement : organites impliqués, caractère complet ou incomplet de la dégradation.
\end{enumerate}
Modèle attendu :
\begin{center}
\begin{tabularx}{\linewidth}{|l|X|X|}\hline
\rowcolor{popblueL} \textbf{Critère} & \textbf{Respiration} & \textbf{Fermentation} \\ \hline
Lieu & Cytoplasme + mitochondrie & Cytoplasme uniquement \\ \hline
Besoin en O$_2$ & Indispensable (accepteur final) & Aucun (anaérobiose) \\ \hline
Produits & CO$_2$ + H$_2$O & Acide lactique, ou éthanol + CO$_2$ \\ \hline
Dégradation du glucose & Complète & Incomplète \\ \hline
Rendement & 36 à 38 ATP / glucose & 2 ATP / glucose \\ \hline
\end{tabularx}
\end{center}
\end{methode}

\courssub{Conséquences physiologiques : essoufflement et dette d'oxygène}

Puisque la fermentation est 18 fois moins rentable et produit un déchet (l'acide lactique), l'organisme cherche toujours à \textbf{rester en respiration} le plus longtemps possible. Lors d'un effort prolongé (course de fond, travail aux champs), le corps améliore l'apport en O$_2$ aux muscles :

\begin{itemize}
\item la \textbf{fréquence respiratoire} augmente : c'est l'\cle{essoufflement} ;
\item la \textbf{fréquence cardiaque} augmente : le sang (transporteur d'O$_2$) circule plus vite ;
\item les vaisseaux sanguins des muscles se \textbf{dilatent}.
\end{itemize}

\begin{aretenir}[La dette d'oxygène]
Après un effort intense, on continue à respirer fortement pendant plusieurs minutes, alors que l'effort est terminé. Cet excès de consommation d'oxygène sert à \textbf{éliminer l'acide lactique} accumulé : une partie est oxydée en CO$_2$ et H$_2$O dans le muscle, une autre est reconvertie en glucose dans le foie (cycle de Cori). On parle de \cle{dette d'oxygène} : l'organisme « rembourse » l'oxygène qui a manqué pendant l'effort. La fermentation lactique est donc \textbf{réversible} : son produit disparaît dès que l'oxygène redevient disponible.
\end{aretenir}

\begin{exoresolu}[Exploitation de données expérimentales]
\textbf{Énoncé.} On mesure chez deux sujets A et B effectuant le même exercice physique la concentration sanguine d'acide lactique (en mmol/L) :

\begin{center}
\begin{tabular}{|l|c|c|c|c|}\hline
Durée de l'effort (min) & 0 & 5 & 10 & 15 \\ \hline
Sujet A (entraîné) & 1 & 2 & 2,5 & 3 \\ \hline
Sujet B (sédentaire) & 1 & 5 & 9 & 14 \\ \hline
\end{tabular}
\end{center}

1. Compare l'évolution de l'acide lactique chez les deux sujets.
2. Explique la différence observée.
3. Précise le devenir de l'acide lactique après l'arrêt de l'effort.

\tcblower
\textbf{Solution détaillée.}

\textbf{1. Comparaison.} Chez les deux sujets, la concentration d'acide lactique augmente au cours de l'effort, mais cette augmentation est beaucoup plus forte et plus rapide chez le sujet B : après 15 minutes, elle atteint 14 mmol/L chez B contre seulement 3 mmol/L chez A (environ 4,7 fois moins).

\textbf{2. Explication.} L'acide lactique est le produit de la fermentation lactique, qui se déclenche lorsque l'apport en O$_2$ aux muscles devient insuffisant. Chez le sujet A, entraîné, le cœur et l'appareil respiratoire sont plus performants, et les muscles possèdent davantage de mitochondries : l'oxygénation reste suffisante plus longtemps, les cellules restent en respiration (36 à 38 ATP par glucose) et produisent peu d'acide lactique. Chez le sujet B, sédentaire, l'apport en O$_2$ est vite dépassé par les besoins : les muscles basculent précocement en fermentation lactique (2 ATP seulement par glucose), d'où l'accumulation rapide d'acide lactique, source de fatigue et de crampes.

\textbf{3. Devenir de l'acide lactique.} Après l'effort, l'essoufflement persistant (dette d'oxygène) fournit l'O$_2$ nécessaire pour oxyder une partie de l'acide lactique en CO$_2$ et H$_2$O dans les muscles, et pour reconvertir l'autre partie en glucose dans le foie. La concentration sanguine revient progressivement à sa valeur de repos (environ 1 mmol/L).
\end{exoresolu}

\begin{attention}[Pièges fréquents à éviter]
\begin{itemize}
\item Ne dis jamais que la fermentation « produit de l'énergie sans glucose » : elle dégrade toujours du glucose, mais de façon incomplète.
\item Ne confonds pas les produits : fermentation \textbf{lactique} $\rightarrow$ acide lactique seul ; fermentation \textbf{alcoolique} $\rightarrow$ éthanol \textbf{+} CO$_2$.
\item Les 2 ATP de la fermentation sont ceux de la \textbf{glycolyse} : l'étape de fermentation proprement dite n'en rapporte aucun.
\item L'acide lactique n'est pas un poison définitif : il est éliminé après l'effort grâce à l'oxygène (dette d'oxygène).
\end{itemize}
\end{attention}

En résumé, la fermentation est une voie de secours remarquable : confinée au cytoplasme, indépendante de l'oxygène et des mitochondries, elle permet à la cellule de survivre à court terme en régénérant le NAD$^+$ indispensable à la glycolyse. Mais son faible rendement (2 ATP contre 36 à 38) et ses déchets (acide lactique) imposent à l'organisme des stratégies d'optimisation de l'oxygénation : c'est précisément l'objet de la section suivante, consacrée aux facteurs qui influencent le rendement énergétique de la cellule.

\coursec{Facteurs influençant le rendement énergétique}

Nous venons de voir que la fermentation lactique est une voie de secours, peu rentable (2 ATP par glucose contre environ 36 en respiration complète) et limitée dans le temps par l'accumulation d'acide lactique. Une question se pose alors naturellement, et elle passionne aussi bien le sportif du lycée qui court le \num{1500} mètres en EPS que le médecin du Centre Hospitalier d'Essos : \emph{comment faire pour qu'une cellule produise le plus d'ATP possible, le plus longtemps possible, sans basculer dans la fermentation ?} Autrement dit, quels sont les facteurs qui influencent le rendement énergétique de l'organisme ? C'est l'objet de cette section.

\courssub{L'oxygénation : le premier facteur limitant}

\textbf{Observation de la vie courante.} Lors d'une course de fond sur le terrain de sport du lycée, tout le monde a remarqué que les élèves respirent de plus en plus vite et de plus en plus fort : c'est l'\cle{essoufflement}. Pourquoi ce phénomène ?

Lors de l'effort musculaire, les fibres musculaires consomment beaucoup plus d'ATP qu'au repos. Pour produire cet ATP par la respiration cellulaire, elles ont besoin de davantage de dioxygène, accepteur final des électrons de la chaîne respiratoire. Or l'apport d'\ce{O2} aux cellules dépend de deux maillons en chaîne :

\begin{itemize}
\item la \cle{ventilation pulmonaire}, qui renouvelle l'air des alvéoles (fréquence et amplitude respiratoires) ;
\item la \cle{circulation sanguine}, qui transporte l'\ce{O2} fixé sur l'hémoglobine des globules rouges jusqu'aux muscles (débit cardiaque).
\end{itemize}

L'essoufflement traduit donc l'augmentation des besoins en \ce{O2} : le système nerveux accélère la ventilation et le cœur bat plus vite pour augmenter l'apport d'oxygène aux cellules musculaires. Tant que cet apport couvre les besoins, la cellule reste en \cle{respiration aérobie}. Mais lorsque l'intensité de l'effort dépasse la capacité d'approvisionnement en \ce{O2}, les mitochondries ne peuvent plus oxyder tout le pyruvate produit par la glycolyse : la cellule bascule partiellement en \cle{fermentation lactique}.

\begin{definition}[Seuil anaérobie]
Le \cle{seuil anaérobie} est l'intensité d'effort à partir de laquelle l'apport d'\ce{O2} aux cellules musculaires devient insuffisant : la fermentation lactique s'installe alors de façon significative et l'acide lactique (lactate) commence à s'accumuler dans le sang (concentration supérieure à environ \SI{4}{mmol/L}).
\end{definition}

\courssub{L'entraînement physique : plus de mitochondries, un seuil reculé}

\textbf{Situation concrète.} Deux élèves de Première C courent ensemble à la même vitesse. L'un, membre du club d'athlétisme, respire calmement ; l'autre, sédentaire, est déjà essoufflé et ressent des brûlures dans les cuisses. Même effort, même vitesse : pourquoi cette différence ?

\textbf{Hypothèse.} L'entraînement régulier modifie les cellules musculaires elles-mêmes, en augmentant leur capacité à utiliser l'\ce{O2}.

\textbf{Données expérimentales.} Des biopsies musculaires (prélèvements microscopiques de muscle) réalisées chez des sportifs avant et après plusieurs mois d'entraînement d'endurance montrent une augmentation nette du \emph{nombre} et de la \emph{taille} des mitochondries dans les fibres musculaires, ainsi qu'une augmentation du réseau de capillaires sanguins qui les irriguent.

\textbf{Interprétation.} Plus de mitochondries signifie plus de chaînes respiratoires fonctionnant en parallèle : la cellule peut oxyder davantage de pyruvate par unité de temps, donc produire plus d'ATP en aérobie. Le recours à la fermentation lactique n'est nécessaire que pour des efforts plus intenses.

\begin{propriete}[Effets de l'entraînement sur le rendement énergétique]
L'entraînement physique régulier (endurance) :
\begin{enumerate}
\item augmente le \textbf{nombre} et la \textbf{taille} des mitochondries dans les cellules musculaires ;
\item améliore la capacité d'utilisation de l'\ce{O2}, mesurée par la \cle{VO2 max} (volume maximal d'\ce{O2} consommé par minute et par kilogramme de masse corporelle) ;
\item \textbf{déplace le seuil anaérobie} vers des efforts plus intenses : le sportif entraîné reste plus longtemps en respiration aérobie et retarde l'apparition de la fermentation lactique.
\end{enumerate}
\end{propriete}

\begin{popfigure}
\begin{center}
\begin{tikzpicture}
\begin{axis}[
width=0.92\linewidth, height=6.5cm,
xlabel={Intensité de l'effort (\% de l'effort maximal)},
ylabel={Consommation d'\ce{O2} (mL/min/kg)},
xmin=0, xmax=100, ymin=0, ymax=85,
xtick={0,20,40,60,80,100},
ytick={0,20,40,60,80},
legend style={at={(0.03,0.97)}, anchor=north west, font=\small},
grid=both, grid style={popline!40},
axis line style={popdark},
]
\addplot[popcurve, domain=0:100, samples=100, color=popblue, very thick] {5 + 0.75*x};
\addlegendentry{Athlète entraîné (VO2 max $\approx 80$)}
\addplot[popcurve, domain=0:100, samples=100, color=poporange, very thick] {5 + 0.35*x};
\addlegendentry{Sujet sédentaire (VO2 max $\approx 40$)}
\draw[dashed, popblue] (axis cs:100,80) -- (axis cs:0,80);
\draw[dashed, poporange] (axis cs:100,40) -- (axis cs:0,40);
\end{axis}
\end{tikzpicture}
\end{center}
\legende{Consommation d'\ce{O2} en fonction de l'intensité de l'effort chez un athlète entraîné et un sujet sédentaire. À effort égal, les deux sujets consomment de l'\ce{O2}, mais le plafond (VO2 max) de l'athlète est environ deux fois plus élevé : il peut soutenir des efforts bien plus intenses en restant aérobie.}
\end{popfigure}

\begin{popfigure}
\begin{center}
\begin{tikzpicture}
\begin{axis}[
width=0.92\linewidth, height=6.5cm,
xlabel={Intensité de l'effort (\% de l'effort maximal)},
ylabel={Lactate sanguin (mmol/L)},
xmin=0, xmax=100, ymin=0, ymax=14,
xtick={0,20,40,60,80,100},
legend style={at={(0.03,0.97)}, anchor=north west, font=\small},
grid=both, grid style={popline!40},
axis line style={popdark},
]
\addplot[popcurve, domain=0:100, samples=100, color=popblue, very thick] {1 + 12*(x/100)^5};
\addlegendentry{Athlète entraîné}
\addplot[popcurve, domain=0:100, samples=100, color=poporange, very thick] {1 + 12*(x/100)^2.6};
\addlegendentry{Sujet sédentaire}
\addplot[domain=0:100, samples=2, dashed, popdark] {4};
\draw[dashed, popblue] (axis cs:82,0) -- (axis cs:82,4);
\draw[dashed, poporange] (axis cs:57,0) -- (axis cs:57,4);
\node[font=\small, popdark, anchor=south west] at (axis cs:2,4.3) {seuil : 4 mmol/L};
\node[font=\small, popblue, anchor=west] at (axis cs:83,1.2) {seuil anaérobie $\approx 82\,\%$};
\node[font=\small, poporange, anchor=east] at (axis cs:55,1.2) {seuil $\approx 57\,\%$};
\end{axis}
\end{tikzpicture}
\end{center}
\legende{Production d'acide lactique (lactate sanguin) en fonction de l'intensité de l'effort. Le seuil anaérobie (lactate = 4 mmol/L) est atteint vers 57 \% de l'effort maximal chez le sédentaire, mais seulement vers 82 \% chez l'athlète : l'entraînement retarde l'installation de la fermentation lactique.}
\end{popfigure}

\begin{savaistu}[Des moteurs biologiques exceptionnels]
Les grands marathoniens et cyclistes de haut niveau ont une VO2 max pouvant dépasser \SI{80}{mL} d'\ce{O2}/min/kg, contre environ \SI{40}{mL}/min/kg chez un adulte sédentaire. Certains champions d'endurance approchent même \SI{90}{mL}/min/kg ! Cette capacité, en partie héréditaire, est surtout le fruit d'années d'entraînement qui ont multiplié les mitochondries de leurs fibres musculaires. Les coureurs de fond camerounais qui s'entraînent sur les pentes du Mont Cameroun lors de la Course de l'Espoir bénéficient en plus de l'altitude, qui stimule la production de globules rouges.
\end{savaistu}

\courssub{L'alimentation : fournir le carburant}

Sans glucose, pas de glycolyse, donc pas de respiration. L'alimentation conditionne directement le rendement énergétique :

\begin{itemize}
\item les \cle{glucides} (riz, igname, manioc, maïs, pain) sont le carburant privilégié de l'effort : ils reconstituent les réserves de glycogène musculaire et hépatique ;
\item pendant un effort prolongé, les \cle{boissons sucrées} apportent du glucose rapidement disponible et retardent l'épuisement des réserves ;
\item une alimentation équilibrée fournit aussi les \cle{vitamines du groupe B} (B1, B2, B3...), indispensables car elles constituent des \emph{cofacteurs} des enzymes du cycle de Krebs et de la chaîne respiratoire (NAD, FAD sont dérivés de vitamines B2 et B3).
\end{itemize}

\courssub{Autres facteurs}

\begin{itemize}
\item \textbf{Température} : les enzymes ont une température optimale (environ \SI{37}{\celsius}) ; la fièvre ou l'hypothermie perturbent leur activité.
\item \textbf{Disponibilité en enzymes et cofacteurs} : une carence en vitamines B ou en fer (anémie, fréquente en cas de paludisme répété) réduit la capacité de production d'ATP ou de transport de l'\ce{O2}.
\item \textbf{État de santé} : les maladies respiratoires, cardiaques ou les infections diminuent l'oxygénation et donc le rendement énergétique.
\end{itemize}

\courssub{Exploitation de données expérimentales : consommation d'\ce{O2} et production de \ce{CO2}}

Au Probatoire, on te donnera souvent un tableau ou une courbe de consommation d'\ce{O2} et de production de \ce{CO2} au repos et à l'effort. Voici la démarche à adopter.

\begin{methode}[Exploiter des données d'échanges gazeux]
\begin{enumerate}
\item \textbf{Lire} les valeurs de consommation d'\ce{O2} et de production de \ce{CO2} dans chaque condition (repos, effort), en vérifiant les unités.
\item \textbf{Comparer} : calculer de combien chaque valeur est multipliée entre le repos et l'effort.
\item \textbf{Calculer le quotient respiratoire} : $QR = \dfrac{\text{volume de } \ce{CO2} \text{ produit}}{\text{volume de } \ce{O2} \text{ consommé}}$.
\item \textbf{Conclure} : si $QR \approx 1$, ce sont des glucides qui sont oxydés (d'après le bilan $\ce{C6H12O6 + 6O2 -> 6CO2 + 6H2O}$, 6 mol de \ce{CO2} pour 6 mol d'\ce{O2}). Relier l'augmentation des échanges gazeux à l'augmentation des besoins en ATP des cellules musculaires.
\end{enumerate}
\end{methode}

\begin{exoresolu}[Calcul d'un quotient respiratoire]
\textbf{Énoncé.} On mesure chez un élève au repos, pendant une minute : consommation d'\ce{O2} = \SI{0,30}{L} ; production de \ce{CO2} = \SI{0,30}{L}. Après ingestion d'un repas riche en glucides, on mesure sur un muscle en activité : 6 mol d'\ce{O2} consommées et 6 mol de \ce{CO2} produites. Calculer les QR et conclure.
\tcblower
\textbf{Solution.} Au repos : $QR = \dfrac{0{,}30}{0{,}30} = 1$. Sur le muscle : $QR = \dfrac{6}{6} = 1$.

Dans les deux cas, $QR \approx 1$ : le substrat oxydé est un \textbf{glucide} (glucose). En effet, le bilan de la respiration cellulaire $\ce{C6H12O6 + 6O2 -> 6CO2 + 6H2O}$ montre que l'oxydation complète du glucose produit autant de \ce{CO2} qu'elle consomme d'\ce{O2}. (À titre de comparaison, l'oxydation des lipides donne $QR \approx 0{,}7$.)
\end{exoresolu}

\begin{attention}[Piège classique]
Une consommation d'\ce{O2} élevée ne signifie \textbf{pas} forcément un bon rendement énergétique ! Ce qui compte, c'est la quantité d'\textbf{ATP produite par molécule de glucose}. La respiration produit environ 36 ATP par glucose, la fermentation seulement 2 ATP \emph{sans consommer d'\ce{O2}}. Un muscle qui fermente peut dégrader énormément de glucose très vite avec un rendement misérable. Raisonne toujours sur le couple (glucose consommé, ATP produit), jamais sur la seule consommation d'\ce{O2}.
\end{attention}

\begin{aretenir}[L'essentiel de la section]
\begin{itemize}
\item L'apport d'\ce{O2} dépend de la ventilation pulmonaire et de la circulation sanguine ; l'essoufflement traduit l'augmentation des besoins en \ce{O2}.
\item Le seuil anaérobie marque le basculement vers la fermentation lactique ; l'entraînement le recule en multipliant les mitochondries et en augmentant la VO2 max.
\item L'alimentation fournit le glucose (carburant) et les vitamines B (cofacteurs enzymatiques).
\item $QR = \ce{CO2}/\ce{O2} \approx 1$ pour l'oxydation des glucides : outil clé d'exploitation des données au Probatoire.
\end{itemize}
\end{aretenir}

\begin{exercice}[Analyse de courbes]
À partir des deux courbes de cette section : (1) déterminer graphiquement la VO2 max de chaque sujet ; (2) indiquer à quelle intensité d'effort le sujet sédentaire atteint son seuil anaérobie ; (3) expliquer en trois lignes pourquoi l'athlète peut courir plus longtemps sans douleurs musculaires.
\end{exercice}

\begin{corrige}[Exercice : Analyse de courbes]
(1) VO2 max : environ \SI{80}{mL}/min/kg pour l'athlète et \SI{40}{mL}/min/kg pour le sédentaire (plateaux des courbes). (2) Le seuil anaérobie du sédentaire est atteint vers 57 \% de l'effort maximal (lactate = 4 mmol/L). (3) L'athlète possède davantage de mitochondries et une meilleure oxygénation : ses cellules musculaires produisent l'ATP par respiration aérobie jusqu'à des efforts intenses, ce qui retarde l'accumulation d'acide lactique responsable des brûlures et de la fatigue musculaire.
\end{corrige}

\coursec{Synthèse et applications}

Nous avons vu dans la section précédente que le rendement énergétique de la cellule n'est pas une fatalité : l'entraînement, l'alimentation et l'oxygénation permettent de l'améliorer sensiblement. Il est temps maintenant de rassembler toutes les pièces du puzzle. Cette dernière section te donne une vision d'ensemble — du glucose à l'ATP — et te montre comment les connaissances acquises s'appliquent à des situations concrètes de la vie quotidienne au Cameroun : la crampe du sprinter lors de la fête du sport, la levée de la pâte de beignets, la fermentation du vin de palme. C'est aussi la section qui te prépare directement aux questions de synthèse du Probatoire.

\courssub{Vision d'ensemble : du glucose à l'ATP}

\begin{aretenir}[L'idée directrice de la leçon]
Toute cellule a besoin d'\cle{ATP} (adénosine triphosphate) pour fonctionner : contraction musculaire, transport actif, synthèses, division cellulaire. Pour produire cet ATP, la cellule dégrade le \cle{glucose} selon deux stratégies :
\begin{itemize}
\item la \cle{respiration cellulaire}, en présence de dioxygène : oxydation complète du glucose jusqu'à \ce{CO2} et \ce{H2O}, avec un rendement élevé (environ 36 à 38 ATP par glucose) ;
\item la \cle{fermentation}, en absence de dioxygène : dégradation incomplète du glucose, avec un rendement faible (2 ATP par glucose) et la production d'un déchet (lactate chez l'Homme, éthanol et \ce{CO2} chez la levure).
\end{itemize}
Dans les deux cas, la première étape est la \cle{glycolyse}, qui se déroule dans le cytoplasme. C'est le destin du pyruvate, produit de la glycolyse, qui détermine la voie suivie.
\end{aretenir}

La carte mentale ci-dessous résume l'architecture complète de la leçon. Prends le temps de la parcourir : chaque flèche correspond à une relation de cause à effet que tu dois être capable de reformuler avec tes propres mots.

\begin{popfigure}
\begin{center}
\begin{tikzpicture}[
  mindmap,
  grow cyclic,
  every node/.style={concept, font=\small},
  root concept/.append style={concept color=popblue, font=\bfseries},
  level 1/.append style={level distance=4.2cm, sibling angle=120},
  level 2/.append style={level distance=2.8cm, sibling angle=50},
  level 1 concept/.append style={font=\small\bfseries},
  level 2 concept/.append style={font=\scriptsize}
]
\node[root concept, align=center]{Glucose \\ \ce{C6H12O6}}
  child[concept color=popgreen]{ node {Respiration (avec \ce{O2})}
    child { node {Glycolyse (cytoplasme) : 2 ATP} }
    child { node {Cycle de Krebs (matrice) : 2 ATP + coenzymes réduits} }
    child { node[align=center]{Chaîne respiratoire (crêtes) : ~34 ATP \\ \ce{O2} accepteur final} }
  }
  child[concept color=poporange]{ node {Fermentation (sans \ce{O2})}
    child { node {Lactique : lactate (muscle)} }
    child { node {Alcoolique : éthanol + \ce{CO2} (levure)} }
    child { node {Rendement : 2 ATP} }
  }
  child[concept color=poppurple]{ node {Optimisation}
    child { node {Entraînement} }
    child { node {Alimentation} }
    child { node {Oxygénation} }
  };
\end{tikzpicture}
\end{center}
\legende{Carte mentale de synthèse : le glucose, au centre, est dégradé selon deux voies (respiration aérobie, fermentation anaérobie) dont le rendement en ATP diffère fortement ; trois facteurs périphériques permettent d'optimiser la production d'énergie.}
\end{popfigure}

\courssub{Tableau comparatif final : respiration contre fermentation}

Ce tableau est un grand classique des sujets de Probatoire. Apprends-le ligne par ligne, et surtout comprends chaque case : un examinateur peut te demander de justifier n'importe laquelle d'entre elles.

\begin{center}
\begin{tabularx}{\linewidth}{|l|X|X|}
\hline
\rowcolor{popblueL} \textbf{Critère} & \textbf{Respiration cellulaire} & \textbf{Fermentation} \\
\hline
Conditions & Présence de dioxygène (aérobie) & Absence de dioxygène (anaérobie) \\
\hline
Localisation & Cytoplasme (glycolyse) puis mitochondrie (cycle de Krebs, chaîne respiratoire) & Cytoplasme uniquement \\
\hline
Organite impliqué & Mitochondrie indispensable & Aucun organite \\
\hline
Substrat dégradé & Glucose totalement oxydé & Glucose partiellement dégradé \\
\hline
Produits finaux & \ce{CO2} + \ce{H2O} & Lactate (cellule musculaire) ou éthanol + \ce{CO2} (levure) \\
\hline
Rendement & Environ 36 à 38 ATP par glucose & 2 ATP par glucose \\
\hline
Vitesse de production & Lente mais durable & Rapide mais épuisable (accumulation de déchet) \\
\hline
\end{tabularx}
\end{center}

\begin{attention}[Les erreurs qui coûtent des points]
\begin{itemize}
\item Ne dis jamais que la fermentation « ne produit pas d'ATP » : elle produit 2 ATP par glucose, grâce à la glycolyse. C'est peu, mais ce n'est pas zéro.
\item La glycolyse est \textbf{commune} aux deux voies : elle n'a pas besoin de dioxygène. Dire « la glycolyse a lieu dans la mitochondrie » est faux : elle se déroule dans le \textbf{cytoplasme}.
\item La fermentation ne produit \textbf{pas} de \ce{CO2} chez la cellule musculaire humaine (fermentation lactique). Le \ce{CO2} n'apparaît que dans la fermentation alcoolique des levures.
\item Le muscle ne « manque pas de glucose » pendant un effort intense : c'est l'apport en \ce{O2} qui devient insuffisant.
\end{itemize}
\end{attention}

\courssub{L'équation-bilan : la phrase à maîtriser parfaitement}

\begin{propriete}[Équation-bilan de la respiration cellulaire]
L'oxydation complète d'une molécule de glucose en présence de dioxygène s'écrit :
\[ \ce{C6H12O6 + 6O2 -> 6CO2 + 6H2O + Energie (ATP)} \]
Cette équation est \textbf{exigible au Probatoire}, avec ses coefficients exacts.
\end{propriete}

\begin{attention}[Équilibre de l'équation : les coefficients sont exigés]
Vérifie toujours l'équilibre atome par atome :
\begin{itemize}
\item Carbone : 6 à gauche (dans \ce{C6H12O6}), 6 à droite (dans $6\,\ce{CO2}$) ;
\item Hydrogène : 12 à gauche, $6 \times 2 = 12$ à droite (dans $6\,\ce{H2O}$) ;
\item Oxygène : $6 + 12 = 18$ à gauche, $12 + 6 = 18$ à droite.
\end{itemize}
Les erreurs fréquentes : écrire \ce{O2} sans le coefficient 6, oublier l'eau, ou inverser réactifs et produits. Une équation non équilibrée est sanctionnée, même si l'idée est bonne.
\end{attention}

\begin{methode}[Rédiger une réponse complète de type Probatoire sur le bilan de la respiration]
Face à une question du type « Décrire les étapes de la production d'ATP par la cellule musculaire à l'effort et donner le bilan énergétique », structure ta réponse en cinq temps :
\begin{enumerate}
\item \textbf{Poser le cadre} : la cellule produit de l'ATP par dégradation du glucose ; en présence d'\ce{O2}, c'est la respiration cellulaire.
\item \textbf{Écrire l'équation-bilan équilibrée} : \ce{C6H12O6 + 6O2 -> 6CO2 + 6H2O} + énergie.
\item \textbf{Détailler les trois étapes} en précisant toujours le lieu : glycolyse (cytoplasme, glucose $\to$ 2 pyruvates, 2 ATP) ; cycle de Krebs (matrice mitochondriale, oxydation du pyruvate, libération de \ce{CO2} et production de coenzymes réduits) ; chaîne respiratoire (membrane interne des crêtes, l'\ce{O2} est l'accepteur final d'électrons, formation d'eau et de la quasi-totalité de l'ATP).
\item \textbf{Donner le rendement chiffré} : environ 36 à 38 ATP par glucose, contre 2 ATP en fermentation.
\item \textbf{Conclure en revenant à la situation} : c'est ce rendement élevé qui permet l'effort prolongé, et c'est l'insuffisance d'apport en \ce{O2} qui déclenche la fermentation et la fatigue.
\end{enumerate}
Rédige en phrases complètes, nomme les structures précisément (matrice, crêtes mitochondriales), et appuie-toi sur les données de l'énoncé s'il y en a.
\end{methode}

\courssub{Applications au quotidien camerounais}

\textbf{La crampe du sprinter.} Lors d'un 100 m à la fête du sport du lycée, l'effort est si violent que le sang ne peut pas apporter l'oxygène assez vite aux fibres musculaires. Celles-ci basculent en \cle{fermentation lactique} : le lactate s'accumule, le pH intracellulaire chute, d'où la brûlure musculaire puis la crampe. Le sprinter reste ensuite essoufflé plusieurs minutes : il « rembourse sa dette d'oxygène » en oxydant le lactate accumulé.

\textbf{L'essoufflement.} L'accélération des battements du cœur et de la ventilation pendant et après l'effort n'est pas un hasard : elle augmente l'apport de \ce{O2} aux mitochondries et l'élimination du \ce{CO2}. C'est la preuve physiologique directe de l'équation-bilan : plus on produit d'ATP par respiration, plus on consomme d'\ce{O2} et plus on rejette de \ce{CO2}.

\textbf{Le vin de palme.} La sève de palmier récoltée dans la calebasse contient des sucres et des levures naturelles. À l'abri de l'air, les levures réalisent la \cle{fermentation alcoolique} : glucose $\to$ éthanol + \ce{CO2} + 2 ATP. L'éthanol donne au vin de palme son degré alcoolique ; le \ce{CO2} produit le pétillement caractéristique. Si la fermentation se prolonge, des bactéries acétiques transforment l'éthanol en acide acétique : le vin « tourne » et devient vinaigre.

\textbf{La levée de la pâte.} Pour les beignets, les puff-puff ou le pain, la levure de boulanger fermente les sucres de la farine. Le \ce{CO2} dégagé forme des bulles prisonnières du réseau de gluten : la pâte gonfle. À la cuisson, l'éthanol s'évapore et les bulles donnent la mie aérée. Même principe, même équation : la fermentation alcoolique.

\textbf{La stratégie de course.} Un coureur de fond qui part trop vite épuise ses réserves en fermentation lactique et « explose » avant la fin. La tactique gagnante consiste à courir juste en dessous du seuil où la ventilation ne suffit plus, afin de rester en filière aérobie le plus longtemps possible, puis de sprinter — en anaérobie — dans les derniers hectomètres.

\begin{exemplebox}[Sprint de 100 m contre marathon : deux stratégies énergétiques]
Comparons les bilans énergétiques de deux efforts extrêmes.
\begin{itemize}
\item \textbf{Sprint de 100 m (environ 10 s)} : la demande en ATP est massive et immédiate. La respiration, trop lente à s'ajuster, ne peut suivre : la \textbf{fermentation lactique domine}. Rendement : 2 ATP par glucose. La consommation de glucose est énorme pour peu d'ATP, et le lactate s'accumule.
\item \textbf{Marathon (environ 2 h 30)} : l'effort est modéré mais prolongé. La \textbf{respiration domine} : chaque glucose fournit 36 à 38 ATP, soit \textbf{18 fois plus} d'énergie par molécule. Le coureur puise aussi dans les lipides, dégradés par la même machinerie mitochondriale.
\end{itemize}
Conclusion chiffrée : pour produire la même quantité d'ATP, la fermentation consomme environ 18 fois plus de glucose que la respiration. C'est pourquoi un effort anaérobie ne peut durer que quelques dizaines de secondes, alors qu'un effort aérobie peut durer des heures.
\end{exemplebox}

\begin{savaistu}[Une masse d'ATP égale à ton poids, chaque jour]
L'ATP n'est pas stocké en grande quantité : une cellule n'en contient que quelques secondes d'autonomie. Il est donc recyclé en permanence : ATP $\to$ ADP + phosphate $\to$ ATP. On estime qu'un adulte au repos hydrolyse et re-synthétise environ \textbf{40 kg d'ATP par jour}, et jusqu'à \num{0,5} kg \emph{par minute} lors d'un effort intense — soit, sur une journée, une masse d'ATP recyclée de l'ordre de la masse corporelle elle-même. La mitochondrie est vraiment une usine qui ne s'arrête jamais.
\end{savaistu}

\courssub{Lien avec la famille de situations : optimiser la production d'énergie cellulaire}

Toute la leçon se résume en un problème pratique : \emph{comment un organisme peut-il produire plus d'ATP, plus longtemps ?} La réponse tient en trois leviers, que tu dois savoir relier aux mécanismes cellulaires :

\begin{itemize}
\item \textbf{Améliorer l'oxygénation} : une ventilation efficace et un sang riche en hémoglobine garantissent l'apport de \ce{O2}, accepteur final d'électrons de la chaîne respiratoire. Sans \ce{O2}, toute la filière aérobie s'arrête.
\item \textbf{S'entraîner régulièrement} : l'entraînement augmente le nombre et la taille des mitochondries dans les fibres musculaires, ainsi que la densité des capillaires sanguins. La cellule produit alors plus d'ATP par respiration et repousse le seuil de bascule vers la fermentation.
\item \textbf{Bien s'alimenter} : les glucides (manioc, igname, maïs, riz) fournissent le glucose, carburant de la glycolyse ; une alimentation équilibrée fournit aussi les vitamines et minéraux nécessaires aux enzymes et coenzymes des étapes énergétiques.
\end{itemize}

\begin{exoresolu}[Exploitation de données expérimentales]
Énoncé. On mesure chez un élève au repos puis après 5 minutes de course intense les échanges gazeux respiratoires :

\begin{center}
\begin{tabular}{|l|c|c|}
\hline
\rowcolor{popblueL} & \textbf{\ce{O2} consommé (L/min)} & \textbf{\ce{CO2} rejeté (L/min)} \\
\hline
Repos & 0,25 & 0,20 \\
\hline
Après l'effort & 2,80 & 3,10 \\
\hline
\end{tabular}
\end{center}

1. Comparer les échanges gazeux au repos et après l'effort. 2. Interpréter ces variations en utilisant l'équation-bilan de la respiration. 3. Expliquer pourquoi le rapport \ce{CO2}/\ce{O2} dépasse 1 après l'effort.

\tcblower
Solution détaillée.

1. Après l'effort, la consommation de \ce{O2} est multipliée par $\num{2,80} / \num{0,25} = \num{11,2}$ et le rejet de \ce{CO2} par $\num{3,10} / \num{0,20} = \num{15,5}$. Les deux échanges gazeux augmentent fortement.

2. Pendant l'effort, les cellules musculaires produisent beaucoup plus d'ATP par respiration cellulaire. D'après l'équation-bilan \ce{C6H12O6 + 6O2 -> 6CO2 + 6H2O}, toute augmentation de l'oxydation du glucose se traduit par une consommation accrue de \ce{O2} et un rejet accru de \ce{CO2}. Les mesures confirment donc l'intensification de la respiration cellulaire.

3. Le rapport $\ce{CO2}/\ce{O2} = \num{3,10} / \num{2,80} \approx \num{1,1} > 1$ : le corps rejette plus de \ce{CO2} qu'il ne consomme d'\ce{O2}. Or, d'après l'équation-bilan, la respiration seule donnerait un rapport de $6/6 = 1$. L'excès de \ce{CO2} provient de la fermentation lactique réalisée pendant l'effort : le lactate accumulé est neutralisé par les bicarbonates du sang, ce qui libère du \ce{CO2} supplémentaire. Ce rapport supérieur à 1 est donc la signature d'une part d'anaérobiose dans l'effort.
\end{exoresolu}

\begin{exercice}[Pour t'entraîner]
Une levure est placée dans deux milieux contenant chacun 10 g de glucose : le milieu A est barboté d'air, le milieu B est maintenu sans oxygène. 1. Prévoir, en le justifiant, quel milieu épuisera son glucose le plus vite pour produire la même quantité d'ATP. 2. Nommer les produits formés dans le milieu B. 3. Citer une application alimentaire camerounaise de ce phénomène.
\end{exercice}

\begin{corrige}[Exercice]
1. Le milieu B (anaérobie) épuisera son glucose plus vite : la fermentation ne fournit que 2 ATP par glucose contre 36 à 38 en respiration ; pour produire la même quantité d'ATP, la levure doit donc dégrader environ 18 fois plus de glucose. 2. Dans le milieu B, la fermentation alcoolique produit de l'éthanol et du \ce{CO2}. 3. La préparation du vin de palme ou la levée de la pâte à beignets.
\end{corrige}

\begin{aretenir}[Ce qu'il faut retenir de toute la leçon]
\begin{itemize}
\item Équation-bilan de la respiration : \ce{C6H12O6 + 6O2 -> 6CO2 + 6H2O} + énergie, à écrire avec ses coefficients.
\item Trois étapes : glycolyse (cytoplasme), cycle de Krebs (matrice), chaîne respiratoire (crêtes mitochondriales, \ce{O2} accepteur final).
\item Rendements : 36 à 38 ATP par glucose en respiration ; 2 ATP en fermentation.
\item La mitochondrie est la centrale énergétique de la cellule ; son abondance augmente avec l'entraînement.
\item Optimiser l'énergie cellulaire = améliorer l'oxygénation, l'entraînement et l'alimentation.
\end{itemize}
\end{aretenir}

\newpage

\coursec{Activité d'intégration}

\courssub{Objectifs de l'activité}

À l'issue de cette activité d'intégration, tu seras capable de :

\begin{itemize}
\item \cle{exploiter des données expérimentales} de consommation de dioxygène ($O_2$) et de production de dioxyde de carbone ($CO_2$) mesurées chez des sujets humains au repos et à l'effort ;
\item \cle{calculer un quotient respiratoire} (QR) et en déduire la nature du nutriment majoritairement oxydé par l'organisme ;
\item \cle{comparer les rendements énergétiques} de la respiration cellulaire et de la fermentation lactique (36 à 38 ATP contre 2 ATP par molécule de glucose) ;
\item \cle{expliquer un phénomène physiologique} concret (crampes, essoufflement, dette d'oxygène) à l'aide du bilan simplifié de la respiration cellulaire ;
\item \cle{formuler des conseils} d'entraînement, d'alimentation et d'oxygénation fondés sur des arguments scientifiques, afin d'optimiser la production d'énergie par l'organisme.
\end{itemize}

\courssub{Situation}

\textbf{Le défi énergétique du sprinteur de Yaoundé.}

Dans le cadre de la préparation des jeux inter-lycées du Centre, le club d'athlétisme du Lycée Général Leclerc de Yaoundé organise une séance de tests physiologiques au laboratoire de SVT. Deux élèves de Première C se portent volontaires :

\begin{itemize}
\item \textbf{Sujet A (Junior)} : membre du club d'athlétisme, il s'entraîne trois fois par semaine depuis deux ans ;
\item \textbf{Sujet B (Rodrigue)} : élève non entraîné, qui ne pratique aucune activité sportive régulière.
\end{itemize}

Les deux sujets, de masse corporelle voisine (environ \SI{60}{kg}), sont équipés d'un masque relié à un analyseur de gaz. On mesure leur consommation de dioxygène ($\dot{V}O_2$) et leur production de dioxyde de carbone ($\dot{V}CO_2$) au repos (assis, 10 minutes), puis pendant une course de 3 minutes à vitesse imposée. On prélève également un échantillon de sang au bout du doigt avant l'effort et 3 minutes après la fin de la course, afin de doser l'acide lactique sanguin. Les résultats sont consignés dans le tableau ci-dessous.

\begin{center}
\rowcolors{1}{popblueL}{white}
\begin{tabularx}{\linewidth}{|l|X|X|}
\hline
\rowcolor{popblue!40} \textbf{Grandeur mesurée} & \textbf{Sujet A (entraîné)} & \textbf{Sujet B (non entraîné)} \\
\hline
$\dot{V}O_2$ au repos (\SI{}{mL/min}) & 250 & 250 \\
\hline
$\dot{V}CO_2$ au repos (\SI{}{mL/min}) & 200 & 200 \\
\hline
$\dot{V}O_2$ pendant la course (\SI{}{mL/min}) & 2400 & 1800 \\
\hline
$\dot{V}CO_2$ pendant la course (\SI{}{mL/min}) & 2280 & 1980 \\
\hline
Acide lactique sanguin avant effort (\SI{}{mmol/L}) & 1,0 & 1,0 \\
\hline
Acide lactique sanguin après effort (\SI{}{mmol/L}) & 4,5 & 12,0 \\
\hline
Observations après la course & Récupération rapide, respiration normale au bout de 2 min & Essoufflement prolongé, crampes douloureuses aux mollets \\
\hline
\end{tabularx}
\end{center}
\legende{Données physiologiques recueillies au repos et lors d'un effort standardisé (course 3 min) chez un sujet entraîné (A) et un sujet non entraîné (B).}

Le professeur de SVT, qui encadre la séance, pose alors le problème suivant à la classe : \emph{« Comment expliquer scientifiquement la différence de performance et d'état final entre Junior et Rodrigue, et que conseiller à Rodrigue pour améliorer son rendement énergétique ? »}

\courssub{Consignes}

Par groupes de quatre élèves, traitez les questions suivantes, puis rédigez un mini-rapport (une page maximum) qui sera affiché et présenté en classe.

\begin{enumerate}
\item Calculez le \cle{quotient respiratoire} (QR) de chaque sujet au repos, défini par $\mathrm{QR} = \dfrac{\dot{V}CO_2}{\dot{V}O_2}$. En vous aidant des valeurs de référence (QR = 1 pour l'oxydation des glucides ; QR = 0,7 pour l'oxydation des lipides), identifiez le nutriment majoritairement oxydé par les deux sujets au repos. Justifiez votre réponse à partir du bilan de la respiration cellulaire.
\item Comparez l'évolution de la consommation de $O_2$ entre le repos et la course chez les deux sujets (calculez le facteur d'augmentation pour chacun). Calculez ensuite le QR de chaque sujet pendant la course. Quel sujet est resté le plus longtemps en \cle{respiration aérobie} ? Justifiez.
\item À l'aide du bilan simplifié de la respiration cellulaire et du rendement de la fermentation lactique (36 à 38 ATP contre 2 ATP par molécule de glucose), expliquez pourquoi le sujet B présente un taux d'acide lactique élevé et des crampes après l'effort.
\item Rédigez un court texte de conseils (8 à 12 lignes) destiné à Rodrigue, portant sur l'entraînement, l'alimentation et l'oxygénation, afin d'améliorer le rendement énergétique de ses cellules musculaires.
\end{enumerate}

\courssub{Ressources à mobiliser}

\begin{aretenir}[Boîte à outils de la leçon]
\begin{itemize}
\item \textbf{Bilan simplifié de la respiration cellulaire} :
\[ \ce{C6H12O6 + 6O2 -> 6CO2 + 6H2O} + \text{énergie (36 à 38 ATP)} \]
\item \textbf{Bilan de la fermentation lactique} (en absence d'oxygène) :
\[ \ce{C6H12O6 -> 2 C3H6O3} \text{ (acide lactique)} + \text{énergie (2 ATP)} \]
\item \textbf{Quotient respiratoire} : $\mathrm{QR} = \dfrac{\dot{V}CO_2 \text{ produit}}{\dot{V}O_2 \text{ consommé}}$. D'après le bilan de la respiration, l'oxydation complète du glucose donne $\mathrm{QR} = \dfrac{6}{6} = 1$ ; celle des lipides donne $\mathrm{QR} \approx 0,7$.
\item \textbf{Localisation} : la glycolyse se déroule dans le \cle{cytosol} (hyaloplasme) ; la décarboxylation oxydative du pyruvate, le cycle de Krebs et la chaîne respiratoire ont lieu dans les \cle{mitochondries} (matrice et crêtes mitochondriales).
\item \textbf{Facteurs du rendement énergétique} : entraînement (augmentation du nombre et de la surface des crêtes mitochondriales), alimentation (apport de glucides complexes pour les réserves de glycogène), oxygénation (ventilation efficace, transport du $O_2$ par l'hémoglobine, capillarisation musculaire).
\end{itemize}
\end{aretenir}

\courssub{Démarche et corrigé commenté}

\begin{exoresolu}[Le défi énergétique du sprinteur de Yaoundé]
Énoncé : à partir des données du tableau ci-dessus, répondre aux quatre consignes et rédiger le mini-rapport.

\tcblower

\textbf{Étape 1 — Calcul et interprétation du quotient respiratoire au repos.}

Pour chaque sujet au repos :
\[ \mathrm{QR} = \frac{\dot{V}CO_2}{\dot{V}O_2} = \frac{200}{250} = 0,8 \]
Les deux sujets ont le même QR de repos, égal à 0,8. Cette valeur est intermédiaire entre 0,7 (oxydation quasi exclusive des lipides) et 1 (oxydation quasi exclusive des glucides) : au repos, l'organisme oxyde un \textbf{mélange de glucides et de lipides}, avec une part notable de lipides, ce qui est normal pour un métabolisme de base. On peut vérifier la cohérence avec le bilan de la respiration : pour le glucose, $\ce{C6H12O6 + 6O2 -> 6CO2 + 6H2O}$, on produit autant de $CO_2$ qu'on consomme de $O_2$, d'où QR = 1 ; les lipides, plus riches en hydrogène, exigent davantage de $O_2$ par $CO_2$ produit, d'où un QR plus faible (environ 0,7 pour le tripalmitine).

\textbf{Étape 2 — Comparaison de la consommation de $O_2$ à l'effort et analyse du QR.}

Facteur d'augmentation de la consommation de $O_2$ (passage repos $\rightarrow$ course) :
\begin{align*}
\text{Sujet A : } & \frac{2400}{250} = 9,6 \\
\text{Sujet B : } & \frac{1800}{250} = 7,2
\end{align*}
Le sujet A multiplie sa consommation de $O_2$ par 9,6, le sujet B seulement par 7,2 : l'organisme entraîné est capable d'apporter beaucoup plus de dioxygène à ses muscles (meilleure ventilation, meilleur transport sanguin, meilleure extraction musculaire).

Quotient respiratoire pendant la course :
\begin{align*}
\text{Sujet A : } & \mathrm{QR} = \frac{2280}{2400} = 0,95 \\
\text{Sujet B : } & \mathrm{QR} = \frac{1980}{1800} = 1,1
\end{align*}
Le sujet A présente un QR proche de 1 : il oxyde presque exclusivement des \cle{glucides}, le carburant privilégié de l'effort intense, et l'essentiel de son énergie provient encore de la respiration cellulaire complète (aérobie).
Le sujet B présente un QR supérieur à 1, ce qui est physiologiquement impossible pour la seule oxydation des nutriments (le bilan chimique ne permet pas de produire plus de $CO_2$ que de $O_2$ consommé). Cet excès de $CO_2$ mesuré provient du \textbf{tamponnage de l'acide lactique} accumulé par les bicarbonates sanguins : $\ce{H+ + HCO3- -> H2CO3 -> CO2 + H2O}$. Ce $CO_2$ « non métabolique » s'ajoute au $CO_2$ respiratoire, faisant artificiellement grimper le QR au-dessus de 1. C'est la signature d'un recours massif à la fermentation lactique.
\textbf{Conclusion : le sujet A est resté le plus longtemps en respiration aérobie.}

\textbf{Étape 3 — Explication des crampes et de l'essoufflement du sujet B.}

Pendant un effort intense, la demande en ATP des fibres musculaires explose. Chez le sujet B, l'apport en $O_2$ (limité à \SI{1800}{mL/min}) ne suffit pas à alimenter la chaîne respiratoire des mitochondries pour régénérer l'ATP au rythme requis. Les cellules musculaires basculent alors vers la \cle{fermentation lactique}, voie alternative sans oxygène :
\[ \ce{C6H12O6 -> 2 C3H6O3} + 2 \text{ ATP} \]
Or cette voie est très peu rentable : 2 ATP par glucose, contre 36 à 38 ATP pour la respiration complète, soit un rendement environ \textbf{18 fois inférieur}. Pour compenser ce faible rendement et tenter de maintenir la production d'ATP, la cellule consomme énormément de glucose (glycogène musculaire) et produit en masse de l'\cle{acide lactique}, qui s'accumule dans le muscle puis diffuse dans le sang : le taux sanguin du sujet B passe de 1,0 à 12,0 mmol/L (contre 4,5 mmol/L seulement chez le sujet A).
L'acidification du milieu musculaire (chute du pH) perturbe le fonctionnement des enzymes glycolytiques, l'interaction actine-myosine et la libération du calcium : c'est la cause directe des \textbf{crampes douloureuses}. L'essoufflement prolongé traduit le remboursement de la \cle{dette d'oxygène} : l'organisme doit consommer un excès d'$O_2$ après l'effort pour réoxyder l'acide lactique accumulé (néoglucogenèse hépatique ou oxydation directe) et restaurer les réserves.

\textbf{Étape 4 — Conseils à Rodrigue (exemple de rédaction attendue).}

\emph{« Rodrigue, pour améliorer ton rendement énergétique et tes performances : (1) \textbf{Entraîne-toi régulièrement} (2 à 3 séances d'endurance par semaine : footing, football, vélo) : l'entraînement d'endurance augmente le nombre, la taille et la densité en crêtes des mitochondries dans tes fibres musculaires, ce qui élève ta $\dot{V}O_2$ max et ta capacité à produire de l'ATP en aérobie. (2) \textbf{Soigne ton alimentation} : consomme des glucides complexes (riz, igname, manioc, banane plantain, maïs) dans les repas précédant l'effort pour maximiser tes réserves de glycogène musculaire et hépatique ; limite les graisses juste avant la course. (3) \textbf{Optimise ton oxygénation} : apprends à respirer profondément (ventilation ample), échauffe-toi progressivement pour préparer le système cardio-vasculaire, et hydrate-toi correctement pour maintenir le volume sanguin et le transport de l'$O_2$ par l'hémoglobine. Avec ces habitudes, tu retarderas le seuil de basculement vers la fermentation lactique, tu produiras plus d'ATP par molécule de glucose, et tu éviteras l'accumulation d'acide lactique responsable des crampes et de la fatigue précoce. »}
\end{exoresolu}

\courssub{Bilan de l'activité}

Cette activité a permis de mettre en évidence, sur une situation authentique de la vie scolaire camerounaise, les points essentiels de la leçon :

\begin{itemize}
\item le \cle{quotient respiratoire} est un outil simple qui renseigne sur la nature des nutriments oxydés (QR $\approx$ 0,7 lipides ; QR $\approx$ 1 glucides) et permet de détecter, lorsqu'il dépasse 1, un recours massif à la fermentation lactique (excès de $CO_2$ tampon) ;
\item la \cle{respiration cellulaire} (glycolyse dans le cytosol, cycle de Krebs dans la matrice mitochondriale, chaîne respiratoire sur les crêtes) est la voie majeure de production d'ATP : 36 à 38 ATP par glucose, soit un rendement environ 18 fois supérieur à celui de la fermentation lactique (2 ATP) ;
\item la \cle{fermentation lactique} n'est qu'une voie de secours anaerobie : elle permet de produire rapidement un peu d'ATP sans oxygène, mais au prix d'une accumulation d'acide lactique, source d'acidose musculaire, de fatigue, de crampes et de dette d'oxygène ;
\item le \cle{rendement énergétique} de la cellule musculaire n'est pas figé : il dépend de l'entraînement (biogenèse mitochondriale, capillarisation), de l'alimentation (disponibilité en substrats glucidiques) et de l'oxygénation (ventilation, diffusion, transport sanguin, extraction musculaire).
\end{itemize}

Tu as ainsi mobilisé les trois savoir-faire exigibles de la leçon — écrire le bilan de la respiration, comparer les rendements respiration/fermentation, exploiter des données de consommation de $O_2$ et de production de $CO_2$ — pour résoudre un problème concret de performance sportive, exactement comme le ferait un préparateur physique ou un médecin du sport.

\newpage

\coursec{Méthodes et exercices résolus}

\courssub{Méthodes et exercices résolus}

\begin{propriete}[Équation bilan de la respiration cellulaire aérobie]
    L'oxydation complète d'une molécule de glucose en présence de dioxygène se résume par l'équation chimique globale :
    \[ \ce{C6H12O6 + 6 O2 -> 6 CO2 + 6 H2O + Energie} \]
    L'énergie libérée est stockée sous forme d'ATP (environ \num{30} à \num{32} ATP par glucose selon le type cellulaire) et dissipée sous forme de chaleur. Cette réaction est exergonique ($\Delta G^\circ \approx -2870\ \text{kJ/mol}$).
\end{propriete}

\begin{methode}[Écrire le bilan simplifié de la respiration cellulaire]
    \textbf{Objectif :} Établir l'équation bilan globale de la respiration aérobie à partir des substrats et des produits finaux.
    \begin{enumerate}
        \item \textbf{Identifier le substrat principal :} Le glucose (\ce{C6H12O6}) est le carburant de référence (issu du glycogène ou de la glycémie).
        \item \textbf{Identifier le comburant :} Le dioxygène (\ce{O2}), capteur final des électrons en fin de chaîne respiratoire mitochondriale (complexe IV).
        \item \textbf{Identifier les produits finaux :} Le dioxyde de carbone (\ce{CO2}), l'eau (\ce{H2O}) et l'énergie utilisable (ATP + chaleur).
        \item \textbf{Équilibrer l'équation chimique :} Respecter la conservation de la matière (atomes C, H, O) et des charges.
        \item \textbf{Préciser le rendement en ATP :} Environ \num{30} à \num{32} ATP par glucose (selon la navette mitochondriale utilisée pour le transfert des NADH cytoplasmiques).
    \end{enumerate}
    \textbf{Formule à retenir :} \ce{C6H12O6 + 6 O2 -> 6 CO2 + 6 H2O} + Énergie ($\approx$ 30-32 ATP)
\end{methode}

\begin{exoresolu}[Bilan global et calcul du rendement ATP]
    \textbf{Énoncé :} Au cours d'un effort musculaire modéré et prolongé, une cellule musculaire consomme \SI{1,80e-6}{mol} de glucose par minute. La respiration cellulaire est complète.
    \begin{enumerate}
        \item Écrire l'équation bilan de la respiration cellulaire du glucose.
        \item Calculer la quantité d'oxygène consommée et de dioxyde de carbone produite par minute.
        \item En admettant un rendement de \num{30} ATP par glucose, calculer le nombre de molécules d'ATP produites par seconde.
    \end{enumerate}
    \tcblower
    \textbf{Solution détaillée :}
    \begin{enumerate}
        \item \textbf{Équation bilan :} L'oxydation complète du glucose en présence d'oxygène suit l'équation :
              \[ \ce{C6H12O6 + 6 O2 -> 6 CO2 + 6 H2O + Energie (ATP)} \]
        \item \textbf{Calculs stœchiométriques :} D'après l'équation, 1 mol de glucose réagit avec 6 mol de \ce{O2} et produit 6 mol de \ce{CO2}.
              \begin{align*}
                  n(\ce{O2}) &= 6 \times n(\text{glucose}) = 6 \times 1,80 \times 10^{-6} = \mathbf{1,08 \times 10^{-5}\ mol/min} \\
                  n(\ce{CO2}) &= 6 \times n(\text{glucose}) = \mathbf{1,08 \times 10^{-5}\ mol/min}
              \end{align*}
        \item \textbf{Production d'ATP :}
              \begin{itemize}
                  \item Quantité de glucose par seconde : $\frac{1,80 \times 10^{-6}}{60} = 3,00 \times 10^{-8}\ \text{mol/s}$.
                  \item Quantité d'ATP par seconde : $30 \times 3,00 \times 10^{-8} = 9,00 \times 10^{-7}\ \text{mol/s}$.
                  \item Nombre de molécules d'ATP par seconde : $N = n \times N_A = 9,00 \times 10^{-7} \times 6,02 \times 10^{23} = \mathbf{5,42 \times 10^{17}\ \text{molécules/s}}$.
              \end{itemize}
    \end{enumerate}
    \textbf{Conclusion :} La cellule consomme \SI{1,08e-5}{mol/min} d'\ce{O2}, produit \SI{1,08e-5}{mol/min} de \ce{CO2} et synthétise environ \num{5,42e17} molécules d'ATP par seconde.
\end{exoresolu}

\begin{definition}[Quotient Respiratoire (QR)]
    Le quotient respiratoire est le rapport entre le volume de dioxyde de carbone produit et le volume de dioxygène consommé sur une même durée :
    \[ QR = \frac{V_{\ce{CO2}}\text{ produit}}{V_{\ce{O2}}\text{ consommé}} = \frac{\dot{V}_{\ce{CO2}}}{\dot{V}_{\ce{O2}}} \]
    Il permet d'identifier le substrat énergétique principal oxydé : Glucides (QR = 1,0), Lipides (QR $\approx$ 0,70), Protéines (QR $\approx$ 0,80-0,90). Un QR > 1 témoigne d'une hyperventilation (effort intense, tamponnage lactique) ou d'une lipogenèse.
\end{definition}

\begin{methode}[Comparer respiration et fermentation : rendement énergétique]
    \textbf{Objectif :} Opposer les deux voies de dégradation du glucose en termes de substrats, produits, localisation, besoin en \ce{O2} et surtout rendement en ATP.
    \begin{enumerate}
        \item \textbf{Construire un tableau comparatif} avec les critères : Voie métabolique, Localisation, Comburant (\ce{O2}), Produits finaux, Rendement ATP/glucose, Devenir du pyruvate.
        \item \textbf{Rappeler le sort du pyruvate :}
              \begin{itemize}
                  \item Respiration : Pyruvate $\xrightarrow{\text{mitochondrie}}$ Acétyl-CoA $\to$ Cycle de Krebs $\to$ Chaîne respiratoire.
                  \item Fermentation lactique : Pyruvate $\xrightarrow{\text{cytoplasme}}$ Lactate (régénération du NAD$^+$ via lactate déshydrogénase).
                  \item Fermentation alcoolique : Pyruvate $\xrightarrow{\text{cytoplasme}}$ Éthanol + \ce{CO2} (régénération du NAD$^+$).
              \end{itemize}
        \item \textbf{Justifier la différence de rendement :} La respiration oxyde totalement le carbone en \ce{CO2} via le cycle de Krebs et exploite le gradient de protons (chimiosmose, complexe ATP synthase) pour produire \num{~30} ATP. La fermentation ne réalise que la glycolyse (2 ATP nets par phosphorylation au niveau du substrat) et « gâche » l'énergie restante dans les molécules réduites (lactate, éthanol).
        \item \textbf{Calculer le ratio :} Rendement respiration / Rendement fermentation $\approx 30 / 2 = 15$. La respiration est 15 fois plus efficace en rendement molaire.
    \end{enumerate}
\end{methode}

\begin{exoresolu}[Comparaison quantitative lors d'un effort intense]
    \textbf{Énoncé :} Un sprinteur camerounais court le 100 m en 10 s. Ses fibres musculaires rapides utilisent principalement la fermentation lactique. Un marathonien court 42 km en 2 h 10 min, utilisant la respiration aérobie. On admet que chaque coureur dégrade \SI{0,50}{mol} de glucose au total pendant sa course.
    \begin{enumerate}
        \item Calculer la quantité totale d'ATP produite par le sprinteur (fermentation : 2 ATP/glucose) et par le marathonien (respiration : 30 ATP/glucose).
        \item Calculer la puissance ATP moyenne (en mol/s) pour chacun.
        \item Expliquer pourquoi le sprinteur ne peut pas maintenir son allure plus de 10-20 s alors que le marathonien tient des heures.
    \end{enumerate}
    \tcblower
    \textbf{Solution détaillée :}
    \begin{enumerate}
        \item \textbf{Quantité totale d'ATP (Rendement) :}
              \begin{itemize}
                  \item Sprinteur (Fermentation) : $n_{\text{ATP}} = 0,50 \times 2 = \mathbf{1,00\ mol\ ATP}$.
                  \item Marathonien (Respiration) : $n_{\text{ATP}} = 0,50 \times 30 = \mathbf{15,0\ mol\ ATP}$.
              \end{itemize}
        \item \textbf{Puissance ATP moyenne (Débit) :}
              \begin{itemize}
                  \item Sprinteur : $t = 10\ \text{s}$. $P = \frac{1,00}{10} = \mathbf{0,100\ mol\ ATP/s}$.
                  \item Marathonien : $t = 2 \times 3600 + 10 \times 60 = 7800\ \text{s}$. $P = \frac{15,0}{7800} \approx \mathbf{1,92 \times 10^{-3}\ mol\ ATP/s}$.
              \end{itemize}
        \item \textbf{Interprétation physiologique :}
              Le sprinteur produit de l'ATP \textbf{50 fois plus vite} (puissance maximale) grâce à la rapidité de la glycolyse (réactions cytoplasmiques, pas de transport mitochondrial, pas de limitation par l'\ce{O2}). Cependant, le rendement est très faible (1 mol ATP pour 0,5 mol glucose) et l'accumulation de lactate (et H$^+$) acidifie le milieu (pH $\downarrow$), inhibant les enzymes glycolytiques (PFK-1) et la contraction (fixation Ca$^{2+}$ sur troponine) $\to$ \textbf{fatigue rapide}.
              Le marathonien a une puissance faible mais un rendement maximal (15 mol ATP pour 0,5 mol glucose). L'approvisionnement en \ce{O2} et en substrats (lipides, glycogène, glucose sanguin) est maintenu par la ventilation, la circulation sanguine et la mobilisation des réserves $\to$ \textbf{endurance}.
    \end{enumerate}
    \textbf{Conclusion :} Fermentation = Puissance maximale, faible rendement, fatigue acide. Respiration = Rendement maximal, puissance modérée, endurance.
\end{exoresolu}

\begin{methode}[Exploiter des données expérimentales : Respirométrie (consommation \ce{O2}, production \ce{CO2})]
    \textbf{Objectif :} Interpréter des courbes de variation de volume gazeux ou de pression dans un respiromètre pour déterminer le substrat oxydé (quotient respiratoire) et le débit de consommation d'\ce{O2}.
    \begin{enumerate}
        \item \textbf{Identifier le dispositif :} Respiromètre à pression constante (volume variable, seringue/capillaire) ou pression variable (volume constant, capteur de pression). Présence de \ce{KOH} (ou soude) solide ou en solution pour absorber le \ce{CO2} produit ($\ce{CO2 + 2 KOH -> K2CO3 + H2O}$).
        \item \textbf{Analyser la courbe $V(t)$ ou $P(t)$ :}
              \begin{itemize}
                  \item \textbf{Sans \ce{KOH} :} Variation due à la consommation d'\ce{O2} ET à la production de \ce{CO2}. Si QR = 1, volumes égaux $\to$ volume/pression constants.
                  \item \textbf{Avec \ce{KOH} :} \ce{CO2} absorbé $\to$ variation due \textbf{uniquement à la consommation d'\ce{O2}} ($\Delta V \propto -\Delta n_{\ce{O2}}$). La pente de la droite donne le débit volumique $\dot{V}_{\ce{O2}}$ (en mL/min ou L/min).
              \end{itemize}
        \item \textbf{Calculer le Quotient Respiratoire (QR) :}
              \[ QR = \frac{V_{\ce{CO2}}}{V_{\ce{O2}}} = \frac{|\text{pente sans KOH}| - |\text{pente avec KOH}|}{|\text{pente avec KOH}|} \]
              (En valeur absolue ; si volume constant sans KOH, pente = 0 $\to$ QR = 1).
        \item \textbf{Interpréter le QR :} Voir tableau ci-contre (Propriété).
        \item \textbf{Passer au débit molaire :} $\dot{n}_{\ce{O2}} = \frac{\dot{V}_{\ce{O2}}}{V_m}$ avec $V_m$ volume molaire (\SI{24,0}{L/mol} à \SI{25}{\celsius} et 1 atm ; \SI{22,4}{L/mol} aux CNTP 0°C).
    \end{enumerate}
\end{methode}

\begin{exoresolu}[Détermination du substrat par respirométrie]
    \textbf{Énoncé :} On place une souris (masse \SI{30}{g}) dans un respiromètre à pression constante thermostaté à \SI{25}{\celsius}. On mesure le volume de gaz $V$ en fonction du temps $t$.
    \begin{itemize}
        \item Expérience 1 (sans \ce{KOH}) : Le volume reste constant ($V = \text{cte}$).
        \item Expérience 2 (avec \ce{KOH} dans la chambre) : Le volume diminue linéairement. La pente de la droite $V(t)$ est de $-0,12\ \text{mL/min}$.
    \end{itemize}
    \begin{enumerate}
        \item Expliquer pourquoi le volume est constant sans \ce{KOH}.
        \item Calculer le débit de consommation d'oxygène $\dot{V}_{\ce{O2}}$ (en mL/min) et $\dot{n}_{\ce{O2}}$ (en $\mu$mol/min).
        \item Déduire le Quotient Respiratoire (QR) et identifier le substrat principal oxydé.
        \item Calculer la consommation massique d'\ce{O2} (en mL/min/kg).
    \end{enumerate}
    \tcblower
    \textbf{Solution détaillée :}
    \begin{enumerate}
        \item \textbf{Interprétation Expérience 1 :} Sans \ce{KOH}, le \ce{CO2} produit reste dans la chambre gazeuse. Le volume ne change pas $\iff$ Volume \ce{O2} consommé = Volume \ce{CO2} produit. Donc $QR = 1$. Cela signe l'oxydation de \textbf{glucides} (glucose, glycogène).
        \item \textbf{Calculs Expérience 2 :} Avec \ce{KOH}, le \ce{CO2} est piégé chimiquement. La baisse de volume correspond \textbf{uniquement} à la consommation d'\ce{O2}.
              \begin{align*}
                  \dot{V}_{\ce{O2}} &= |\text{pente}| = \mathbf{0,12\ mL/min} \\
                  \dot{n}_{\ce{O2}} &= \frac{\dot{V}_{\ce{O2}}}{V_m} = \frac{0,12 \times 10^{-3}\ \text{L/min}}{24,0\ \text{L/mol}} = 5,0 \times 10^{-6}\ \text{mol/min} = \mathbf{5,0\ \mu mol/min}
              \end{align*}
        \item \textbf{QR et Substrat :} D'après Expérience 1, $QR = 1$. Le substrat principal est le \textbf{glucide} (glycogène hépatique/musculaire, glucose sanguin).
        \item \textbf{Consommation massique :} $\frac{0,12\ \text{mL/min}}{0,030\ \text{kg}} = \mathbf{4,0\ mL/min/kg}$.
    \end{enumerate}
    \textbf{Conclusion :} La souris oxyde principalement des glucides (QR=1) avec une consommation massique d'\ce{O2} de 4,0 mL/min/kg, valeur cohérente pour un petit mammifère au repos (métabolisme de base élevé, surface/volume grand).
\end{exoresolu}

\begin{propriete}[Rendement ATP détaillé par glucose (P/O = 2,5 pour NADH ; 1,5 pour FADH2)]
    Le bilan comptable exact dépend de la navette mitochondriale utilisée pour transférer les électrons des NADH cytoplasmiques (glycolyse) vers la matrice mitochondriale :
    \begin{center}
        \begin{tabular}{|l|c|c|}\hline
        \textbf{Étape / Localisation} & \textbf{Navette Malate-Aspartate (Foie, Cœur, Rein)} & \textbf{Navette Glycérol-3-P (Muscle, Cerveau)} \\ \hline
        Glycolyse (Cytoplasme) : 2 ATP (SLP) + 2 NADH$_{\text{cyt}}$ & 2 NADH$_{\text{cyt}}$ $\to$ 2 NADH$_{\text{mat}}$ (5 ATP) & 2 NADH$_{\text{cyt}}$ $\to$ 2 FADH$_{2\text{mat}}$ (3 ATP) \\ \hline
        Pyruvate $\to$ Acétyl-CoA (Matrice) : 2 NADH$_{\text{mat}}$ & 5 ATP & 5 ATP \\ \hline
        Cycle de Krebs $\times$ 2 (Matrice) : 2 GTP + 6 NADH$_{\text{mat}}$ + 2 FADH$_{2\text{mat}}$ & 2 + 15 + 3 = 20 ATP & 2 + 15 + 3 = 20 ATP \\ \hline
        \textbf{TOTAL} & \textbf{32 ATP} & \textbf{30 ATP} \\ \hline
    \end{tabular}
    \end{center}
    \textbf{Nota Bene :} Au Probatoire, on retient \textbf{30 ATP} pour le muscle squelettique (navette G3P) et \textbf{32 ATP} pour le foie/cœur (navette Malate-Aspartate). Préciser toujours l'hypothèse. L'ancien chiffre de 38 ATP (P/O = 3 et 2) est obsolète.
\end{propriete}

\begin{methode}[Calculer le rendement ATP détaillé (Glycolyse + Krebs + Chaîne respiratoire)]
    \textbf{Objectif :} Reconstituer le compte exact des ATP, NADH, FADH2 par glucose pour justifier le chiffre de 30-32 ATP.
    \begin{enumerate}
        \item \textbf{Glycolyse (Cytoplasme) :} Bilan net : 2 ATP (SLP) + 2 NADH$_{\text{cyt}}$.
              \emph{Piège :} Le transport des 2 NADH$_{\text{cyt}}$ dans la matrice mitochondriale coûte de l'énergie selon la navette.
              \begin{itemize}
                  \item Navette Malate-Aspartate (cœur, foie, rein) : 2 NADH$_{\text{cyt}}$ $\to$ 2 NADH$_{\text{mat}}$ $\to$ 2 $\times$ 2,5 = 5 ATP.
                  \item Navette Glycérol-3-P (muscle, cerveau) : 2 NADH$_{\text{cyt}}$ $\to$ 2 FADH$_{2\text{mat}}$ $\to$ 2 $\times$ 1,5 = 3 ATP.
              \end{itemize}
        \item \textbf{Pyruvate $\to$ Acétyl-CoA (Matrice) :} 2 NADH$_{\text{mat}}$ $\to$ 5 ATP.
        \item \textbf{Cycle de Krebs (Matrice) $\times$ 2 :} 2 ATP (SLP via GTP) + 6 NADH$_{\text{mat}}$ (15 ATP) + 2 FADH$_{2\text{mat}}$ (3 ATP).
        \item \textbf{Somme :}
              \begin{align*}
                  \text{SLP (Glycolyse + Krebs)} &= 4\ \text{ATP} \\
                  \text{Oxydative (Navette Malate-Aspartate)} &= 5 + 5 + 15 + 3 = 28\ \text{ATP} \to \mathbf{Total = 32\ ATP} \\
                  \text{Oxydative (Navette Glycérol-3-P)} &= 3 + 5 + 15 + 3 = 26\ \text{ATP} \to \mathbf{Total = 30\ ATP}
              \end{align*}
    \end{enumerate}
\end{methode}

\begin{exoresolu}[Bilan comptable complet dans le muscle au repos]
    \textbf{Énoncé :} Une fibre musculaire au repos oxyde du glucose via la navette glycérol-3-phosphate. Établir le bilan comptable détaillé du nombre d'ATP formés par molécule de glucose. On rappelle : 1 NADH$_{\text{mat}}$ $\to$ 2,5 ATP ; 1 FADH$_{2\text{mat}}$ $\to$ 1,5 ATP ; Phosphorylation au niveau du substrat (SLP) : 1 ATP (ou GTP) direct.
    \tcblower
    \textbf{Solution détaillée :}
    \begin{center}
        \begin{tabularx}{\linewidth}{|l|X|c|}\hline
        \textbf{Étape / Localisation} & \textbf{Produits réduits / ATP direct} & \textbf{Gain ATP} \\ \hline
        Glycolyse (Cytoplasme) & 2 ATP (SLP) + 2 NADH$_{\text{cyt}}$ & +2 \\
        Navette Glycérol-3-P & 2 NADH$_{\text{cyt}}$ $\to$ 2 FADH$_{2\text{mat}}$ & 2 $\times$ 1,5 = +3 \\
        Décarboxylation oxydative du Pyruvate (Matrice) & 2 NADH$_{\text{mat}}$ & 2 $\times$ 2,5 = +5 \\
        Cycle de Krebs $\times$ 2 (Matrice) & 2 GTP (SLP) + 6 NADH$_{\text{mat}}$ + 2 FADH$_{2\text{mat}}$ & +2 + 15 + 3 \\ \hline
        \textbf{TOTAL} & & \textbf{30 ATP} \\ \hline
        \end{tabularx}
    \end{center}
    \textbf{Justification du coût de la navette :} Dans le muscle, les électrons des NADH cytoplasmiques sont transférés au FAD mitochondrial via la glycérol-3-phosphate déshydrogénase (face externe de la membrane interne), formant du FADH$_2$ mitochondrial qui entre dans la chaîne respiratoire au niveau du complexe III (donc 1,5 ATP au lieu de 2,5 pour un NADH entrant au complexe I).
    \textbf{Conclusion :} Le rendement théorique maximal dans le muscle squelettique est de \textbf{30 ATP par glucose}.
\end{exoresolu}

\begin{definition}[VO2max, Seuil lactique, Économie de course]
    \begin{itemize}
        \item \textbf{VO2max (Consommation maximale d'O2) :} Débit maximal d'oxygène qu'un organisme peut capter, transporter et utiliser par unité de temps (mL/min/kg). Limité par le système cardio-respiratoire (débit cardiaque max, hémoglobine) et la capacité mitochondriale musculaire.
        \item \textbf{Seuil lactique (ou seuil anaérobie) :} Intensité d'effort au-delà de laquelle la production de lactate dépasse sa capacité d'élimination (oxydation, néoglucogenèse). Correspond à une élévation non linéaire de la lactatémie et souvent à QR $\approx$ 1,0 (hyperventilation compensatrice).
        \item \textbf{Économie de course (Coût énergétique) :} Consommation d'O2 ($\dot{V}_{\ce{O2}}$) nécessaire pour maintenir une vitesse ou une puissance donnée. Une économie meilleure = $\dot{V}_{\ce{O2}}$ plus faible pour la même charge externe.
    \end{itemize}
\end{definition}

\begin{methode}[Analyser l'influence de l'entraînement, l'alimentation et l'oxygénation sur le rendement]
    \textbf{Objectif :} Relier des paramètres physiologiques/environnementaux à l'efficacité de la production d'ATP (VO2max, seuils, substrats, mitochondries).
    \begin{enumerate}
        \item \textbf{Entraînement (Endurance) :}
              \begin{itemize}
                  \item $\uparrow$ Densité mitochondriale (biogenèse via PGC-1$\alpha$).
                  \item $\uparrow$ Activité enzymes oxydatives (Citrate Synthase, $\beta$-HAD, Cytochrome c oxydase).
                  \item $\uparrow$ Capillarisation $\to$ meilleure diffusion \ce{O2} (réduction distance de diffusion).
                  \item \textbf{Effet :} Déplacement du seuil lactique vers des intensités plus élevées (\% VO2max), $\uparrow$ VO2max, $\uparrow$ oxydation des lipides à iso-charge (épargne glycogène), $\uparrow$ rendement effectif (moins de fermentation pour une même charge).
              \end{itemize}
        \item \textbf{Alimentation :}
              \begin{itemize}
                  \item Glucides : QR=1, rendement ATP/\ce{O2} élevé ($\approx$ 6,3 ATP/\ce{O2}), puissance max. Indispensable pour l'intense (glycolyse rapide).
                  \item Lipides : QR=0,7, rendement ATP/\ce{O2} plus faible ($\approx$ 5,6 ATP/\ce{O2}), stock immense (adipocytes). Favori au repos/effort modéré longue durée.
                  \item \textbf{Effet :} Adapter l'apport à l'effort (charge glycémique avant, réapprovisionnement glycogène après). Jeûne prolongé $\to$ néoglucogenèse, production de corps cétoniques (carburant alternatif cerveau/cœur).
              \end{itemize}
        \item \textbf{Oxygénation (Altitude, Anémie, Pathologie) :}
              \begin{itemize}
                  \item Hypoxie (altitude, anémie, BPCO) : Limite la chaîne respiratoire (complexe IV, cytochrome c oxydase). $\to$ Bascule vers fermentation (lactate $\uparrow$), $\downarrow$ rendement ATP/glucose, fatigue précoce.
                  \item Acclimatation altitude : $\uparrow$ EPO (rein) $\to$ $\uparrow$ Hématies/Hb $\to$ $\uparrow$ transport \ce{O2} (polyglobulie).
              \end{itemize}
    \end{enumerate}
\end{methode}

\begin{exoresolu}[Effet de l'entraînement sur le seuil lactique et le substrat]
    \textbf{Énoncé :} On compare deux sujets de même masse (70 kg) sur un cyclo-ergomètre : Sujet S (sédentaire) et Sujet E (entraîné, coureur de fond). On mesure la consommation d'\ce{O2} ($\dot{V}_{\ce{O2}}$) et le quotient respiratoire (QR) en fonction de la puissance développée.
    \begin{center}
        \begin{tabular}{|c|c|c|c|c|}\hline
        Puissance (W) & $\dot{V}_{\ce{O2}}$ S (L/min) & QR S & $\dot{V}_{\ce{O2}}$ E (L/min) & QR E \\ \hline
        50 & 0,8 & 0,85 & 0,7 & 0,80 \\ \hline
        100 & 1,4 & 0,92 & 1,2 & 0,82 \\ \hline
        150 & 2,1 & 1,05 & 1,7 & 0,88 \\ \hline
        200 & 2,8 & 1,15 & 2,3 & 0,95 \\ \hline
        250 & -- & -- & 2,9 & 1,02 \\ \hline
    \end{tabular}
    \end{center}
    \begin{enumerate}
        \item Pour chaque sujet, identifier la puissance correspondant au \textbf{seuil lactique} (transition aérobie/anaérobie) à l'aide du QR.
        \item Comparer l'économie de course (coût énergétique) des deux sujets à 100 W.
        \item Expliquer pourquoi le sujet E a un QR plus bas à 150 W.
        \item Calculer la VO2max estimée du sujet E si 250 W correspond à 85\% de sa VO2max.
    \end{enumerate}
    \tcblower
    \textbf{Solution détaillée :}
    \begin{enumerate}
        \item \textbf{Seuil lactique (QR $\approx$ 1,00 / point de bascule vers QR > 1) :}
              \begin{itemize}
                  \item Sujet S : QR passe de 0,92 (100W) à 1,05 (150W). Seuil $\approx$ \textbf{120-130 W}.
                  \item Sujet E : QR passe de 0,88 (150W) à 0,95 (200W) à 1,02 (250W). Seuil $\approx$ \textbf{220-230 W}.
              \end{itemize}
              L'entraîné repousse le seuil à une puissance bien supérieure (gain \~80\%).
        \item \textbf{Économie de course à 100 W :} $\dot{V}_{\ce{O2}}$ S = 1,4 L/min ; $\dot{V}_{\ce{O2}}$ E = 1,2 L/min. Le sujet E consomme \textbf{moins d'\ce{O2} pour la même puissance externe} $\to$ meilleur rendement mécanique / couplage mitochondrial plus efficace (moins de proton leak, meilleur couplage oxydatif).
        \item \textbf{QR bas du sujet E à 150 W (0,88 vs 1,05) :} Le sujet E oxyde davantage de \textbf{lipides} (QR $\approx$ 0,7-0,85) et épargne son glycogène musculaire. Le sujet S, moins mitochondrial, est déjà en déficit d'\ce{O2} relatif au niveau musculaire et doit recourir massivement aux glucides (QR $\to$ 1) et à la fermentation lactique (hyperventilation compensatrice $\to$ QR > 1).
        \item \textbf{VO2max Sujet E :} À 250 W, $\dot{V}_{\ce{O2}}$ = 2,9 L/min = 85\% VO2max.
              \[ \text{VO2max} = \frac{2,9}{0,85} \approx \mathbf{3,41\ L/min} \quad (\text{soit } \approx 48,7\ \text{mL/min/kg}) \]
    \end{enumerate}
    \textbf{Conclusion :} L'entraînement d'endurance augmente la puissance au seuil lactique, améliore l'économie d'\ce{O2}, favorise la lipolyse (épargne glycogène) et augmente la VO2max.
\end{exoresolu}

\begin{popfigure}
\begin{tikzpicture}[scale=0.8, every node/.style={font=\small}]
    % Axes
    \draw[->] (0,0) -- (10,0) node[right] {$t$ (min)};
    \draw[->] (0,0) -- (0,5) node[above] {$\dot{V}_{\ce{O2}}$ (L/min)};
    % Repères
    \draw[dashed, popmuted] (0,1.3) node[left] {Steady-state} -- (8,1.3);
    \draw[dashed, popmuted] (0,0.25) node[left] {Repos} -- (8,0.25);
    % Phase 1: Cardiodynamique (fast rise)
    \draw[popblue, thick] (0,0.25) -- (0.3,0.8) node[midway, above, sloped] {Phase 1};
    % Phase 2: Primary exponential
    \draw[popgreen, thick, domain=0.3:4, samples=50] plot (\x, {1.3 - 1.05*exp(-(\x-0.3)*60/25) + 0.25}); % approx visual
    % Actually draw exponential properly: V(t) = Vss - A*exp(-(t-t0)/tau)
    % Let's draw a schematic curve
    \draw[poporange, thick] (0.3,0.8) .. controls (1,1.1) and (2,1.25) .. (4,1.29) node[midway, above, sloped] {Phase 2 (exponentielle)};
    % Steady state
    \draw[popdark, thick] (4,1.3) -- (6,1.3) node[midway, above] {Steady-state};
    % Recovery (EPOC)
    \draw[poppurple, thick] (6,1.3) .. controls (6.5,1.0) and (7.5,0.5) .. (8,0.3) node[midway, above, sloped] {Récupération (EPOC)};
    % Deficit area shading
    \fill[popblue!20] (0,0.25) -- (0.3,0.8) -- plot[domain=0.3:4, samples=20] (\x, {1.3 - 1.05*exp(-(\x-0.3)*60/25) + 0.25}) -- (4,1.3) -- (0,1.3) -- cycle;
    % Labels
    \node[popblue] at (1.5, 1.4) {Déficit O$_2$};
    \node[poppurple] at (7, 1.4) {Dette O$_2$ (EPOC)};
\end{tikzpicture}
\legende{Cinétique type de la consommation d'oxygène ($\dot{V}_{\ce{O2}}$) au cours d'un effort constant modéré (sous seuil) et de la récupération. Phase 1 : cardiodynamique. Phase 2 : exponentielle primaire (constante de temps $\tau$). Steady-state : palier. Zone hachurée bleue = Déficit d'O2. Zone hachurée violette (récup) = Dette d'O2 (EPOC).}
\end{popfigure}

\begin{methode}[Interpréter une courbe de cinétique de la consommation d'\ce{O2} (VO2) à l'effort]
    \textbf{Objectif :} Analyser les phases de la cinétique $\dot{V}_{\ce{O2}}(t)$ au départ d'un effort constant (phase cardiodynamique, phase exponentielle primaire, composante lente, steady-state) et en déduire le déficit/dette d'\ce{O2}.
    \begin{enumerate}
        \item \textbf{Phase 1 (0-15-20 s) : Phase cardiodynamique.} Augmentation brutale du débit cardiaque $\to$ augmentation $\dot{V}_{\ce{O2}}$ sans extraction musculaire accrue. Non représentative du métabolisme musculaire aérobie.
        \item \textbf{Phase 2 (20 s - 2-3 min) : Phase exponentielle primaire.} $\dot{V}_{\ce{O2}}$ monte exponentiellement vers un palier (steady-state) pour les efforts \textbf{modérés} (sous seuil lactique). Modèle : $\dot{V}_{\ce{O2}}(t) = \dot{V}_{\ce{O2}}(\text{base}) + \Delta \dot{V}_{\ce{O2}} (1 - e^{-(t-t_0)/\tau})$. Constante de temps $\tau \approx$ 20-30 s (entraîné) à 40-60 s (sédentaire).
        \item \textbf{Phase 3 (Efforts intenses/sévères > Seuil critique) : Composante lente.} $\dot{V}_{\ce{O2}}$ continue d'augmenter lentement après la phase primaire, ne stabilisant pas (ou atteignant VO2max). Signe de recrutement fibres rapides (type IIx), fatigue, coût \ce{O2} supplémentaire (pompe Na$^+$/K$^+$, instabilité ionique).
        \item \textbf{Déficit d'\ce{O2} (Début effort) :} Surface entre la demande d'\ce{O2} théorique (instantanée = steady-state) et la $\dot{V}_{\ce{O2}}$ réelle (intégrale de la différence). Reflète la part anaérobie (phosphagènes PCr + Glycolyse) au démarrage. Unité : mL ou L.
        \item \textbf{Dette d'\ce{O2} / EPOC (Récupération) :} $\dot{V}_{\ce{O2}}$ reste élevée après l'arrêt. Surface entre $\dot{V}_{\ce{O2}}$ récupération et $\dot{V}_{\ce{O2}}$ repos. Sert à : resynthétiser PCr (majoritaire), oxyder lactate (Corrie cycle, oxydation in situ), restaurer \ce{O2} sang (Hb) et myoglobine (Mb), thermorégulation, ventilation élevée, pompe Na$^+$/K$^+$. \textbf{La dette est toujours supérieure au déficit.}
    \end{enumerate}
\end{methode}

\begin{exoresolu}[Analyse de la cinétique VO2 et déficit/dette d'oxygène]
    \textbf{Énoncé :} Un sujet effectue un exercice constant à 100 W (intensité modérée, sous seuil). Sa $\dot{V}_{\ce{O2}}$ au repos est de 0,25 L/min. La cinétique de la phase primaire suit : $\dot{V}_{\ce{O2}}(t) = 0,25 + 1,05 (1 - e^{-t/25})$ ($t$ en secondes, $\dot{V}_{\ce{O2}}$ en L/min). L'exercice dure 6 min.
    \begin{enumerate}
        \item Calculer la $\dot{V}_{\ce{O2}}$ de steady-state (palier).
        \item Calculer le déficit d'\ce{O2} (en mL) pendant les 3 premières minutes (180 s).
        \item L'exercice s'arrête à 6 min. La récupération suit : $\dot{V}_{\ce{O2}}(t_{\text{rec}}) = 0,25 + 1,05 e^{-t_{\text{rec}}/40}$. Calculer la dette d'\ce{O2} (EPOC) sur les 5 premières minutes de récupération (300 s).
        \item Comparer déficit et dette. Commenter.
    \end{enumerate}
    \tcblower
    \textbf{Solution détaillée :}
    \begin{enumerate}
        \item \textbf{Steady-state :} Pour $t \to \infty$, $e^{-t/25} \to 0$.
              $\dot{V}_{\ce{O2}}(\infty) = 0,25 + 1,05 = \mathbf{1,30\ L/min}$.
        \item \textbf{Déficit d'\ce{O2} (0 à 180 s) :} C'est l'intégrale de [Demande - Réelle]. La demande est instantanée à 1,30 L/min.
              \emph{Attention aux unités :} $\dot{V}_{\ce{O2}}$ en L/min, $t$ en secondes. Il faut convertir $\dot{V}_{\ce{O2}}$ en mL/s ou $dt$ en min.
              $\Delta \dot{V}_{\ce{O2}} = 1,05\ \text{L/min} = 17,5\ \text{mL/s}$.
              $\dot{V}_{\ce{O2}}(t) = 4,167 + 17,5 (1 - e^{-t/25})\ \text{mL/s}$.
              Demande = $21,667\ \text{mL/s}$.
              \[ \text{Déficit} = \int_0^{180} \left[ 21,667 - \left(4,167 + 17,5(1-e^{-t/25})\right) \right] dt = \int_0^{180} 17,5 e^{-t/25} dt \]
              \[ = 17,5 \times \left[ -25 e^{-t/25} \right]_0^{180} = 437,5 \times (1 - e^{-7,2}) \approx \mathbf{437\ mL} \]
              (L'intégrale de 0 à $\infty$ donne exactement $A \times \tau = 17,5 \times 25 = 437,5$ mL. À 3 min, le déficit est quasi total).
        \item \textbf{Dette d'\ce{O2} / EPOC (0 à 300 s rec) :} Intégrale de l'excès par rapport au repos (4,167 mL/s).
              Excès = $17,5 e^{-t/40}\ \text{mL/s}$.
              \[ \text{EPOC} = \int_0^{300} 17,5 e^{-t/40} dt = 17,5 \times \left[ -40 e^{-t/40} \right]_0^{300} = 700 \times (1 - e^{-7,5}) \approx \mathbf{700\ mL} \]
        \item \textbf{Comparaison :} Dette (700 mL) > Déficit (437 mL).
              \textbf{Explication :} Le déficit ne couvre que le coût de la resynthèse des phosphagènes (PCr $\approx$ 70\%) et de la glycolyse initiale (lactate $\approx$ 30\%). La dette inclut \textbf{en plus} : l'oxydation du lactate produit (même modéré), la restauration des réserves d'\ce{O2} (hémoglobine, myoglobine), le coût énergétique de la ventilation et de la thermorégulation maintenues élevées, et le coût de la pompe Na$^+$/K$^+$ pour restaurer les gradients ioniques perturbés par l'effort.
    \end{enumerate}
    \textbf{Conclusion :} La dette d'\ce{O2} (EPOC) est systématiquement supérieure au déficit d'\ce{O2} car la récupération implique des processus métaboliques plus larges que la simple compensation du démarrage anaérobie.
\end{exoresolu}

\begin{methode}[Résoudre un problème de génétique mitochondriale (Hérédité maternelle)]
    \textbf{Objectif :} Résoudre un exercice de transmission d'une maladie mitochondriale (ex: MELAS, LHON, Myopathie, Surdité A1555G) en appliquant la règle de l'hérédité maternelle stricte.
    \begin{enumerate}
        \item \textbf{Rappeler la règle fondamentale :} Les mitochondries (et leur ADNmt circulaire) sont apportées quasi exclusivement par l'ovocyte (cytoplasme volumineux). Le spermatozoïde n'apporte que son noyau haploïde et son centriole ; ses mitochondries (flagelle) sont ubiquitinées et détruites dans le zygote. $\to$ \textbf{Seules les mères transmettent l'ADNmt à TOUS leurs enfants (garçons et filles). Les pères ne transmettent jamais leur ADNmt.}
        \item \textbf{Analyser l'arbre généalogique :}
              \begin{itemize}
                  \item Femme atteinte $\to$ Tous ses enfants (M et F) héritent de l'ADNmt mutant (homoplasie ou hétéroplasie).
                  \item Homme atteint $\to$ \textbf{Aucun} de ses enfants n'hérite de la mutation.
                  \item Femme saine $\to$ Enfants sains (sauf mutation \emph{de novo} rare).
              \end{itemize}
        \item \textbf{Gérer l'hétéroplasie :} Mélange ADNmt mutant / sauvage dans une cellule. Seuil de pathogénicité (ex: 60-80\% mutants) pour expression phénotypique. Distribution aléatoire lors de la mitose (bouteille génétique oocytaire) $\to$ variabilité intra-familiale (pénétrance incomplète, expressivité variable).
        \item \textbf{Ne pas confondre :} Maladie nucléaire (autosomique récessive/dominante ou liée X) touchant une protéine mitochondriale codée par l'ADN nucléaire (ex: Frataxine dans l'ataxie de Friedreich, POLG) $\to$ Hérédité mendélienne classique.
    \end{enumerate}
\end{methode}

\begin{exoresolu}[Arbre généalogique : Surdité mitochondriale (mutation A1555G)]
    \textbf{Énoncé :} La mutation A1555G du gène \emph{12S rRNA} (ADNmt) prédispose à la surdité neurosensorielle, surtout après prise d'aminoglycosides. On étudie la famille ci-dessous :
    \begin{itemize}
        \item Génération I : Grand-mère (I-1) atteinte. Grand-père (I-2) sain.
        \item Génération II : 3 filles (II-1, II-2, II-3) et 1 fils (II-4). II-1 et II-3 atteintes. II-2 et II-4 sains.
        \item Génération III :
              \begin{itemize}
                  \item II-1 (atteinte) $\times$ Homme sain $\to$ 1 fils (III-1) atteint, 1 fille (III-2) atteinte.
                  \item II-2 (saine) $\times$ Homme sain $\to$ 1 fille (III-3) saine.
                  \item II-3 (atteinte) $\times$ Homme sain $\to$ 1 fils (III-4) sain (!), 1 fille (III-5) atteinte.
                  \item II-4 (sain) $\times$ Femme saine $\to$ 1 fils (III-6) sain.
              \end{itemize}
    \end{itemize}
    \begin{enumerate}
        \item Ce mode de transmission est-il compatible avec une maladie mitochondriale ? Justifier.
        \item Expliquer le cas de III-4 (fils d'une mère atteinte, mais sain).
        \item Calculer la probabilité qu'un enfant de III-2 (atteinte) soit atteint.
        \item Pourquoi le grand-père (I-2) n'a-t-il pas transmis la maladie ?
    \end{enumerate}
    \tcblower
    \textbf{Solution détaillée :}
    \begin{enumerate}
        \item \textbf{Compatibilité :} \textbf{OUI, parfaitement.}
              \begin{itemize}
                  \item I-1 (femme atteinte) transmet à \textbf{tous} ses enfants (II-1, II-2, II-3, II-4) $\to$ Cohérent (hétéroplasie possible expliquant II-2 et II-4 sains).
                  \item II-1 (femme atteinte) transmet à \textbf{tous} ses enfants (III-1, III-2) $\to$ Cohérent.
                  \item II-3 (femme atteinte) transmet à sa fille (III-5) mais pas à son fils (III-4) $\to$ Possible par hétéroplasie (voir q2).
                  \item II-4 (homme sain, fils d'atteinte) $\times$ Femme saine $\to$ Enfant (III-6) sain. \textbf{Preuve clé :} Un homme ne transmet pas l'ADNmt. Si c'était autosomique dominant, II-4 aurait 50\% de chance de transmettre l'allèle mutant. Ici 0\% de transmission $\to$ Exclut l'autosomique dominant. Compatible mitochondrial.
                  \item II-2 (femme saine) $\to$ Enfant sain. Cohérent.
              \end{itemize}
        \item \textbf{Cas III-4 :} Phénomène d'\textbf{hétéroplasie}. La mère II-3 porte un mélange d'ADNmt mutant et sauvage. Lors de l'oogénèse (bouteille génétique mitotique), l'échantillon d'ADNmt transmis à l'ovocyte fécondé donnant III-4 a eu une proportion de mutants \textbf{inférieure au seuil pathogène} (ex: < 60\%), alors que l'échantillon pour III-5 était au-dessus. Distribution stochastique (aléatoire) lors des divisions cellulaires.
        \item \textbf{Probabilité pour enfant de III-2 :} III-2 est atteinte (phénotype exprimé). Elle a donc une charge mutante élevée (probablement > seuil dans la plupart des tissus, y compris la lignée germinale). Probabilité de transmission à un enfant $\approx$ \textbf{1 (ou 100\%)} (quasiment tous les ovocytes auront une charge > seuil). \emph{Nuance :} Si hétéroplasie résiduelle faible, probabilité < 1 mais très élevée.
        \item \textbf{Grand-père I-2 :} Les mitochondries du spermatozoïde sont situées dans la pièce intermédiaire du flagelle et sont \textbf{activement détruites/éliminées} après la fécondation (marquage par ubiquitine, dégradation par le protéasome). Seules les mitochondries de l'ovocyte (maternelles) persistent dans le zygote. Donc I-2 ne peut transmettre son ADNmt, qu'il soit mutant ou non.
    \end{enumerate}
    \textbf{Conclusion :} L'arbre confirme une transmission maternelle stricte avec hétéroplasie expliquant la non-pénétrance chez III-4 et II-2/II-4.
\end{exoresolu}

\begin{attention}[Les 5 erreurs qui coûtent des points au Probatoire]
    \begin{enumerate}
        \item \textbf{Confondre Rendement et Puissance :} Écrire « La fermentation produit plus d'ATP que la respiration car elle est plus rapide ». \textbf{FAUX.} La fermentation a une \textbf{puissance} (débit ATP/s) plus élevée mais un \textbf{rendement} (quantité totale ATP/mol glucose) 15 fois plus faible (2 vs 30 ATP). \emph{Astuce : Préciser toujours l'unité : ATP/mol glucose (rendement) vs ATP/s (puissance).}
        \item \textbf{Oublier le coût des navettes mitochondriales :} Écrire « 32 ATP » pour le muscle ou « 30 ATP » pour le foie sans justification, ou pire, écrire « 38 ATP » (vieux chiffre théorique P/O = 3 et 2). \textbf{Au Probatoire :} On utilise P/O = 2,5 (NADH) et 1,5 (FADH2). \textbf{Muscle = 30 ATP (navette G3P), Foie/Cœur = 32 ATP (navette Malate-Aspartate).} Il faut le préciser explicitement.
        \item \textbf{Erreur sur le Quotient Respiratoire (QR) :} Calculer $QR = \frac{\ce{O2}}{\ce{CO2}}$ au lieu de $\frac{\ce{CO2}}{\ce{O2}}$. Ou affirmer « QR > 1 signifie oxydation de lipides ». \textbf{FAUX.} QR > 1 = Hyperventilation / Lipogenèse / Effort intense (tamponnage lactate $\to$ \ce{CO2} excédentaire). Lipides $\to$ QR $\approx$ 0,7. Glucides $\to$ QR = 1,0.
        \item \textbf{Confondre Déficit et Dette d'Oxygène :} Dire « La dette d'O2 sert à payer le déficit ». \textbf{Imprécis.} La dette (EPOC) est \textbf{toujours supérieure} au déficit. Le déficit = anaérobie de démarrage (PCr + Glycolyse). La dette = Récupération PCr + Oxydation Lactate + Restauration \ce{O2} sang/myoglobine + Thermorégulation + Ventilation + Pompe Na$^+$/K$^+$. Il faut lister les composantes de la dette.
        \item \textbf{Hérédité mitochondriale : « Le père transmet à ses fils » ou « Transmission liée au sexe ».} \textbf{FAUX.} Transmission \textbf{strictement maternelle} (cytoplasmique). Tous les enfants d'une mère atteinte héritent de l'ADNmt (hétéroplasie). Aucun enfant d'un père atteint n'hérite de l'ADNmt. Ne pas confondre avec l'hérédité liée au chromosome X (homme transmet à toutes ses filles).
    \end{enumerate}
\end{attention}

\begin{savaistu}[Le lactate n'est pas un déchet, mais un carburant !]
    Contrairement à l'idée reçue (« acide lactique = fatigue »), le lactate (ion lactate$^-$ + H$^+$) est un \textbf{métabolite énergétique majeur}. Produit par les fibres rapides (glycolyse rapide), il diffuse dans le sang et est :
    \begin{itemize}
        \item Oxydé sur place par les fibres lentes voisines (navette lactate, oxydation mitochondriale).
        \item Transporté vers le foie pour la néoglucogenèse (Cycle de Cori : Lactate $\to$ Glucose $\to$ Sang $\to$ Muscle).
        \item Utilisé par le cœur (carburant préféré au repos).
    \end{itemize}
    L'entraînement d'endurance augmente l'expression des transporteurs MCT1 (import lactate) et de la LDH-1 (isoforme oxydative), transformant le muscle en « éponge à lactate ». C'est pour cela que le seuil lactique recule : on élimine le lactate aussi vite qu'on le produit.
\end{savaistu}

\newpage

\coursec{Exercices}

\coursec{Leçon NCT01 : Amélioration de la production de l'énergie par l'organisme}

\courssub{Je vérifie mes connaissances}

\begin{exercice}[QCM : la respiration cellulaire]
Pour chaque question, une seule réponse est exacte. Entoure la lettre correspondante.

\textbf{1.} La glycolyse se déroule :
\begin{itemize}
\item[a)] dans la matrice mitochondriale ;
\item[b)] dans le hyaloplasme (cytosol) ;
\item[c)] dans la crête mitochondriale ;
\item[d)] dans le noyau.
\end{itemize}

\textbf{2.} L'accepteur final des électrons de la chaîne respiratoire est :
\begin{itemize}
\item[a)] le \ce{CO2} ;
\item[b)] le glucose ;
\item[c)] le dioxygène \ce{O2} ;
\item[d)] l'acide pyruvique.
\end{itemize}

\textbf{3.} Le bilan énergétique de la respiration cellulaire d'une molécule de glucose est d'environ :
\begin{itemize}
\item[a)] 2 ATP ;
\item[b)] 12 ATP ;
\item[c)] 36 à 38 ATP ;
\item[d)] 100 ATP.
\end{itemize}

\textbf{4.} La fermentation lactique produit :
\begin{itemize}
\item[a)] de l'éthanol et du \ce{CO2} ;
\item[b)] de l'acide lactique ;
\item[c)] de l'eau et du \ce{CO2} ;
\item[d)] de l'acide pyruvique.
\end{itemize}
\end{exercice}

\begin{exercice}[Vrai ou faux justifié]
Pour chacune des affirmations suivantes, indique si elle est \textbf{vraie} ou \textbf{fausse}, puis justifie ta réponse en une ou deux phrases.

\begin{enumerate}
\item « La glycolyse nécessite la présence de dioxygène. »
\item « La fermentation produit autant d'ATP que la respiration cellulaire. »
\item « Le cycle de Krebs se déroule dans la matrice des mitochondries. »
\item « Une cellule privée de mitochondries ne peut produire aucun ATP. »
\end{enumerate}
\end{exercice}

\begin{exercice}[Définitions]
Définis en une ou deux phrases chacun des termes suivants :
\begin{enumerate}
\item respiration cellulaire ;
\item fermentation ;
\item quotient respiratoire (QR) ;
\item chaîne respiratoire.
\end{enumerate}
\end{exercice}

\begin{exercice}[Calcul direct : bilan énergétique]
Une cellule musculaire d'un coureur dégrade complètement, par respiration cellulaire, \num{0,5}~mol de glucose.

\textbf{Données :} la respiration d'une molécule de glucose produit environ 36 ATP ; la masse molaire du glucose est $M(\ce{C6H12O6}) = \SI{180}{g.mol^{-1}}$.

\begin{enumerate}
\item Écris le bilan simplifié et équilibré de la respiration cellulaire.
\item Calcule le nombre de mol d'ATP produites.
\item Calcule la masse de glucose dégradée.
\item Calcule le volume de dioxygène consommé, sachant que dans les conditions de l'expérience le volume molaire gazeux vaut $V_m = \SI{24}{L.mol^{-1}}$.
\end{enumerate}
\end{exercice}

\courssub{Je m'entraîne}

\begin{exercice}[Comparaison des rendements énergétiques]
On souhaite comparer les deux voies de production d'ATP par la cellule.

\textbf{Données :} respiration : 36 ATP par molécule de glucose ; fermentation : 2 ATP par molécule de glucose ; $M(\text{glucose}) = \SI{180}{g.mol^{-1}}$.

\begin{enumerate}
\item Écris le bilan simplifié et équilibré de la respiration cellulaire.
\item Calcule la quantité de matière (en mol) puis la masse de glucose nécessaire pour produire 180~mol d'ATP par respiration.
\item Calcule la quantité de matière puis la masse de glucose nécessaire pour produire 180~mol d'ATP par fermentation.
\item Calcule le rapport des masses de glucose consommées (fermentation / respiration). Conclus sur le rendement énergétique comparé des deux voies.
\end{enumerate}
\end{exercice}

\begin{exercice}[Quotient respiratoire de graines germées]
Au laboratoire du lycée, on place des graines de haricot germées dans un appareil permettant de mesurer les échanges gazeux. On relève chaque heure les volumes cumulés de \ce{O2} consommé et de \ce{CO2} produit :

\begin{center}
\begin{tabularx}{\linewidth}{|l|*{6}{>{\centering\arraybackslash}X|}}
\hline
\rowcolor{popblueL} Durée (h) & 1 & 2 & 3 & 4 & 5 & 6 \\
\hline
\ce{O2} consommé (mL) & 12 & 24 & 36 & 48 & 60 & 72 \\
\hline
\ce{CO2} produit (mL) & 12 & 24 & 36 & 48 & 60 & 72 \\
\hline
\end{tabularx}
\end{center}

\begin{enumerate}
\item Rappelle la formule du quotient respiratoire (QR).
\item Calcule le QR à chaque heure de mesure.
\item Trace la courbe $\mathrm{QR} = f(\text{temps})$ sur papier millimétré (échelles : \SI{2}{cm} pour 1~h en abscisse ; \SI{4}{cm} pour 1 unité de QR en ordonnée).
\item Interprète la valeur du QR obtenue en t'aidant du bilan de la respiration cellulaire. Que respirent les graines germées ?
\end{enumerate}
\end{exercice}

\begin{exercice}[Légende d'un schéma de mitochondrie]
Le schéma ci-dessous représente une mitochondrie en coupe.

\begin{popfigure}
\begin{tikzpicture}[scale=1.0]
% membrane externe
\draw[thick, popblue] (0,0) ellipse (4.2 and 1.9);
% membrane interne avec crêtes
\draw[thick, poporange, decorate, decoration={coil, segment length=6pt, amplitude=3pt}] (0,0) ellipse (3.6 and 1.4);
% matrice
\fill[poporange!15] (0,0) ellipse (3.6 and 1.4);
% espace intermembranaire
\fill[popblue!10, even odd rule] (0,0) ellipse (4.2 and 1.9) (0,0) ellipse (3.6 and 1.4);
% labels
\draw[thin, popmuted] (3.2,1.55) -- (5.0,2.3) node[right, align=left] {\textbf{1}};
\draw[thin, popmuted] (2.6,1.05) -- (5.0,1.3) node[right, align=left] {\textbf{2}};
\draw[thin, popmuted] (0,-0.6) -- (5.0,-0.6) node[right, align=left] {\textbf{3}};
\draw[thin, popmuted] (-3.0,1.2) -- (-5.0,1.8) node[left, align=right] {\textbf{4}};
\end{tikzpicture}
\legende{Schéma simplifié d'une mitochondrie en coupe. Les pointés 1 à 4 désignent : la membrane externe, l'espace intermembranaire, la matrice et une crête mitochondriale (dans le désordre).}
\end{popfigure}

\begin{enumerate}
\item Identifie les structures désignées par les numéros 1 à 4.
\item Recopie le tableau ci-dessous et place chaque étape de la respiration cellulaire dans le compartiment correct (certains compartiments n'accueillent aucune étape) :
\begin{center}
\begin{tabularx}{\linewidth}{|X|X|}
\hline
\rowcolor{popblueL} Étape & Compartiment (numéro et nom) \\
\hline
Glycolyse & \\
\hline
Cycle de Krebs & \\
\hline
Chaîne respiratoire & \\
\hline
\end{tabularx}
\end{center}
\item Explique pourquoi la glycolyse ne peut pas être placée sur ce schéma.
\end{enumerate}
\end{exercice}

\begin{exercice}[Tableau comparatif respiration / fermentation]
\begin{enumerate}
\item Construis un tableau comparatif de la respiration cellulaire et de la fermentation, comportant les critères suivants : localisation dans la cellule ; besoin en dioxygène ; produits obtenus ; nombre d'ATP produits par molécule de glucose.
\item Rédige un paragraphe de 8 à 10 lignes expliquant l'intérêt de la respiration cellulaire pour un organisme pluricellulaire comme l'Homme, et pourquoi la fermentation ne peut être qu'une solution temporaire.
\end{enumerate}
\end{exercice}

\courssub{Je me prépare au Probatoire}

\begin{exercice}[L'acide lactique du sprinteur — étude de cas]
Lors d'une compétition inter-lycées à Yaoundé, un athlète de Première C court un 400~m en 52~s. On mesure son taux d'acide lactique sanguin avant l'effort, à l'arrivée, puis pendant la récupération :

\begin{center}
\begin{tabularx}{\linewidth}{|l|*{5}{>{\centering\arraybackslash}X|}}
\hline
\rowcolor{popblueL} Temps & avant course & arrivée & 5 min après & 15 min après & 30 min après \\
\hline
Acide lactique sanguin (mmol/L) & 1,0 & 12,5 & 9,0 & 4,5 & 1,2 \\
\hline
\end{tabularx}
\end{center}

\begin{enumerate}
\item Rappelle dans quel compartiment cellulaire et par quel processus l'acide lactique est produit dans les muscles pendant l'effort intense.
\item Explique pourquoi, malgré la présence de mitochondries dans ses fibres musculaires, l'athlète produit de l'acide lactique pendant le sprint.
\item Trace la courbe donnant le taux d'acide lactique en fonction du temps (échelles : \SI{1}{cm} pour 5~min ; \SI{1}{cm} pour 2~mmol/L).
\item Explique pourquoi le taux d'acide lactique diminue après l'effort, en précisant le devenir de l'acide lactique et le rôle du dioxygène consommé pendant la récupération.
\item Propose deux moyens, fondés sur les connaissances du cours, permettant à cet athlète d'améliorer le rendement énergétique de ses muscles.
\end{enumerate}
\end{exercice}

\begin{exercice}[Expérience sur les échanges gazeux de la souris]
Dans un centre de recherche, on place une souris dans une enceinte close reliée à des capteurs. L'expérience est conduite à température constante. On obtient les résultats suivants :

\begin{center}
\begin{tabularx}{\linewidth}{|l|*{4}{>{\centering\arraybackslash}X|}}
\hline
\rowcolor{popblueL} Durée (min) & 0 & 10 & 20 & 30 \\
\hline
\ce{O2} dans l'enceinte (\%) & 21,0 & 19,5 & 18,0 & 16,5 \\
\hline
\ce{CO2} dans l'enceinte (\%) & 0,03 & 1,2 & 2,4 & 3,6 \\
\hline
\end{tabularx}
\end{center}

Dans une seconde expérience, on répète le protocole avec une enceinte dont la température est abaissée de \SI{10}{\degreeCelsius} : la consommation de \ce{O2} de la souris augmente nettement.

\begin{enumerate}
\item Analyse les variations des teneurs en \ce{O2} et en \ce{CO2} au cours de la première expérience.
\item Calcule le quotient respiratoire de la souris entre 0 et 30 min. Que peut-on en déduire sur la nature des nutriments majoritairement utilisés ?
\item Nomme l'organite cellulaire responsable de la consommation de \ce{O2} et précise le rôle exact du dioxygène dans la respiration cellulaire.
\item Explique, à l'aide de tes connaissances, pourquoi la consommation de \ce{O2} augmente lorsque la température de l'enceinte diminue.
\item Prédis ce qui se passerait pour la souris si la teneur en \ce{O2} de l'enceinte devenait très faible : quelle voie métabolique serait sollicitée et avec quelles conséquences ?
\end{enumerate}
\end{exercice}

\begin{exercice}[Alimentation et énergie chez le sportif camerounais]
Un préparateur physique conseille à des footballeurs de consommer, quelques heures avant un match, un repas riche en glucides (foufou de maïs, riz, igname). On donne les informations suivantes :

\begin{itemize}
\item la respiration complète d'une molécule de glucose produit environ 36 ATP ;
\item la masse molaire du glucose vaut \SI{180}{g.mol^{-1}} ;
\item un match de football de 90 min nécessite environ \num{1200}~mol d'ATP pour l'ensemble des muscles sollicités ;
\item le volume molaire gazeux dans les conditions de l'effort vaut $V_m = \SI{25}{L.mol^{-1}}$.
\end{itemize}

\begin{enumerate}
\item Écris le bilan simplifié et équilibré de la respiration cellulaire, puis identifie sur ce bilan le nutriment énergétique et le comburant.
\item Calcule la quantité de matière puis la masse de glucose que les muscles doivent dégrader par respiration pour fournir les \num{1200}~mol d'ATP du match.
\item Calcule le volume de dioxygène correspondant consommé par les muscles pendant le match.
\item Explique pourquoi le préparateur recommande un repas riche en glucides plutôt qu'en lipides juste avant le match, en t'appuyant sur la notion de rendement énergétique.
\item En fin de match, les joueurs ressentent des courbatures. Un coéquipier affirme : « Les courbatures prouvent que nos muscles ont respiré sans oxygène. » Corrige cette affirmation en distinguant clairement respiration et fermentation, et explique l'origine réelle des douleurs musculaires.
\end{enumerate}
\end{exercice}



\newpage

\coursec{Corrigés des exercices}

\begin{corrige}[Exercice 1 — QCM : la respiration cellulaire]
\textbf{1. Réponse b).} La glycolyse se déroule dans le \cle{cytosol} (hyaloplasme) : c'est une voie commune à toutes les cellules, y compris celles dépourvues de mitochondries (procaryotes, hématies). Les réponses a) et c) correspondent aux étapes mitochondriales (cycle de Krebs et chaîne respiratoire), et d) est le siège du matériel génétique.

\textbf{2. Réponse c).} Le dioxygène \ce{O2} est l'\cle{accepteur final des électrons} de la chaîne respiratoire : il capte les électrons et les protons pour former de l'eau selon $\ce{O2 + 4 H+ + 4 e- -> 2 H2O}$. Sans \ce{O2}, la chaîne s'arrête et la production d'ATP mitochondriale cesse.

\textbf{3. Réponse c).} Le bilan théorique est d'environ \textbf{36 à 38 ATP} par molécule de glucose (selon les coefficients P/O classiques : 3 ATP/NADH, 2 ATP/FADH2) : 2 ATP nets de la glycolyse, 2 ATP du cycle de Krebs (via GTP), et environ 32 à 34 ATP issus de la réoxydation des coenzymes réduits (\ce{NADH}, \ce{FADH2}) au niveau de la chaîne respiratoire.

\textbf{4. Réponse b).} La \cle{fermentation lactique} (muscles en effort intense, bactéries lactiques) réduit le \textbf{pyruvate} (acide pyruvique) en \textbf{lactate} (acide lactique). La réponse a) décrit la fermentation alcoolique des levures ; c) décrit la respiration ; d) est le produit de la glycolyse, pas de la fermentation.

\textbf{Point de vigilance :} ne pas confondre fermentation lactique et fermentation alcoolique — au Probatoire, préciser toujours laquelle est en jeu.
\end{corrige}

\begin{corrige}[Exercice 2 — Vrai ou faux justifié]
\begin{enumerate}
\item \textbf{FAUX.} La glycolyse est une voie \cle{anaérobie} : elle se déroule dans le cytosol sans aucune intervention du dioxygène. C'est d'ailleurs pour cela qu'elle peut fonctionner chez les organismes vivant sans \ce{O2}.
\item \textbf{FAUX.} La fermentation ne produit que \textbf{2 ATP} par glucose (ceux de la glycolyse seule), contre environ \textbf{36 ATP} pour la respiration complète. Le rendement de la respiration est donc environ \textbf{18 fois supérieur}.
\item \textbf{VRAI.} Les enzymes du \cle{cycle de Krebs} sont localisées dans la \textbf{matrice mitochondriale} (compartiment délimité par la membrane interne). Seule la succinate déshydrogénase est insérée dans la membrane interne.
\item \textbf{FAUX.} Une cellule sans mitochondries peut encore produire de l'ATP \emph{par glycolyse} (2 ATP/glucose). Exemple : l'\cle{hématie} (globule rouge) humaine, dépourvue de mitochondries, tire tout son ATP de la glycolyse suivie de fermentation lactique.
\end{enumerate}
\textbf{Point de vigilance :} une justification de type « vrai/faux » doit toujours citer le lieu cellulaire et/ou un chiffre de rendement ; une réponse sans justification ne vaut aucun point.
\end{corrige}

\begin{corrige}[Exercice 3 — Définitions]
\begin{enumerate}
\item \textbf{Respiration cellulaire :} ensemble des réactions d'oxydation \emph{aérobies} par lesquelles la cellule dégrade complètement un substrat organique (glucose) en \ce{CO2} et \ce{H2O}, en récupérant l'énergie libérée sous forme d'ATP. Bilan : $\ce{C6H12O6 + 6 O2 -> 6 CO2 + 6 H2O} + \text{énergie (ATP)}$.
\item \textbf{Fermentation :} voie métabolique \emph{anaérobie} qui prolonge la glycolyse : le pyruvate (acide pyruvique) est réduit (en lactate / acide lactique ou en éthanol + \ce{CO2}) afin de régénérer le \ce{NAD+} indispensable à la poursuite de la glycolyse. Elle ne produit aucun ATP supplémentaire (bilan : 2 ATP/glucose).
\item \textbf{Quotient respiratoire (QR) :} rapport du volume de \ce{CO2} produit au volume de \ce{O2} consommé pendant le même temps : $\mathrm{QR} = \dfrac{V_{\ce{CO2}}}{V_{\ce{O2}}}$. Il renseigne sur la nature du substrat oxydé (QR $\approx 1$ pour les glucides, $\approx \num{0,7}$ pour les lipides).
\item \textbf{Chaîne respiratoire :} ensemble de transporteurs d'électrons (complexes protéiques et cytochromes) insérés dans la \textbf{membrane interne mitochondriale}, qui transfèrent les électrons des coenzymes réduits (\ce{NADH}, \ce{FADH2}) jusqu'au dioxygène ; l'énergie libérée sert à créer un gradient de protons utilisé par l'ATP synthase pour fabriquer l'ATP.
\end{enumerate}
\textbf{Point de vigilance :} dans une définition, toujours préciser le \emph{lieu} (compartiment), la \emph{condition} (aérobie/anaérobie) et le \emph{résultat} (ATP, produits).
\end{corrige}

\begin{corrige}[Exercice 4 — Calcul direct : bilan énergétique]
\begin{enumerate}
\item Bilan simplifié et équilibré :
\[ \ce{C6H12O6 + 6 O2 -> 6 CO2 + 6 H2O} + \text{énergie } (\approx 36\ \text{ATP}) \]
Vérification de l'équilibre : 6 C, 12 H et 18 O de chaque côté.
\item Chaque mole de glucose produit 36 mol d'ATP :
\[ n_{\mathrm{ATP}} = \num{0,5} \times 36 = \mathbf{18\ mol\ d'ATP}. \]
\item Masse de glucose dégradée :
\[ m = n \times M = \num{0,5} \times 180 = \mathbf{90\ g}. \]
\item D'après le bilan, 1 mol de glucose consomme 6 mol de \ce{O2} :
\[ n_{\ce{O2}} = \num{0,5} \times 6 = 3\ \text{mol} \quad\Longrightarrow\quad V_{\ce{O2}} = n \times V_m = 3 \times 24 = \mathbf{72\ L} \quad (\text{à CNTP, } V_m = \num{24}\ \text{L/mol}). \]
\end{enumerate}
\textbf{Point de vigilance :} toujours écrire la relation littérale avant l'application numérique, et ne pas oublier les unités (mol, g, L).
\end{corrige}

\begin{corrige}[Exercice 5 — Comparaison des rendements énergétiques]
\begin{enumerate}
\item Bilan de la respiration : $\ce{C6H12O6 + 6 O2 -> 6 CO2 + 6 H2O} + \approx 36\ \text{ATP}$.
\item \textbf{Par respiration :} pour produire \num{180} mol d'ATP (quantité de référence), il faut $\dfrac{\num{180}}{36} = 5$ mol de glucose, soit
\[ m = 5 \times 180 = \mathbf{900\ g\ de\ glucose}. \]
\item \textbf{Par fermentation :} pour produire la même quantité (\num{180} mol d'ATP), il faut $\dfrac{\num{180}}{2} = 90$ mol de glucose, soit
\[ m = 90 \times 180 = \num{16200}\ \text{g} = \mathbf{16,2\ kg\ de\ glucose}. \]
\item Rapport des masses : $\dfrac{\num{16200}}{900} = \mathbf{18}$.
\end{enumerate}
\textbf{Conclusion :} pour fournir la même quantité d'ATP, la fermentation consomme \textbf{18 fois plus de glucose} que la respiration. Son rendement énergétique est très faible : elle ne peut être qu'une solution d'urgence et de courte durée, car elle épuiserait rapidement les réserves glucidiques de l'organisme.

\textbf{Point de vigilance :} la conclusion doit être rédigée en phrases et reposer explicitement sur le rapport calculé (18), pas sur une affirmation générale non chiffrée.
\end{corrige}

\begin{corrige}[Exercice 6 — Quotient respiratoire de graines germées]
\begin{enumerate}
\item Formule : $\mathrm{QR} = \dfrac{V_{\ce{CO2}}\ \text{produit}}{V_{\ce{O2}}\ \text{consommé}}$.
\item À chaque heure, $V_{\ce{CO2}} = V_{\ce{O2}}$ (\num{12}/\num{12}, \num{24}/\num{24}, \ldots, \num{72}/\num{72}), donc :
\[ \mathrm{QR} = \mathbf{1,0} \quad \text{à chaque mesure (constante).} \]
\item Courbe $\mathrm{QR} = f(t)$ : droite horizontale d'ordonnée 1.
\begin{popfigure}
\begin{tikzpicture}
\begin{axis}[
width=0.85\linewidth, height=5.5cm,
xlabel={Temps (h)}, ylabel={QR},
xmin=0, xmax=6.5, ymin=0, ymax=1.6,
xtick={0,1,2,3,4,5,6}, ytick={0,0.5,1,1.5},
grid=major, grid style={popline}]
\addplot[popcurve, thick, mark=*, mark size=2pt] coordinates {(1,1) (2,1) (3,1) (4,1) (5,1) (6,1)};
\end{axis}
\end{tikzpicture}
\legende{Quotient respiratoire de graines de haricot germées : QR constant et égal à 1.}
\end{popfigure}
\item \textbf{Interprétation :} le bilan de la respiration du glucose $\ce{C6H12O6 + 6 O2 -> 6 CO2 + 6 H2O}$ donne $\mathrm{QR} = \frac{6}{6} = 1$. Un QR égal à 1 signifie donc que les graines germées oxydent majoritairement des \cle{glucides} (glucose issu de l'hydrolyse de l'amidon de réserve de la graine).
\end{enumerate}
\textbf{Point de vigilance (type Probatoire) :} la courbe exige des axes gradués, légendés avec unités, et un titre ; l'interprétation doit s'appuyer sur le bilan équilibré cité. Barème indicatif : formule 1 pt ; calculs 2 pts ; courbe 3 pts ; interprétation 4 pts.
\end{corrige}

\begin{corrige}[Exercice 7 — Légende d'un schéma de mitochondrie]
\begin{enumerate}
\item Identification des structures :
\begin{itemize}
\item \textbf{1 : membrane externe} (enveloppe lisse, perméable aux petites molécules) ;
\item \textbf{2 : espace intermembranaire} (entre les deux membranes ; les protons y sont accumulés pendant le fonctionnement de la chaîne respiratoire) ;
\item \textbf{3 : matrice mitochondriale} (gel interne contenant les enzymes du cycle de Krebs, l'ADN mitochondrial et les ribosomes) ;
\item \textbf{4 : crête mitochondriale} (repli de la membrane interne portant la chaîne respiratoire et les ATP synthases).
\end{itemize}
\item Tableau complété :
\begin{center}
\begin{tabularx}{\linewidth}{|X|X|}
\hline
\rowcolor{popblueL} Étape & Compartiment (numéro et nom) \\
\hline
Glycolyse & Aucun : cytosol (hyaloplasme), hors mitochondrie \\
\hline
Cycle de Krebs & 3 — Matrice mitochondriale \\
\hline
Chaîne respiratoire & 4 — Crête mitochondriale (membrane interne) \\
\hline
\end{tabularx}
\end{center}
\item La glycolyse ne peut pas être placée sur ce schéma car elle se déroule dans le \textbf{cytosol}, à l'extérieur de la mitochondrie ; or le schéma ne représente que l'organite lui-même.
\end{enumerate}
\textbf{Point de vigilance :} l'erreur classique consiste à placer la glycolyse dans la matrice ; retenir : \emph{glycolyse = cytosol, Krebs = matrice, chaîne respiratoire = crêtes}.
\end{corrige}

\begin{corrige}[Exercice 8 — Tableau comparatif respiration / fermentation]
\begin{enumerate}
\item Tableau comparatif :
\begin{center}
\begin{tabularx}{\linewidth}{|X|X|X|}
\hline
\rowcolor{popblueL} \textbf{Critère} & \textbf{Respiration cellulaire} & \textbf{Fermentation} \\
\hline
Localisation & Cytosol (glycolyse) + mitochondrie (Krebs, chaîne respiratoire) & Cytosol uniquement \\
\hline
Besoin en \ce{O2} & Oui (voie aérobie, \ce{O2} accepteur final d'électrons) & Non (voie anaérobie) \\
\hline
Produits obtenus & \ce{CO2} + \ce{H2O} (oxydation complète) & Lactate / acide lactique (lactique) ou éthanol + \ce{CO2} (alcoolique) \\
\hline
ATP par glucose & $\approx 36$ ATP & 2 ATP \\
\hline
\end{tabularx}
\end{center}
\item \textbf{Paragraphe de synthèse (corrigé type) :} Un organisme pluricellulaire comme l'Homme a des besoins énergétiques considérables et permanents : contraction musculaire, transports actifs, biosynthèses, maintien de la température corporelle. Seule la respiration cellulaire, avec environ 36 ATP par glucose, peut couvrir durablement ces besoins en exploitant la quasi-totalité de l'énergie du glucose grâce au dioxygène apporté par la ventilation et la circulation sanguine. La fermentation, limitée à 2 ATP par glucose, impose une consommation de glucose 18 fois plus élevée : elle épuiserait vite les réserves de glycogène. De plus, la fermentation lactique accumule du lactate qui acidifie le milieu cellulaire et perturbe le fonctionnement des enzymes et des fibres musculaires. Elle ne constitue donc qu'une \textbf{voie d'urgence temporaire}, utilisée lors d'efforts brefs et intenses, le temps que l'apport en \ce{O2} redevienne suffisant.
\end{enumerate}
\textbf{Point de vigilance (type Probatoire) :} un paragraphe argumenté doit mobiliser au moins trois savoirs chiffrés (36 vs 2 ATP, rapport 18, acidification) et aboutir à une conclusion explicite. Barème indicatif : tableau 4 pts (1 pt par ligne correcte) ; paragraphe 6 pts (exactitude scientifique 3 pts, argumentation chiffrée 2 pts, cohérence rédactionnelle 1 pt).
\end{corrige}

\begin{corrige}[Exercice 9 — L'acide lactique du sprinteur (étude de cas)]
\begin{enumerate}
\item Le lactate (acide lactique) est produit dans le \textbf{cytosol} des fibres musculaires, par la \cle{fermentation lactique} : le pyruvate issu de la glycolyse est réduit en lactate, ce qui régénère le \ce{NAD+} nécessaire à la poursuite de la glycolyse.
\item Pendant un 400 m, l'effort est maximal : la demande en ATP dépasse la vitesse maximale de la filière aérobie, limitée par l'apport en \ce{O2} aux muscles (débit cardiaque, ventilation). Les mitochondries sont présentes mais \textbf{saturées} ; la glycolyse produit alors du pyruvate plus vite que le cycle de Krebs ne peut l'oxyder. L'excès de pyruvate est détourné vers la fermentation lactique, qui fournit rapidement (mais faiblement) de l'ATP.
\item Courbe (échelles demandées : 1 cm / 5 min ; 1 cm / 2 mmol.L$^{-1}$) :
\begin{popfigure}
\begin{tikzpicture}
\begin{axis}[
width=0.9\linewidth, height=7cm,
xlabel={Temps (min)}, ylabel={Lactate sanguin (mmol/L)},
xmin=-6, xmax=33, ymin=0, ymax=14,
xtick={0,5,10,15,20,25,30}, ytick={0,2,4,6,8,10,12,14},
grid=major, grid style={popline}]
\addplot[popcurve, thick, mark=*, mark size=2pt] coordinates {(-5,1) (0,12.5) (5,9) (15,4.5) (30,1.2)};
\end{axis}
\end{tikzpicture}
\legende{Évolution du taux de lactate sanguin : pic à l'arrivée (12,5 mmol/L), puis décroissance pendant la récupération jusqu'à la valeur de repos.}
\end{popfigure}
\item Après l'effort, la ventilation et le débit cardiaque restent élevés : c'est la \cle{dette d'oxygène}. Le dioxygène ainsi consommé en excès permet : (a) la \textbf{réoxydation} d'une partie du lactate en \ce{CO2} et \ce{H2O} dans les mitochondries (cœur, muscles, foie), qui fournit de l'ATP ; (b) la \textbf{reconversion} du reste en glucose par le foie (néoglucogenèse). Le taux sanguin redescend donc progressivement vers la valeur de repos (environ \num{1} mmol/L à 30 min).
\item Deux moyens fondés sur le cours :
\begin{itemize}
\item \textbf{L'entraînement d'endurance} : il augmente le nombre et la taille des mitochondries musculaires et la capillarisation, ce qui améliore l'utilisation de l'\ce{O2} et retarde le recours à la fermentation ;
\item \textbf{Une alimentation riche en glucides} avant l'épreuve : elle maximise les réserves de glycogène musculaire et hépatique, substrat de la respiration cellulaire.
\end{itemize}
\end{enumerate}
\textbf{Point de vigilance (type Probatoire) :} à la question 2, la justification attendue est la \emph{saturation de la filière aérobie par l'insuffisance d'apport en \ce{O2}}, et non « l'absence de mitochondries ». Barème indicatif : Q1 : 2 pts ; Q2 : 3 pts ; courbe : 4 pts (axes + échelles + points + allure) ; Q4 : 4 pts ; Q5 : 2 pts.
\end{corrige}

\begin{corrige}[Exercice 10 — Expérience sur les échanges gazeux de la souris]
\begin{enumerate}
\item \textbf{Analyse :} entre 0 et 30 min, la teneur en \ce{O2} diminue régulièrement (de \num{21,0}\,\% à \num{16,5}\,\%, soit $-\num{4,5}$ points) tandis que celle en \ce{CO2} augmente régulièrement (de \num{0,03}\,\% à \num{3,6}\,\%, soit $+\num{3,57}$ points). Ces variations opposées et quasi proportionnelles au temps traduisent une \textbf{consommation continue de \ce{O2}} et une \textbf{production continue de \ce{CO2}} par la souris : elles mettent en évidence la respiration cellulaire de l'animal.
\item Quotient respiratoire :
\[ \mathrm{QR} = \frac{\Delta \ce{CO2}}{\Delta \ce{O2}} = \frac{\num{3,6} - \num{0,03}}{\num{21,0} - \num{16,5}} = \frac{\num{3,57}}{\num{4,5}} \approx \mathbf{0,79}. \]
Un QR voisin de \num{0,8}, intermédiaire entre \num{0,7} (lipides) et 1 (glucides), indique que la souris oxyde un \textbf{mélange de nutriments, avec une part importante de lipides}, ce qui est typique d'un animal au repos.
\item L'organite responsable est la \textbf{mitochondrie}. Le rôle exact du dioxygène : il est l'\cle{accepteur final des électrons} de la chaîne respiratoire ; il capte les électrons et les protons pour former de l'eau ($\ce{O2 + 4 H+ + 4 e- -> 2 H2O}$), ce qui maintient le fonctionnement de la chaîne et donc la synthèse d'ATP.
\item La souris est un animal \textbf{homéotherme} : elle maintient sa température corporelle constante. Quand la température ambiante baisse de \SI{10}{\degreeCelsius}, les pertes de chaleur augmentent ; l'animal accroît sa \textbf{thermogenèse} (frissons, métabolisme de production de chaleur), ce qui exige davantage d'ATP, donc une respiration cellulaire plus intense, donc une consommation de \ce{O2} plus élevée.
\item Si la teneur en \ce{O2} devient très faible (hypoxie), la chaîne respiratoire ralentit puis s'arrête faute d'accepteur final d'électrons. Les cellules basculent vers la \textbf{fermentation lactique} : rendement effondré (2 ATP/glucose au lieu de 36), épuisement rapide des réserves glucidiques et accumulation de lactate (acidose), pouvant entraîner la mort de l'animal si l'hypoxie persiste.
\end{enumerate}
\textbf{Point de vigilance (type Probatoire) :} pour le QR, utiliser les \emph{variations} ($\Delta$) des teneurs et non les valeurs brutes ; citer explicitement « accepteur final des électrons » à la question 3. Barème indicatif : Q1 : 3 pts ; Q2 : 3 pts ; Q3 : 3 pts ; Q4 : 3 pts ; Q5 : 3 pts.
\end{corrige}

\begin{corrige}[Exercice 11 — Alimentation et énergie chez le sportif camerounais]
\begin{enumerate}
\item Bilan simplifié et équilibré :
\[ \ce{C6H12O6 + 6 O2 -> 6 CO2 + 6 H2O} + \text{énergie } (\approx 36\ \text{ATP}) \]
Le \textbf{nutriment énergétique} est le glucose (\ce{C6H12O6}) ; le \textbf{comburant} est le dioxygène (\ce{O2}).
\item Quantité de glucose nécessaire :
\[ n_{\text{glucose}} = \frac{n_{\mathrm{ATP}}}{36} = \frac{\num{1200}}{36} \approx \mathbf{33,3\ mol} \]
\[ m = n \times M = \num{33,3} \times 180 \approx \mathbf{6000\ g} \approx \mathbf{6,0\ kg\ de\ glucose}. \]
\item D'après le bilan, 1 mol de glucose consomme 6 mol de \ce{O2} :
\[ n_{\ce{O2}} = 6 \times \num{33,3} \approx \num{200}\ \text{mol} \quad\Longrightarrow\quad V_{\ce{O2}} = \num{200} \times \num{24} = \mathbf{4800\ L} \quad (\text{à CNTP, } V_m = \num{24}\ \text{L/mol}). \]
\item Les glucides sont recommandés car, à quantité d'ATP produite, leur oxydation exige \textbf{moins de dioxygène} que celle des lipides : pour les glucides, $\frac{36\ \text{ATP}}{6\ \text{mol}\ \ce{O2}} = 6$ ATP par mole d'\ce{O2}, contre environ \num{4,6} pour les lipides. Or pendant un match, l'apport en \ce{O2} est le facteur limitant : les glucides sont donc le carburant le plus efficace en conditions d'effort intense. De plus, ils sont stockés sous forme de glycogène musculaire, immédiatement mobilisable, et leur digestion est plus rapide que celle des lipides.
\item \textbf{Correction de l'affirmation :} elle est doublement inexacte.
\begin{itemize}
\item D'une part, « respirer sans oxygène » est une contradiction : la \textbf{respiration} est par définition \emph{aérobie} ; la voie sans oxygène s'appelle la \textbf{fermentation} (lactique dans le muscle).
\item D'autre part, le lactate produit pendant l'effort est éliminé en moins d'une heure de récupération : il ne peut pas causer des douleurs apparaissant 24 à 48 h plus tard. Les \cle{courbatures} sont dues à des \textbf{microlésions des fibres musculaires} (micro-déchirures des membranes et des structures contractiles lors des efforts inhabituels) suivies d'une réaction inflammatoire locale.
\end{itemize}
\end{enumerate}
\textbf{Point de vigilance (type Probatoire) :} distinguer rigoureusement respiration (aérobie, 36 ATP, \ce{CO2} + \ce{H2O}) et fermentation (anaérobie, 2 ATP, lactate) ; dans les calculs, annoncer la relation issue du bilan stœchiométrique avant toute application numérique. Barème indicatif : Q1 : 2 pts ; Q2 : 3 pts ; Q3 : 2 pts ; Q4 : 3 pts ; Q5 : 5 pts (correction 2 pts, distinction respiration/fermentation 2 pts, origine réelle 1 pt).
\end{corrige}

\newpage

\coursec{Fiche bilan}

\begin{popfigure}
\centering
\begin{tikzpicture}[mindmap, grow cyclic,
  every node/.style={concept, font=\small},
  root concept/.append style={concept color=popblue, font=\small\bfseries, text=white},
  level 1/.append style={level distance=3.6cm, sibling angle=60, font=\scriptsize},
  level 1 concept/.append style={concept color=poporange!80},
  text=white]
\node[root concept]{Production d'énergie cellulaire}
  child[concept color=popgreen!85]{ node {Échanges gazeux : consommation de \ce{O2}, rejet de \ce{CO2}} }
  child[concept color=popblue!75]{ node {Mitochondrie : double membrane, crêtes, matrice} }
  child[concept color=poporange!90]{ node {3 étapes : glycolyse, cycle de Krebs, chaîne respiratoire} }
  child[concept color=poppurple!85]{ node {Fermentation : voie sans \ce{O2}, faible rendement} }
  child[concept color=poppink!85]{ node {Rendement : 36 ATP (aérobie) vs 2 ATP (anaérobie)} }
  child[concept color=popgold!90]{ node {Facteurs : entraînement, alimentation, oxygénation} };
\end{tikzpicture}
\legende{Carte mentale de la leçon : l'organisme améliore sa production d'énergie grâce à la respiration cellulaire mitochondriale (voie aérobie, très rentable) et, à défaut d'oxygène, à la fermentation (voie anaérobie, peu rentable).}
\end{popfigure}

\begin{bilan}
\textbf{1. Définitions clés}
\begin{itemize}
\item \cle{Respiration cellulaire} : ensemble des réactions d'oxydation du glucose, en présence de dioxygène, fournissant de l'énergie stockée sous forme d'\cle{ATP} (adénosine triphosphate).
\item \cle{Glycolyse} : dégradation du glucose (6 C) en deux pyruvates (3 C), dans le \textbf{cytoplasme}, sans oxygène. Bilan : 2 ATP nets.
\item \cle{Cycle de Krebs} : oxydation complète de l'acétyl-CoA dans la \textbf{matrice mitochondriale}, produisant \ce{CO2} et des coenzymes réduits (NADH,H$^+$ et FADH$_2$).
\item \cle{Chaîne respiratoire} : transfert d'électrons le long des \textbf{crêtes mitochondriales} ; le dioxygène est l'accepteur final (réduit en \ce{H2O}) ; l'énergie libérée sert à synthétiser l'ATP (phosphorylation oxydative).
\item \cle{Fermentation} : dégradation incomplète du glucose \textbf{sans oxygène}, dans le cytoplasme ; elle régénère le NAD$^+$ nécessaire à la glycolyse. Bilan : seulement 2 ATP.
\end{itemize}

\textbf{2. Formules et bilans à connaître par c\oe ur}
\begin{center}
\begin{tabularx}{\linewidth}{|l|X|}
\hline
\rowcolor{popblueL} \textbf{Processus} & \textbf{Bilan} \\
\hline
Respiration cellulaire & \ce{C6H12O6 + 6O2 -> 6CO2 + 6H2O} + énergie ($\approx$ 36 ATP) \\
\hline
Fermentation lactique (muscle) & glucose $\rightarrow$ 2 acide lactique + 2 ATP \\
\hline
Fermentation alcoolique (levure) & glucose $\rightarrow$ 2 éthanol + 2 \ce{CO2} + 2 ATP \\
\hline
Quotient respiratoire & $QR = \dfrac{\text{volume de }\ce{CO2}\text{ produit}}{\text{volume de }\ce{O2}\text{ consommé}}$ ($QR = 1$ pour le glucose) \\
\hline
\end{tabularx}
\end{center}

\textbf{3. La mitochondrie en 4 points}
\begin{itemize}
\item \textbf{Membrane externe} lisse ; \textbf{membrane interne} repliée en \textbf{crêtes} (siège de la chaîne respiratoire).
\item \textbf{Matrice} : siège du cycle de Krebs ; \textbf{espace intermembranaire} : accumulation des H$^+$.
\item L'ATP synthase laisse repasser les H$^+$ et fabrique l'ATP.
\item Les cellules très actives (muscle cardiaque, spermatozoïdes) sont riches en mitochondries.
\end{itemize}

\textbf{4. Méthodes express}
\begin{itemize}
\item \emph{Exploiter une expérience} (tube avec germes de haricot + potasse) : la potasse absorbe le \ce{CO2} ; la remontée du liquide prouve la consommation de \ce{O2}.
\item \emph{Comparer les rendements} : toujours rapporter le nombre d'ATP produit par molécule de glucose (36 contre 2, soit 18 fois plus).
\item \emph{Interpréter un effort musculaire} : essoufflement = augmentation de la ventilation pour apporter plus de \ce{O2} ; courbatures = acide lactique issu de la fermentation quand l'apport en \ce{O2} ne suffit plus.
\end{itemize}

\textbf{5. Facteurs améliorant le rendement énergétique}
\begin{itemize}
\item \textbf{Entraînement} régulier : augmentation du nombre et de la taille des mitochondries, meilleure irrigation sanguine.
\item \textbf{Alimentation} équilibrée : apport de glucides (substrats énergétiques) et de fer (hémoglobine, transport de \ce{O2}).
\item \textbf{Oxygénation} : bonne ventilation, absence de tabac (le CO bloque le transport de \ce{O2}).
\end{itemize}
\end{bilan}

\begin{aretenir}[Checklist avant le Probatoire]
\begin{itemize}
\item[$\square$] Je sais écrire le bilan global de la respiration cellulaire : \ce{C6H12O6 + 6O2 -> 6CO2 + 6H2O}.
\item[$\square$] Je sais situer chaque étape : glycolyse (cytoplasme), cycle de Krebs (matrice), chaîne respiratoire (crêtes).
\item[$\square$] Je sais légender une coupe de mitochondrie : membrane externe, membrane interne, crêtes, matrice.
\item[$\square$] Je connais le rôle du dioxygène : accepteur final d'électrons, réduit en eau.
\item[$\square$] Je connais les rendements : environ 36 ATP en aérobie contre 2 ATP en anaérobie.
\item[$\square$] Je sais distinguer fermentation lactique (muscle) et fermentation alcoolique (levure).
\item[$\square$] Je sais expliquer pourquoi la fermentation est indispensable sans \ce{O2} : elle régénère le NAD$^+$.
\item[$\square$] Je sais exploiter une expérience montrant la consommation de \ce{O2} et la production de \ce{CO2}.
\item[$\square$] Je sais calculer et interpréter le quotient respiratoire $QR$.
\item[$\square$] Je sais relier essoufflement et courbatures aux phénomènes cellulaires.
\item[$\square$] Je sais citer trois facteurs améliorant le rendement énergétique (entraînement, alimentation, oxygénation).
\item[$\square$] Je ne confonds jamais respiration cellulaire (cellule) et respiration pulmonaire (ventilation).
\end{itemize}
\end{aretenir}

\findecours

\end{document}
